\chapter{The Quot scheme}
\label{chap:Picard_QuotScheme}
% archon:covers AlgebraicJacobian/Picard/QuotScheme.lean AlgebraicJacobian/Picard/GradedHilbertSerre.lean

% STRATEGY NOTE. This chapter packages the Quot-foundations layer of the
% Grothendieck--Altman--Kleiman Quot-scheme construction
% \(\mathrm{Quot}^{\Phi,L}_{E/X/S}\): the per-fiber Hilbert polynomial, the
% Quot functor of \(T\)-flat coherent quotients of \(E_T\) on
% \(X_T = X \times_S T\), and the Grassmannian \emph{scheme} together with
% its representability. The Quot functor underlies the relative Picard
% functor via the Hilbert scheme \(\Hilb_{C/k} = \Quot_{\mathcal{O}_C/C/k}\)
% -- the case \(E = \mathcal{O}_X\) of the Quot functor.
%
% LOAD-BEARING SUB-PREREQUISITE. The Quot construction embeds Quot as a
% locally closed subfunctor of a Grassmannian
% \(\mathrm{Grass}(W \otimes_{\mathcal{O}_S} \mathrm{Sym}^r V, \Phi(r))\).
% Mathlib (at the pinned commit) carries no Grassmannian \emph{scheme} -- only
% a linear-algebra Grassmannian as a finite-rank-quotient variety. The
% scheme-level Grassmannian over an arbitrary noetherian base is therefore an
% in-project sub-build. This chapter introduces it as
% \cref{def:grassmannian_scheme} and states its representability
% (\cref{thm:grassmannian_representable}).

\section{Setup and motivation}
\label{sec:quot_setup}

Throughout this chapter, let \(S\) be a noetherian scheme, \(\pi : X \to S\) a
projective morphism, \(L\) a relatively very ample line bundle on \(X\) over \(S\),
and \(E\) a coherent sheaf on \(X\). The Quot construction packages the
\(T\)-flat coherent quotients \(E_T \twoheadrightarrow \mathcal{F}\) of the
pull-back \(E_T\) to \(X_T = X \times_S T\), modulo the equivalence
\(\langle \mathcal{F}, q\rangle = \langle \mathcal{F}', q'\rangle\) defined by
\(\ker(q) = \ker(q')\). Fixing the Hilbert polynomial
\(\Phi(\lambda) = \chi(\mathcal{F}|_{X_t}(\lambda)) \in \mathbb{Q}[\lambda]\)
(constant on connected \(T\) by flatness)
stratifies the Quot functor into pieces \(\Quot^{\Phi,L}_{E/X/S}\), each of
which Grothendieck proved is representable by a projective \(S\)-scheme.
This chapter develops the foundations of that picture: the Hilbert
polynomial (\cref{def:hilbert_polynomial}), the Quot functor
(\cref{def:quot_functor}), and the Grassmannian scheme
(\cref{def:grassmannian_scheme}) with its representability
(\cref{thm:grassmannian_representable}).

% SOURCE: [Nitsure], "Construction of Hilbert and Quot Schemes", \S 1;
% FGA Explained Ch.~5, AMS MSM 123 (read from
% references/nitsure-hilbert-quot-src/nitsure-hilbert-quot.tex, L287-L500).
% The Hartshorne reference for the polynomial-eventually
% property of Euler characteristics is [Hartshorne], III.5.2 (the scanned
% PDF at references/hartshorne-algebraic-geometry.pdf has no text layer, so
% the verbatim source-quote is taken from
% Nitsure \S 1; Hartshorne III.5.2 is cited in prose as the textbook
% reference).

\section{Hilbert polynomial of a coherent sheaf}
\label{sec:hilbert_polynomial}

\begin{definition}\leanok
[Hilbert polynomial]
  \label{def:hilbert_polynomial}
  \lean{AlgebraicGeometry.Scheme.hilbertPolynomial}
  \uses{thm:hilbertPoly_of_sectionModule}
  % ENCODING. `def:hilbert_polynomial` is formalized via the GRADED HILBERT
  % FUNCTION of the section module ⊕_{m≥0} Γ(X_s, F_s ⊗ L_s^{⊗m})
  % (def:sectionGradedModule), whose polynomiality is extracted in
  % thm:hilbertPoly_of_sectionModule from the graded Hilbert–Serre rationality
  % lemma (lem:gradedHilbertSerre_rational) together with the Mathlib lemma
  % Polynomial.existsUnique_hilbertPoly (lem:hilbertPoly_exists_mathlib). The
  % construction uses only H^0; for m≫0 it agrees with the Euler characteristic
  % χ by Serre vanishing, so it is the classical Hilbert polynomial and no
  % numerical invariant is lost.
  % SOURCE: [Nitsure], \S 1, sub-section "Stratification by Hilbert
  % Polynomials" (read from
  % references/nitsure-hilbert-quot-src/nitsure-hilbert-quot.tex, L453-L478).
  % SOURCE QUOTE:
  % "Let \(X\) be a finite type scheme over
  %  a field \(k\), together with a line bundle \(L\).
  %  Recall that if \(F\) is a coherent sheaf on
  %  \(X\) whose support is proper over \(k\), then
  %  the \textbf{Hilbert polynomial} \({\Phi}\in \Q[\lambda]\) of \(F\) is defined by
  %  the function
  %  \(\){\Phi}(m) = \chi(F(m)) =
  %  \sum_{i=0}^n (-1)^i \dim_k H^i(X, F\otimes L^{\otimes m})\(\)
  %  where the dimensions of the cohomologies are finite because of the
  %  coherence and properness conditions. The fact that \(\chi(F(m))\)
  %  is indeed a polynomial in \(m\) under the above assumption is a special
  %  case of what is known as Snapper's Lemma (see Kleiman [K] for a proof).
  %
  %  Let \(X\to S\) be a finite type morphism of noetherian schemes,
  %  and let \(L\) be a line bundle on \(X\). Let \(\F\)
  %  be any coherent sheaf on
  %  \(X\) whose schematic support is proper over \(S\). Then for each \(s\in S\),
  %  we get a polynomial \({\Phi}_s \in  \Q[\lambda]\)
  %  which is the Hilbert polynomial of
  %  the restriction \(\F_s = F|_{X_s}\) of \(\F\) to the fiber \(X_s\) over \(s\),
  %  calculated with respect to the line bundle \(L_s = L|_{X_s}\).
  %  If \(\F\) is flat over \(S\) then the function \(s\mapsto {\Phi}_s\)
  %  from the set of points of \(S\) to the polynomial ring \(\Q[\lambda]\) is
  %  known to be locally constant on \(S\)."
  \textit{Source: [Nitsure], \S 1 (FGA Explained Ch.~5); cf.\ [Hartshorne],
  III.5.2.}
  Let \(X \to S\) be a finite-type morphism of noetherian schemes and let \(L\)
  be a line bundle on \(X\). For a coherent sheaf \(\mathcal{F}\) on \(X\) whose
  schematic support is proper over \(S\) and a point \(s \in S\), write \(X_s\)
  for the fibre over \(s\), \(\mathcal{F}_s = \mathcal{F}|_{X_s}\) and
  \(L_s = L|_{X_s}\). The \emph{Hilbert polynomial of \(\mathcal{F}\) at \(s\)
  relative to \(L\)} is the unique polynomial
  \(\Phi_{\mathcal{F},s} \in \mathbb{Q}[\lambda]\) that agrees, for all
  \(m \gg 0\), with the graded Hilbert function
  \[
    m \;\longmapsto\; \dim_{\kappa(s)} \Gamma\bigl(X_s,\, \mathcal{F}_s \otimes L_s^{\otimes m}\bigr)
  \]
  of the section module
  \(\bigoplus_{m \ge 0} \Gamma(X_s, \mathcal{F}_s \otimes L_s^{\otimes m})\)
  (\cref{def:sectionGradedModule}), viewed as a finitely generated graded module
  over the homogeneous section ring
  \(\bigoplus_{m \ge 0} \Gamma(X_s, L_s^{\otimes m})\)
  (\cref{def:sectionGradedRing}). Existence and uniqueness of
  \(\Phi_{\mathcal{F},s}\) are the content of
  \cref{thm:hilbertPoly_of_sectionModule}.

  This is the classical Hilbert polynomial. When \(\mathcal{F}_s\) is coherent
  and \(L_s\) is ample on the proper \(\kappa(s)\)-scheme \(X_s\), Serre vanishing
  gives \(H^i(X_s, \mathcal{F}_s \otimes L_s^{\otimes m}) = 0\) for all \(i > 0\)
  and \(m \gg 0\), so for \(m \gg 0\)
  \[
    \dim_{\kappa(s)} \Gamma(X_s, \mathcal{F}_s \otimes L_s^{\otimes m})
    \;=\; \chi(X_s,\, \mathcal{F}_s \otimes L_s^{\otimes m})
    \;=\; \sum_{i \ge 0} (-1)^i \dim_{\kappa(s)} H^i(X_s,\, \mathcal{F}_s \otimes L_s^{\otimes m}),
  \]
  and hence \(\Phi_{\mathcal{F},s}\) coincides with the Euler-characteristic
  Hilbert polynomial of Snapper's Lemma -- no numerical invariant is lost. The
  construction above, however, uses only the global sections \(H^0\): the higher
  cohomology enters solely through the classical statement that the two functions
  agree for \(m \gg 0\), never through the definition of \(\Phi_{\mathcal{F},s}\)
  itself. When \(\mathcal{F}\) is \(S\)-flat the function
  \(s \mapsto \Phi_{\mathcal{F},s}\) is locally constant on \(S\), so \(\Phi\) is
  well-defined on connected components.
\end{definition}

\section{Graded Hilbert polynomial}
\label{sec:graded_hilbert_polynomial}

This section gives the graded-Hilbert-function construction that underlies
\cref{def:hilbert_polynomial}. Working over a single fibre \(X_s\) with its
ample line bundle \(L_s\) and coherent sheaf \(\mathcal{F}_s\), it builds the
section graded ring (\cref{def:sectionGradedRing}) and module
(\cref{def:sectionGradedModule}), records their finite generation
(\cref{lem:sectionGradedModule_fg}), and extracts the Hilbert polynomial
(\cref{thm:hilbertPoly_of_sectionModule}) by combining the graded
Hilbert--Serre rationality of the Hilbert series
(\cref{lem:gradedHilbertSerre_rational}) with Mathlib's polynomial-extraction
lemma (\cref{lem:hilbertPoly_exists_mathlib}). The route uses only \(H^0\)
(global sections), so it is independent of higher coherent cohomology.

\begin{definition}
[Section graded ring]
  \label{def:sectionGradedRing}
  \lean{AlgebraicGeometry.sectionGradedRing}
  \uses{lem:sectionGradedRing_gcommSemiring}
  % NOTE: Formalization is blocked on the absence of tensor products for sheaves of modules in
  % Mathlib at the pinned commit. The degree-\(m\) component \(\Gamma(X_s, L_s^{\otimes m})\)
  % requires a tensor power \(L_s^{\otimes m}\) of a line bundle; this infrastructure (tensor
  % product of sheaves of modules with lax-monoidal global sections) is not yet available, making
  % this definition a prerequisite sub-build.
  % SOURCE: [Nitsure], "Construction of Hilbert and Quot Schemes", \S 1,
  % sub-section "Stratification by Hilbert Polynomials" (read from
  % references/nitsure-hilbert-quot-src/nitsure-hilbert-quot.tex, L453-L464).
  % SOURCE QUOTE:
  % "Let $X$ be a finite type scheme over
  %  a field $k$, together with a line bundle $L$.
  %  Recall that if $F$ is a coherent sheaf on
  %  $X$ whose support is proper over $k$, then
  %  the \textbf{Hilbert polynomial} ${\Phi}\in \Q[\lambda]$ of $F$ is defined by
  %  the function
  %  $${\Phi}(m) = \chi(F(m)) =
  %  \sum_{i=0}^n (-1)^i \dim_k H^i(X, F\otimes L^{\otimes m})$$
  %  where the dimensions of the cohomologies are finite because of the
  %  coherence and properness conditions."
  \textit{Source: [Nitsure], \S 1 (FGA Explained Ch.~5); the section graded
  ring is the homogeneous coordinate ring of [Hartshorne], I.7.}
  Let \(s \in S\) and let \(L_s\) be a line bundle on the fibre \(X_s\), a scheme
  of finite type over the residue field \(\kappa(s)\). The \emph{section graded
  ring} of \(L_s\) is the graded commutative ring
  \[
    R(X_s, L_s) \;=\; \bigoplus_{m \ge 0} \Gamma\bigl(X_s,\, L_s^{\otimes m}\bigr),
  \]
  with degree-\(m\) component \(\Gamma(X_s, L_s^{\otimes m})\) and multiplication
  the tensor product of sections
  \(\Gamma(X_s, L_s^{\otimes m}) \otimes_{\kappa(s)}
  \Gamma(X_s, L_s^{\otimes m'}) \to \Gamma(X_s, L_s^{\otimes (m+m')})\),
  \(\sigma \otimes \tau \mapsto \sigma \cdot \tau\). Its degree-zero component is
  \(\Gamma(X_s, \mathcal{O}_{X_s})\), which contains \(\kappa(s)\). When \(X_s\) is
  proper over \(\kappa(s)\) and \(L_s\) is ample, \(R(X_s, L_s)\) is a finitely
  generated -- hence Noetherian -- graded \(\kappa(s)\)-algebra (a suitable
  Veronese power of \(L_s\) is very ample and embeds \(X_s\) into a projective
  space, exhibiting \(R(X_s, L_s)\) up to a finite extension as a quotient of a
  polynomial ring).
\end{definition}

\begin{definition}
[Section graded module]
  \label{def:sectionGradedModule}
  \lean{AlgebraicGeometry.sectionGradedModule}
  \uses{def:sectionGradedRing, lem:sectionGradedModule_gmodule}
  % SOURCE: [Nitsure], \S 1, sub-section "Stratification by Hilbert
  % Polynomials" (read from
  % references/nitsure-hilbert-quot-src/nitsure-hilbert-quot.tex, L453-L464).
  % SOURCE QUOTE:
  % "Recall that if $F$ is a coherent sheaf on
  %  $X$ whose support is proper over $k$, then
  %  the \textbf{Hilbert polynomial} ${\Phi}\in \Q[\lambda]$ of $F$ is defined by
  %  the function
  %  $${\Phi}(m) = \chi(F(m)) =
  %  \sum_{i=0}^n (-1)^i \dim_k H^i(X, F\otimes L^{\otimes m})$$
  %  where the dimensions of the cohomologies are finite because of the
  %  coherence and properness conditions."
  \textit{Source: [Nitsure], \S 1 (FGA Explained Ch.~5); cf.\ [Hartshorne],
  I.7.}
  Let \(\mathcal{F}_s\) be a coherent sheaf on the fibre \(X_s\) and \(L_s\) a
  line bundle on \(X_s\). The \emph{section graded module} of \(\mathcal{F}_s\)
  twisted by \(L_s\) is the graded \(R(X_s, L_s)\)-module
  (\cref{def:sectionGradedRing})
  \[
    M(X_s, \mathcal{F}_s, L_s)
    \;=\; \bigoplus_{m \ge 0} \Gamma\bigl(X_s,\, \mathcal{F}_s \otimes L_s^{\otimes m}\bigr),
  \]
  with degree-\(m\) component the \(\kappa(s)\)-vector space
  \(\Gamma(X_s, \mathcal{F}_s \otimes L_s^{\otimes m})\) and scalar action given
  by the section tensor product
  \(\Gamma(X_s, L_s^{\otimes m}) \otimes_{\kappa(s)}
  \Gamma(X_s, \mathcal{F}_s \otimes L_s^{\otimes m'})
  \to \Gamma(X_s, \mathcal{F}_s \otimes L_s^{\otimes (m+m')})\). Its graded Hilbert
  function is \(m \mapsto \dim_{\kappa(s)}
  \Gamma(X_s, \mathcal{F}_s \otimes L_s^{\otimes m})\).
\end{definition}

\begin{lemma}
[Finite generation of the section graded module]
  \label{lem:sectionGradedModule_fg}
  \lean{AlgebraicGeometry.sectionGradedModule_fg}
  \uses{def:sectionGradedModule, def:sectionGradedRing}
  % SOURCE: [Nitsure], \S 1, sub-section "Stratification by Hilbert
  % Polynomials" (read from
  % references/nitsure-hilbert-quot-src/nitsure-hilbert-quot.tex, L453-L464).
  % SOURCE QUOTE:
  % "Recall that if $F$ is a coherent sheaf on
  %  $X$ whose support is proper over $k$, then
  %  the \textbf{Hilbert polynomial} ${\Phi}\in \Q[\lambda]$ of $F$ is defined by
  %  the function
  %  $${\Phi}(m) = \chi(F(m)) =
  %  \sum_{i=0}^n (-1)^i \dim_k H^i(X, F\otimes L^{\otimes m})$$
  %  where the dimensions of the cohomologies are finite because of the
  %  coherence and properness conditions."
  \textit{Source: [Nitsure], \S 1 (FGA Explained Ch.~5); Serre's theorem on
  ample line bundles, [Hartshorne], I.7 / II.5.}
  Let \(X_s\) be proper over the field \(\kappa(s)\), let \(L_s\) be an ample line
  bundle on \(X_s\), and let \(\mathcal{F}_s\) be a coherent sheaf on \(X_s\).
  Then the section graded ring \(R(X_s, L_s)\) (\cref{def:sectionGradedRing}) is a
  Noetherian graded ring, and the section graded module
  \(M(X_s, \mathcal{F}_s, L_s)\) (\cref{def:sectionGradedModule}) is a finitely
  generated graded \(R(X_s, L_s)\)-module. In particular each component
  \(\Gamma(X_s, \mathcal{F}_s \otimes L_s^{\otimes m})\) is a finite-dimensional
  \(\kappa(s)\)-vector space.
\end{lemma}

\begin{proof}
  \uses{def:sectionGradedRing, def:sectionGradedModule}
  Since \(L_s\) is ample and \(X_s\) is proper over \(\kappa(s)\), a Veronese
  power \(L_s^{\otimes e}\) is very ample and gives a closed immersion
  \(X_s \hookrightarrow \mathbb{P}^N_{\kappa(s)}\). By Serre's theorem the section
  graded ring \(R(X_s, L_s)\) is a finitely generated graded
  \(\kappa(s)\)-algebra, so it is Noetherian. For the coherent sheaf
  \(\mathcal{F}_s\), the same ampleness yields, again by Serre, that
  \(\bigoplus_{m \ge 0} \Gamma(X_s, \mathcal{F}_s \otimes L_s^{\otimes m})\) is a
  finitely generated graded module over \(R(X_s, L_s)\): the twists
  \(\mathcal{F}_s \otimes L_s^{\otimes m}\) are globally generated for
  \(m \gg 0\), and the finitely many generating sections in low degrees together
  with the multiplication maps generate \(M(X_s, \mathcal{F}_s, L_s)\) over
  \(R(X_s, L_s)\). Finite dimensionality of each component over \(\kappa(s)\)
  follows from coherence and properness (finiteness of cohomology of coherent
  sheaves on a proper scheme over a field).
\end{proof}

\begin{lemma}[Existence and uniqueness of the Hilbert polynomial of a rational
  graded series]
  \label{lem:hilbertPoly_exists_mathlib}
  \lean{Polynomial.existsUnique_hilbertPoly}
  \mathlibok
  \textit{Provided by Mathlib
  (\texttt{Mathlib.RingTheory.Polynomial.HilbertPoly}).}
  Let \(F\) be a field of characteristic zero, let \(p \in F[X]\), and let
  \(d \in \mathbb{N}\). Then there is a unique polynomial \(h \in F[X]\) such that
  for some \(N \in \mathbb{N}\) and all \(n > N\) the coefficient of \(X^n\) in the
  power-series expansion of \(p \cdot (1 - X)^{-d} \in F\llbracket X\rrbracket\)
  equals \(h(n)\):
  \[
    \exists!\, h \in F[X],\ \exists N,\ \forall n > N : \;
    \bigl[ X^n \bigr]\bigl( p \cdot (1-X)^{-d} \bigr) \;=\; h(n).
  \]
  Equivalently, the eventual coefficient function of any rational power series
  with denominator a power of \(1 - X\) is, for \(n \gg 0\), given by a unique
  polynomial \(h\) (Mathlib's \(\mathrm{hilbertPoly}\,p\,d\)). The hypothesis
  \([\mathrm{CharZero}\,F]\) is satisfied by \(F = \mathbb{Q}\), which is the
  instance used to obtain \(\Phi_{\mathcal{F},s} \in \mathbb{Q}[\lambda]\).
\end{lemma}

\begin{lemma}[Additivity of \(\dim_\kappa\) on a short exact sequence
  (rank--nullity)]
  \label{lem:finrank_ses_additive_mathlib}
  \lean{Submodule.finrank_quotient_add_finrank}
  \mathlibok
  \textit{Provided by Mathlib
  (\texttt{Mathlib.LinearAlgebra.Dimension.RankNullity}).}
  Let \(\kappa\) be a field, \(V\) a finite-dimensional \(\kappa\)-vector space and
  \(W \subseteq V\) a subspace. Then
  \[
    \dim_\kappa (V / W) + \dim_\kappa W \;=\; \dim_\kappa V .
  \]
  Equivalently, \(\dim_\kappa = \operatorname{finrank}_\kappa\) is additive on a
  short exact sequence \(0 \to W \to V \to V/W \to 0\) of finite-dimensional
  \(\kappa\)-vector spaces. This is the additivity that converts the degreewise
  short exact sequences of graded components into additive relations among the
  coefficients of the Hilbert series.
\end{lemma}

\begin{lemma}[The rational power series \((1 - X)^{-d}\)]
  \label{lem:invOneSubPow_mathlib}
  \lean{PowerSeries.invOneSubPow}
  \mathlibok
  \textit{Provided by Mathlib
  (\texttt{Mathlib.RingTheory.PowerSeries.WellKnown}).}
  Let \(S\) be a commutative ring and \(d \in \mathbb{N}\). Then \((1 - X)^d\) is a
  unit in the power series ring \(S\llbracket X\rrbracket\), and Mathlib names its
  inverse
  \[
    \operatorname{invOneSubPow}(S, d) \;=\; (1 - X)^{-d}
      \;\in\; S\llbracket X\rrbracket^{\times}.
  \]
  Together with the coercion
  \(\operatorname{toPowerSeries} : F[X] \to F\llbracket X\rrbracket\) of a
  polynomial to its power series, this is the data in which a rational Hilbert
  series \(p \cdot (1 - X)^{-d}\) is expressed; it is the same rational-series shape
  whose eventual coefficients are extracted by
  \cref{lem:hilbertPoly_exists_mathlib}.
\end{lemma}

\begin{lemma}\leanok
[Graded Hilbert--Serre: rationality of the Hilbert series]
  \label{lem:gradedHilbertSerre_rational}
  \lean{AlgebraicGeometry.gradedModule_hilbertSeries_rational}
  \uses{lem:finrank_ses_additive_mathlib, lem:invOneSubPow_mathlib,
    def:ratHilb, lem:ratHilb_ofEventuallyZero, lem:ratHilb_ofDiffEq,
    lem:ratHilb_bump, def:graded_subquotientHilb,
    lem:graded_subquotient_isRatHilb}
  % SOURCE: [Stacks Project], "Commutative Algebra", \S "Noetherian graded
  % rings", Proposition \texttt{proposition-graded-hilbert-polynomial} (Tag
  % 00K1) (read from references/hilbert-serre-algebra.tex, L13893-L13905).
  % SOURCE QUOTE:
  % "Suppose that $S$ is a Noetherian graded ring
  %  and $M$ a finite graded $S$-module. Consider the
  %  function
  %  $$
  %  \mathbf{Z} \longrightarrow K'_0(S_0), \quad
  %  n \longmapsto [M_n]
  %  $$
  %  see Lemma \ref{lemma-graded-module-fg}.
  %  If $S_{+}$ is generated by elements of degree $1$,
  %  then this function is a numerical polynomial."
  \textit{Source: [Stacks Project], "Commutative Algebra", Tag 00K1
  (proposition on graded Hilbert polynomials); cf.\ [Nitsure], \S 1 (Snapper's
  Lemma), [Hartshorne], I.7.}
  Let \(R\) be a Noetherian graded ring whose degree-zero part is a field
  \(\kappa\) and which is generated as a \(\kappa\)-algebra by finitely many
  elements, and let \(M\) be a finitely generated graded \(R\)-module with
  finite-dimensional graded components \(M_n\). Then there exist
  \(p \in \mathbb{Q}[X]\) and \(d \in \mathbb{N}\) and an \(N \in \mathbb{N}\)
  such that for all \(n > N\)
  \[
    \dim_{\kappa} M_n \;=\; \bigl[ X^n \bigr]\bigl( p \cdot (1 - X)^{-d}\bigr);
  \]
  that is, the Hilbert series \(\sum_{n \ge 0} (\dim_\kappa M_n)\, X^n\) is, up to
  a polynomial correction in low degrees, the rational function
  \(p \cdot (1 - X)^{-d}\), expressed in the power-series form
  \(\operatorname{toPowerSeries}(p) \cdot \operatorname{invOneSubPow}(\mathbb{Q}, d)\)
  (\cref{lem:invOneSubPow_mathlib}).

  \emph{Status.} This is the substantive half of the graded Hilbert--Serre
  theorem (rationality of the Hilbert series of a finitely generated graded
  module). It is not covered by Mathlib: \cref{lem:hilbertPoly_exists_mathlib}
  supplies only the converse polynomial-extraction step from a rational series, not
  the production of the rational series from the module. The argument is realised by
  an \emph{ambient subquotient induction} (\cref{lem:graded_subquotient_isRatHilb}):
  the whole induction ranges over pairs of homogeneous submodules \(N' \le N\) of the
  single fixed graded module \(M\), and the Hilbert function is the ambient difference
  \(n \mapsto \dim_\kappa(N \cap M_n) - \dim_\kappa(N' \cap M_n)\). The class of such
  pairs is closed under the kernel and cokernel of multiplication by a degree-one
  element (each presented again as an ambient subquotient pair), so the inductive step
  never forms a graded quotient ring \(R/(x)\), a graded quotient module \(M/xM\), or a
  regrading of a derived carrier. The genuine Mathlib gap is thereby \emph{avoided}
  rather than filled; the residual obligation is the short ambient bookkeeping of the
  subquotient blocks (\cref{def:graded_subquotientHilb},
  \cref{lem:graded_subquotient_ker_coker},
  \cref{lem:graded_subquotient_degreewise_diff},
  \cref{lem:graded_subquotient_finite_transfer},
  \cref{lem:graded_subquotient_isRatHilb}), each of which manipulates only ambient
  families \(n \mapsto N_{\mathrm{aux}} \cap M_n\) of submodules of \(M\).
  % NOTE: The power-series rationality toolkit (\cref{subsec:isRatHilb}) is formalized
  %   and axiom-clean; the present theorem is a closed, axiom-clean Lean declaration
  %   (gradedModule_hilbertSeries_rational). The ambient subquotient route
  %   (\cref{lem:graded_subquotient_isRatHilb}) sidesteps missing graded-quotient API by
  %   working with ambient families of submodules of \(M\).
\end{lemma}

% SOURCE QUOTE PROOF: [Stacks Project], Tag 00K1, proof, inductive step (read
% from references/hilbert-serre-algebra.tex, L13938-L13947).
% "Let $\overline{M} = M/xM$. Note that the map $x : M \to M$
%  is {\it not} a map of graded $S$-modules, since it does
%  not map $M_d$ into $M_d$. Namely, for each $d$ we have the
%  following short exact sequence
%  $$
%  0 \to M_d \xrightarrow{x} M_{d + 1} \to \overline{M}_{d + 1} \to 0
%  $$
%  This proves that $[M_{d + 1}] - [M_d] = [\overline{M}_{d + 1}]$.
%  Hence we win by Lemma \ref{lemma-numerical-polynomial}."
\begin{proof}
  \uses{def:ratHilb, lem:ratHilb_ofEventuallyZero, lem:ratHilb_ofDiffEq,
    lem:ratHilb_bump, lem:graded_degreewise_finrank_diff,
    def:graded_subquotientHilb, lem:graded_subquotient_ker_coker,
    lem:graded_subquotient_degreewise_diff,
    lem:graded_subquotient_finite_transfer, lem:graded_subquotient_isRatHilb,
    lem:graded_homogeneousSubmodule_iSupIndep,
    lem:graded_homogeneousSubmodule_iSup_eq, lem:ideal_homogeneous_span_mathlib,
    lem:isHomogeneousElem_graded_smul_mathlib,
    lem:fg_restrictScalars_of_surjective_mathlib}
    \leanok
  We argue with the abstract graded \(\kappa\)-algebra
  \(R = \bigoplus_{n \ge 0} R_n\), whose degree-zero part \(R_0 = \kappa\) is a
  field, and the finitely generated graded \(R\)-module
  \(M = \bigoplus_{n \ge 0} M_n\) with finite-dimensional components \(M_n\);
  there is no appeal to any geometric realisation of \(R\) or \(M\).

  \emph{Reduction to generation in degree \(1\).} Since \(R\) is Noetherian and
  finitely generated as a \(\kappa\)-algebra, it is generated by finitely many
  homogeneous elements of positive degree. Replacing \(R\) by a Veronese
  subalgebra \(R^{(e)} = \bigoplus_n R_{en}\) for a suitable \(e\) and re-grading
  (and \(M\) by the corresponding Veronese module), we may assume \(R\) is
  generated by its degree-one part, \(\kappa[R_1] = R\). The degree-one part
  \(R_1\) is then a finite-dimensional \(\kappa\)-vector space whose \(\kappa\)-basis
  \(x_1, \dots, x_r\) generates \(R\) as a \(\kappa\)-algebra; \(r = \dim_\kappa R_1\)
  is the number on which we induct.
  % NOTE: The existence of finitely many degree-one generators (the number r the
  % induction is on) is expected to follow from the standing hypotheses --
  % "finitely generated κ-algebra" together with "generated in degree 1" -- so no
  % new hypothesis on gradedModule_hilbertSeries_rational is anticipated. To be
  % confirmed when the prover formalizes the assembly: degree-1 generation gives
  % R₁ as a finite-dimensional κ-vector space spanning R as a κ-algebra, and a
  % κ-basis of R₁ is a finite degree-one generating set.

  \emph{Set-up of the ambient data.} Regard \(M\) as the fixed ambient graded
  \(\kappa\)-module with internal decomposition \(n \mapsto \mathcal{M}_n := M_n\)
  (\cref{def:sectionGradedModule}); the components are finite-dimensional. For each
  basis vector \(x_i \in R_1\) let \(t_i : M \to M\) be multiplication by \(x_i\), a
  \(\kappa\)-linear endomorphism that raises degree by one (it sends \(M_n\) into
  \(M_{n+1}\), \cref{lem:isHomogeneousElem_graded_smul_mathlib}). Since \(R\) is
  commutative the \(t_i\) pairwise commute, and they trivially preserve the pair of
  homogeneous submodules \((N, N') = (\top, \bot)\). Because \(R\) is generated by
  \(x_1, \dots, x_r\) and \(M\) is finitely generated over \(R\), \(M\) is finitely
  generated as a module over the polynomial ring \(\kappa[x_1, \dots, x_r]\) acting
  through the \(t_i\); this is the finiteness condition required below.

  \emph{Invoking the ambient subquotient induction.} Apply
  \cref{lem:graded_subquotient_isRatHilb} to the pair \((N, N') = (\top, \bot)\)
  with the \(r\) commuting degree-\((+1)\) endomorphisms \(t_1, \dots, t_r\). It
  yields \(\mathrm{IsRatHilb}\) of order \(r\) for the ambient subquotient Hilbert
  function of this pair (\cref{def:graded_subquotientHilb}),
  \[
    h(n) \;=\; \dim_\kappa(\top \cap \mathcal{M}_n) - \dim_\kappa(\bot \cap \mathcal{M}_n)
      \;=\; \dim_\kappa M_n - 0 \;=\; \dim_\kappa M_n .
  \]
  Hence \(\mathrm{IsRatHilb}(n \mapsto \dim_\kappa M_n,\, r)\) (\cref{def:ratHilb}):
  there are \(p \in \mathbb{Q}[X]\), \(d := r \in \mathbb{N}\) and \(N \in \mathbb{N}\)
  with
  \[
    \dim_\kappa M_n \;=\; \bigl[X^n\bigr]\bigl(p \cdot (1 - X)^{-d}\bigr)
    \qquad (n > N),
  \]
  which is the claim, with pole order \(d = r\). The internal mechanics of
  \cref{lem:graded_subquotient_isRatHilb} -- the base case via
  \cref{lem:ratHilb_ofEventuallyZero}, and the inductive step assembling the
  kernel/cokernel closure (\cref{lem:graded_subquotient_ker_coker}), the degreewise
  difference identity (\cref{lem:graded_subquotient_degreewise_diff}, built on
  \(D_5\), \cref{lem:graded_degreewise_finrank_diff}) and the finiteness transfer
  (\cref{lem:graded_subquotient_finite_transfer}) into
  \cref{lem:ratHilb_ofDiffEq} and \cref{lem:ratHilb_bump} -- are the Stacks~00K1
  induction realised entirely with ambient homogeneous submodules of \(M\).
\end{proof}

\subsection{The rationality toolkit \texorpdfstring{\(\mathrm{IsRatHilb}\)}{IsRatHilb}}
\label{subsec:isRatHilb}
% NOTE: Declarations in this subsection are currently \texttt{private} in the Lean
% source; \lean{} pins may not resolve until they are moved to a dedicated module.

The proof of \cref{lem:gradedHilbertSerre_rational} runs the Stacks~00K1 induction
on the number of degree-one generators. Each inductive step manipulates the
\emph{shape} of the Hilbert series -- a numerator polynomial over a power of
\((1-X)\) -- rather than the module itself. We isolate that bookkeeping in a
single predicate \(\mathrm{IsRatHilb}\) on functions \(f : \mathbb{N} \to
\mathbb{Q}\), together with the closure lemmas the induction consumes. These are
project-internal helpers realising standard steps of the Stacks~00K1 argument at
the level of power series; Mathlib supplies only the converse extraction
(\cref{lem:hilbertPoly_exists_mathlib}).

\begin{lemma}\leanok
[Coefficients of \((1-X)^{-1}\) times a series are partial sums]
  \label{lem:coeff_invOneSubPow_one_mul}
  \lean{AlgebraicGeometry.coeff_invOneSubPow_one_mul}
  \uses{lem:invOneSubPow_mathlib}
  Let \(F \in \mathbb{Q}\llbracket X\rrbracket\) be a power series. Then for every
  \(n \in \mathbb{N}\) the \(n\)-th coefficient of
  \(\operatorname{invOneSubPow}(\mathbb{Q}, 1)\cdot F = (1-X)^{-1}\cdot F\)
  (\cref{lem:invOneSubPow_mathlib}) is the partial sum of the coefficients of
  \(F\):
  \[
    \bigl[X^n\bigr]\bigl((1-X)^{-1}\cdot F\bigr)
    \;=\; \sum_{k=0}^{n} \bigl[X^k\bigr] F .
  \]
  Equivalently, multiplication by \((1-X)^{-1} = \sum_{m\ge 0} X^m\) is the
  summation operator on coefficient sequences. Proved directly in Lean from the
  identity \(\operatorname{invOneSubPow}(\mathbb{Q},1) = \sum_{m\ge 0} X^m\) and
  the Cauchy product.
\end{lemma}

\begin{lemma}\leanok
[Antidifference of a rational Hilbert series]
  \label{lem:rationalHilbert_antidiff}
  \lean{AlgebraicGeometry.rationalHilbert_antidiff}
  \uses{lem:coeff_invOneSubPow_one_mul, lem:invOneSubPow_mathlib}
  Let \(H, \delta : \mathbb{N} \to \mathbb{Q}\), let \(q \in \mathbb{Q}[X]\) and
  \(e, N \in \mathbb{N}\). Suppose that for all \(n > N\) the sequence \(\delta\)
  is the coefficient sequence of the rational series \(q\cdot(1-X)^{-e}\),
  \[
    \delta(n) \;=\; \bigl[X^n\bigr]\bigl(q\cdot(1-X)^{-e}\bigr),
  \]
  and that \(\delta\) is the first difference of \(H\) beyond \(N\), i.e.
  \(H(n+1) - H(n) = \delta(n+1)\) for all \(n > N\). Then there is a numerator
  polynomial \(p \in \mathbb{Q}[X]\) such that for all \(n > N\)
  \[
    H(n) \;=\; \bigl[X^n\bigr]\bigl(p\cdot(1-X)^{-(e+1)}\bigr).
  \]
  That is, summing a coefficient sequence (anti-differencing) raises the pole
  order by one: the partial-sum series of \(q\cdot(1-X)^{-e}\) equals
  \(p\cdot(1-X)^{-(e+1)}\) up to a polynomial correction in low degrees. This is
  the power-series heart of the Stacks~00K1 inductive step.
\end{lemma}

\begin{proof}
  \uses{lem:coeff_invOneSubPow_one_mul, lem:invOneSubPow_mathlib}
  \leanok
  Write \(F = q\cdot(1-X)^{-e}\). Multiplying by
  \((1-X)^{-1}\) and using the factorisation
  \((1-X)^{-(e+1)} = (1-X)^{-1}\cdot(1-X)^{-e}\)
  (\cref{lem:invOneSubPow_mathlib}), \cref{lem:coeff_invOneSubPow_one_mul} gives
  \(\bigl[X^m\bigr]\bigl(q\cdot(1-X)^{-(e+1)}\bigr) = \sum_{k=0}^{m}[X^k]F\), the
  partial sums of \(F\). On the other hand, telescoping the first-difference
  hypothesis from \(N+1\) expresses \(H(N+1+j)\) as \(H(N+1)\) plus a window sum
  of the coefficients of \(F\). Splitting the partial sum
  \(\sum_{k=0}^{n}[X^k]F\) at \(N+2\) and absorbing the resulting constant
  discrepancy into a numerator \(C\,(1-X)^{e}\) (a constant function is the
  order-\((e+1)\) coefficient sequence of \(C\,(1-X)^{e}\cdot(1-X)^{-(e+1)} =
  C\,(1-X)^{-1}\)), the explicit numerator
  \(p = C\,(1-X)^{e} + q\) realises \(H(n)\) as
  \([X^n]\bigl(p\cdot(1-X)^{-(e+1)}\bigr)\) for all \(n > N\).
\end{proof}

\begin{definition}\leanok
[Rational Hilbert function]
  \label{def:ratHilb}
  \lean{AlgebraicGeometry.IsRatHilb}
  \uses{lem:invOneSubPow_mathlib}
  \textit{Source: Stacks 00K1.}
  For a function \(f : \mathbb{N} \to \mathbb{Q}\) and \(d \in \mathbb{N}\), say
  that \(f\) is a \emph{rational Hilbert function of order \(d\)}, written
  \(\mathrm{IsRatHilb}(f, d)\), when there exist a numerator polynomial
  \(p \in \mathbb{Q}[X]\) and an \(N \in \mathbb{N}\) such that for all \(n > N\)
  \[
    f(n) \;=\; \bigl[X^n\bigr]\bigl(p\cdot(1-X)^{-d}\bigr)
    \;=\; \bigl[X^n\bigr]\bigl(\operatorname{toPowerSeries}(p)\cdot
      \operatorname{invOneSubPow}(\mathbb{Q}, d)\bigr)
  \]
  (\cref{lem:invOneSubPow_mathlib}). This is the ``eventually equal to the
  coefficient sequence of a rational series \(p\cdot(1-X)^{-d}\)'' predicate that
  the graded Hilbert--Serre induction tracks; its conclusion shape is exactly the
  one consumed by \cref{lem:hilbertPoly_exists_mathlib}.
\end{definition}

\begin{lemma}\leanok
[Base case: eventually-zero functions are rational of order \(0\)]
  \label{lem:ratHilb_ofEventuallyZero}
  \lean{AlgebraicGeometry.IsRatHilb.ofEventuallyZero}
  \uses{def:ratHilb}
  If \(f : \mathbb{N} \to \mathbb{Q}\) is eventually zero -- there is
  \(N\) with \(f(n) = 0\) for all \(n > N\) -- then \(\mathrm{IsRatHilb}(f, 0)\)
  (\cref{def:ratHilb}). Indeed \(f\) is then the coefficient sequence of the
  numerator polynomial \(0\) over \((1-X)^{0} = 1\). This is the base case of the
  Stacks~00K1 induction (no degree-one generators). Proved directly in Lean.
\end{lemma}

\begin{lemma}\leanok
[Raising the pole order]
  \label{lem:ratHilb_bump}
  \lean{AlgebraicGeometry.IsRatHilb.bump}
  \uses{def:ratHilb, lem:invOneSubPow_mathlib}
  If \(\mathrm{IsRatHilb}(f, d)\) (\cref{def:ratHilb}) then
  \(\mathrm{IsRatHilb}(f, d+1)\): multiplying the numerator by \((1-X)\) cancels
  one factor of \((1-X)^{-1}\), so the same coefficient sequence is realised at
  pole order \(d+1\). This lets two rational Hilbert functions of different
  orders be brought to a common denominator. Proved directly in Lean using
  \((1-X)\cdot(1-X)^{-(d+1)} = (1-X)^{-d}\) (\cref{lem:invOneSubPow_mathlib}).
\end{lemma}

\begin{lemma}\leanok
[Closure under difference]
  \label{lem:ratHilb_sub}
  \lean{AlgebraicGeometry.IsRatHilb.sub}
  \uses{def:ratHilb}
  If \(\mathrm{IsRatHilb}(f, d)\) and \(\mathrm{IsRatHilb}(g, d)\)
  (\cref{def:ratHilb}) at the same order \(d\), then
  \(\mathrm{IsRatHilb}(n \mapsto f(n) - g(n),\, d)\): the numerator of the
  difference is the difference of the numerators (over the common denominator
  \((1-X)^{-d}\)). Proved directly in Lean.
\end{lemma}

\begin{lemma}\leanok
[Closure under right shift]
  \label{lem:ratHilb_shiftRight}
  \lean{AlgebraicGeometry.IsRatHilb.shiftRight}
  \uses{def:ratHilb}
  If \(\mathrm{IsRatHilb}(f, d)\) (\cref{def:ratHilb}) then
  \(\mathrm{IsRatHilb}(n \mapsto f(n-1),\, d)\): the right shift on coefficient
  sequences is multiplication of the series by \(X\), so the numerator
  \(p\) is replaced by \(X\cdot p\) at the same pole order. Proved directly in
  Lean.
\end{lemma}

\begin{lemma}\leanok
[Antidifference step for the predicate]
  \label{lem:ratHilb_antidiff}
  \lean{AlgebraicGeometry.IsRatHilb.antidiff}
  \uses{def:ratHilb, lem:rationalHilbert_antidiff}
  Let \(H, g : \mathbb{N} \to \mathbb{Q}\) and \(e, N \in \mathbb{N}\). If
  \(\mathrm{IsRatHilb}(g, e)\) (\cref{def:ratHilb}) and \(g\) is the first
  difference of \(H\) beyond \(N\) (\(H(n+1) - H(n) = g(n+1)\) for all
  \(n > N\)), then \(\mathrm{IsRatHilb}(H, e+1)\). This is the predicate-level
  packaging of \cref{lem:rationalHilbert_antidiff}: a finite difference raises
  the pole order by one, the inductive heart of graded Hilbert--Serre.
\end{lemma}

\begin{proof}
  \uses{def:ratHilb, lem:rationalHilbert_antidiff}
  \leanok
  Unfolding \(\mathrm{IsRatHilb}(g, e)\) supplies a numerator \(q\) and a
  threshold \(N_g\) with \(g(n) = [X^n](q\cdot(1-X)^{-e})\) for \(n > N_g\).
  Applying \cref{lem:rationalHilbert_antidiff} with this \(q\) and threshold
  \(\max(N, N_g)\) -- whose two hypotheses are exactly the rational form of \(g\)
  and the first-difference identity for \(H\) -- yields a numerator \(p\) with
  \(H(n) = [X^n](p\cdot(1-X)^{-(e+1)})\) for \(n > \max(N, N_g)\), i.e.
  \(\mathrm{IsRatHilb}(H, e+1)\).
\end{proof}

\begin{lemma}\leanok
[Inductive-step engine from a degreewise short exact sequence]
  \label{lem:ratHilb_ofDiffEq}
  \lean{AlgebraicGeometry.IsRatHilb.ofDiffEq}
  \uses{def:ratHilb, lem:ratHilb_antidiff, lem:ratHilb_sub, lem:ratHilb_shiftRight}
  Let \(h_M, h_C, h_K : \mathbb{N} \to \mathbb{Q}\) and \(d, N \in \mathbb{N}\).
  Suppose \(\mathrm{IsRatHilb}(h_C, d)\) and \(\mathrm{IsRatHilb}(h_K, d)\)
  (\cref{def:ratHilb}), and that for all \(n > N\)
  \[
    h_M(n+1) - h_M(n) \;=\; h_C(n+1) - h_K(n).
  \]
  Then \(\mathrm{IsRatHilb}(h_M, d+1)\). This packages the entire power-series
  side of the Stacks~00K1 inductive step: \(h_M\) is the Hilbert function of
  \(M\), and \(h_C, h_K\) are the Hilbert functions of the cokernel
  \(C = M/xM\) and kernel \(K = \ker(x : M \to M(1))\) of multiplication by a
  degree-one element \(x\); the displayed identity is the alternating-sum
  consequence of the degreewise short exact sequence
  \(0 \to K_n \to M_n \xrightarrow{x} M_{n+1} \to C_{n+1} \to 0\). Given the
  inductive hypothesis that \(h_C, h_K\) are rational of order \(d\), the
  conclusion is rationality of \(h_M\) at order \(d+1\).
\end{lemma}

\begin{proof}
  \uses{lem:ratHilb_antidiff, lem:ratHilb_sub, lem:ratHilb_shiftRight}
  \leanok
  Set \(g(n) = h_C(n) - h_K(n-1)\). By closure under right shift
  (\cref{lem:ratHilb_shiftRight}) \(n \mapsto h_K(n-1)\) is rational of order
  \(d\), and by closure under difference (\cref{lem:ratHilb_sub}) so is \(g\). The
  hypothesis rewrites as \(h_M(n+1) - h_M(n) = g(n+1)\) for \(n > N\) (since
  \(g(n+1) = h_C(n+1) - h_K(n)\)). Applying the antidifference step
  (\cref{lem:ratHilb_antidiff}) to \(H = h_M\) and this \(g\) gives
  \(\mathrm{IsRatHilb}(h_M, d+1)\).
\end{proof}

\begin{theorem}[Hilbert polynomial from the section module]
  \label{thm:hilbertPoly_of_sectionModule}
  \lean{AlgebraicGeometry.Scheme.hilbertPolynomialOfSectionModule}
  \uses{lem:sectionGradedModule_fg, lem:gradedHilbertSerre_rational,
    lem:hilbertPoly_exists_mathlib}
  % SOURCE: [Nitsure], \S 1, sub-section "Stratification by Hilbert
  % Polynomials" (read from
  % references/nitsure-hilbert-quot-src/nitsure-hilbert-quot.tex, L468-L478).
  % SOURCE QUOTE:
  % "Let $X\to S$ be a finite type morphism of noetherian schemes,
  %  and let $L$ be a line bundle on $X$. Let $\F$
  %  be any coherent sheaf on
  %  $X$ whose schematic support is proper over $S$. Then for each $s\in S$,
  %  we get a polynomial ${\Phi}_s \in  \Q[\lambda]$
  %  which is the Hilbert polynomial of
  %  the restriction $\F_s = F|_{X_s}$ of $\F$ to the fiber $X_s$ over $s$,
  %  calculated with respect to the line bundle $L_s = L|_{X_s}$."
  \textit{Source: [Nitsure], \S 1 (FGA Explained Ch.~5).}
  Let \(X_s\) be proper over the field \(\kappa(s)\), let \(L_s\) be an ample line
  bundle on \(X_s\), and let \(\mathcal{F}_s\) be a coherent sheaf on \(X_s\).
  Then there is a unique polynomial \(\Phi_{\mathcal{F},s} \in \mathbb{Q}[\lambda]\)
  such that for all \(m \gg 0\)
  \[
    \Phi_{\mathcal{F},s}(m)
    \;=\; \dim_{\kappa(s)} \Gamma\bigl(X_s,\, \mathcal{F}_s \otimes L_s^{\otimes m}\bigr).
  \]
  This \(\Phi_{\mathcal{F},s}\) is the Hilbert polynomial of
  \cref{def:hilbert_polynomial}.
\end{theorem}

\begin{proof}
  \uses{lem:sectionGradedModule_fg, lem:gradedHilbertSerre_rational,
    lem:hilbertPoly_exists_mathlib}
  By \cref{lem:sectionGradedModule_fg} the section graded module
  \(M(X_s, \mathcal{F}_s, L_s)\) is a finitely generated graded module over the
  Noetherian graded ring \(R(X_s, L_s)\), with finite-dimensional components and
  degree-zero part the field \(\kappa(s)\). Applying the abstract graded
  Hilbert--Serre rationality lemma (\cref{lem:gradedHilbertSerre_rational}) to the
  graded \(\kappa(s)\)-algebra \(R = R(X_s, L_s)\) and the finitely generated
  graded module \(M = M(X_s, \mathcal{F}_s, L_s)\) exhibits the graded Hilbert
  function
  \(m \mapsto \dim_{\kappa(s)} \Gamma(X_s, \mathcal{F}_s \otimes L_s^{\otimes m})\)
  as the eventual coefficient function of a rational series: there are
  \(p \in \mathbb{Q}[X]\) and \(d \in \mathbb{N}\) and an \(N\) with
  \(\dim_{\kappa(s)} \Gamma(X_s, \mathcal{F}_s \otimes L_s^{\otimes m})
  = [X^m]\,(p \cdot (1 - X)^{-d})\) for \(m > N\). Apply
  \cref{lem:hilbertPoly_exists_mathlib} over the characteristic-zero field
  \(F = \mathbb{Q}\) to this pair \((p, d)\): it yields a unique polynomial
  \(h \in \mathbb{Q}[X]\) whose values \(h(n)\) eventually equal
  \([X^n]\,(p \cdot (1 - X)^{-d})\). Setting
  \(\Phi_{\mathcal{F},s} := h\) (with the variable renamed \(\lambda\)) gives a
  polynomial with \(\Phi_{\mathcal{F},s}(m) = \dim_{\kappa(s)}
  \Gamma(X_s, \mathcal{F}_s \otimes L_s^{\otimes m})\) for \(m \gg 0\). Uniqueness
  is the \(\exists!\) clause of \cref{lem:hilbertPoly_exists_mathlib}: any two
  polynomials agreeing with the Hilbert function for \(m \gg 0\) agree at
  infinitely many points, hence are equal.
\end{proof}

\subsection{The ambient subquotient induction for the Stacks~00K1 step}
\label{subsec:gradedModuleApi}

The power-series side of \cref{lem:gradedHilbertSerre_rational} is fully isolated
in the \(\mathrm{IsRatHilb}\) toolkit (\cref{subsec:isRatHilb}). What remains is
the \emph{graded-module} side of the Stacks~00K1 induction. Rather than build the
kernel and cokernel of multiplication by a degree-one element as graded modules
over a smaller quotient ring \(R/(x)\) -- which would force a graded structure on a
derived carrier -- we range the induction over \emph{pairs of homogeneous
submodules of a single fixed ambient graded module} \(M\). A pair \((N, N')\) with
\(N' \le N\) homogeneous submodules of \(M\) stands for the subquotient \(N/N'\); its
ambient Hilbert function is the difference \(n \mapsto \dim_\kappa(N \cap M_n) -
\dim_\kappa(N' \cap M_n)\). The class of such pairs, equipped with finitely many
pairwise-commuting degree-one endomorphisms preserving \(N\) and \(N'\), is
\emph{closed} under the kernel and cokernel of one of those endomorphisms, both
again presented as ambient subquotient pairs. Consequently no graded quotient ring,
graded quotient module, or regrading of a subtype/quotient carrier is ever formed:
every object that appears is an ambient family \(n \mapsto N_{\mathrm{aux}} \cap M_n\)
of submodules of the fixed \(M\).

This subsection records that machinery as a thin layer over existing Mathlib
scaffold: Mathlib already provides the notion of a homogeneous submodule
(\cref{lem:submodule_isHomogeneous_mathlib}), the homogeneity of a span of
homogeneous elements (\cref{lem:ideal_homogeneous_span_mathlib}), the degree-shift
of multiplication by a homogeneous element
(\cref{lem:isHomogeneousElem_graded_smul_mathlib}), the rank--nullity identity
(\cref{lem:finrank_range_add_finrank_ker_mathlib}), and the finite-generation
transfer along a surjection (\cref{lem:fg_restrictScalars_of_surjective_mathlib}).
What is supplied here is the internal direct-sum decomposition of a homogeneous
submodule as an ambient family (the split \(G_1\)), the subquotient Hilbert data,
the kernel/cokernel closure, the degreewise difference identity, the ambient
finiteness transfer, and the induction itself.

\emph{Setting for this subsection.} Let \(R = \bigoplus_{n \ge 0} R_n\) be a graded
ring whose degree-zero part \(R_0 = \kappa\) is a field and which is generated as a
\(\kappa\)-algebra in degree one, and let \(M = \bigoplus_{n \ge 0} M_n\) be a graded
\(\kappa\)-vector space (the fixed ambient module) with finite-dimensional
components \(M_n\). All submodules \(N, N', \dots\) below are homogeneous submodules
of \(M\), and the relevant endomorphisms are \(\kappa\)-linear maps \(M \to M\)
raising degree by one -- typically multiplication by a homogeneous element
\(x \in R_1\).

\subsubsection*{Mathlib dependency anchors}

\begin{lemma}[Homogeneous submodule]
  \label{lem:submodule_isHomogeneous_mathlib}
  \lean{Submodule.IsHomogeneous}
  \mathlibok
  \textit{Provided by Mathlib
  (\texttt{Mathlib.RingTheory.GradedAlgebra.Homogeneous.Submodule}).}
  Let \(R = \bigoplus_n R_n\) be a graded ring and \(M = \bigoplus_n M_n\) an
  internally graded \(R\)-module. A submodule \(p \subseteq M\) is
  \emph{homogeneous} when it is closed under taking homogeneous components:
  \(m \in p\) implies each graded piece \(\mathrm{proj}_n(m) \in p\). Equivalently
  \(p = \bigoplus_n (p \cap M_n)\) as a set. Mathlib bundles this into a homogeneous
  submodule type providing order and extensionality structure; the ambient internal
  direct-sum decomposition \(p = \bigoplus_n (M_n \cap p)\) is supplied by the split
  \(G_1\) (\cref{lem:graded_homogeneousSubmodule_iSupIndep},
  \cref{lem:graded_homogeneousSubmodule_iSup_eq}).
\end{lemma}

\begin{lemma}[A span of homogeneous elements is homogeneous]
  \label{lem:ideal_homogeneous_span_mathlib}
  \lean{Ideal.homogeneous_span}
  \mathlibok
  \textit{Provided by Mathlib
  (\texttt{Mathlib.RingTheory.GradedAlgebra.Homogeneous.Ideal}).}
  Let \(R = \bigoplus_n R_n\) be a graded ring and \(s \subseteq R\) a set each of
  whose elements is homogeneous. Then the ideal \(\mathrm{span}\,s\) is a
  homogeneous ideal. In particular, for \(x \in R_1\) the principal ideal \((x)\)
  is homogeneous, hence \(xM\) and the auxiliary submodule \(N' + x \cdot N\) are
  homogeneous; this homogeneity feeds the kernel/cokernel closure
  (\cref{lem:graded_subquotient_ker_coker}).
\end{lemma}

\begin{lemma}[Rank--nullity]
  \label{lem:finrank_range_add_finrank_ker_mathlib}
  \lean{LinearMap.finrank_range_add_finrank_ker}
  \mathlibok
  \textit{Provided by Mathlib
  (\texttt{Mathlib.LinearAlgebra.FiniteDimensional.Lemmas}).}
  Let \(\kappa\) be a field, \(V\) a finite-dimensional \(\kappa\)-vector space and
  \(\varphi : V \to W\) a \(\kappa\)-linear map. Then
  \[
    \dim_\kappa(\operatorname{range}\varphi) + \dim_\kappa(\ker\varphi)
      \;=\; \dim_\kappa V .
  \]
  This is the rank--nullity identity underlying the degreewise difference identity
  \(D_5\) (\cref{lem:graded_degreewise_finrank_diff}).
\end{lemma}

\begin{lemma}[Finite generation transfers along a surjection]
  \label{lem:fg_restrictScalars_of_surjective_mathlib}
  \lean{Submodule.FG.restrictScalars_of_surjective}
  \mathlibok
  \textit{Provided by Mathlib (\texttt{Mathlib.RingTheory.Finiteness.Basic}).}
  Let \(\varphi : R \to A\) be a surjective ring homomorphism and \(M\) an
  \(A\)-module that is finitely generated when viewed as an \(R\)-module via
  \(\varphi\). Then \(M\) is finitely generated as an \(A\)-module. Applied to the
  surjection \(\kappa[t_0, \dots, t_{r-1}] \twoheadrightarrow \kappa[t_0, \dots,
  t_{r-2}]\) killing the endomorphism \(x = t_{r-1}\) that annihilates the two
  subquotients, this transfers their finite generation down one generator
  (\cref{lem:graded_subquotient_finite_transfer}).
\end{lemma}

\begin{lemma}[Multiplication by a homogeneous element shifts degree]
  \label{lem:isHomogeneousElem_graded_smul_mathlib}
  \lean{SetLike.IsHomogeneousElem.graded_smul}
  \mathlibok
  \textit{Provided by Mathlib (\texttt{Mathlib.Algebra.GradedMulAction}).}
  Let \(\mathcal{A}\) be a graded family acting on a graded family \(\mathcal{M}\)
  via a graded scalar multiplication. If \(a \in R\) is homogeneous of degree
  \(i\) -- in particular \(x \in R_1\) -- then for \(m \in M_n\) one has
  \(a \cdot m \in M_{i+n}\). This is the degree-shifting property that makes \(xM\)
  a homogeneous submodule (with degree-\((n+1)\) part \(x \cdot M_n\)) and that
  underlies the ambient identity \(x \cdot N \cap M_{n+1} = x \cdot (N \cap M_n)\)
  used in the subquotient closure (\cref{lem:graded_subquotient_ker_coker}) and
  difference identity (\cref{lem:graded_subquotient_degreewise_diff}).
\end{lemma}

\begin{lemma}[Commutative subalgebra generated by commuting elements]
  \label{lem:adjoinCommRingOfComm_mathlib}
  \lean{Algebra.isMulCommutative_adjoin}
  \mathlibok
  \textit{Provided by Mathlib
  (\texttt{Mathlib.Algebra.Algebra.Subalgebra.Lattice}).}
  Let \(R\) be a commutative semiring and \(A\) an \(R\)-algebra, possibly
  noncommutative. If \(s \subseteq A\) is a set of pairwise-commuting elements
  (\(a b = b a\) for all \(a, b \in s\)), then the \(R\)-subalgebra
  \(\mathrm{adjoin}_R(s)\) generated by \(s\) has commutative multiplication; the
  commutative-ring structure on it follows. Applied to the \(r\) pairwise-commuting degree-raising endomorphisms
  \(x_0, \dots, x_{r-1} \in \operatorname{End}_\kappa(M)\), it yields a
  \emph{commutative} \(\kappa\)-subalgebra \(A = \mathrm{adjoin}_\kappa\{x_0, \dots,
  x_{r-1}\} \subseteq \operatorname{End}_\kappa(M)\); the polynomial action of
  \(\kappa[t_0, \dots, t_{r-1}]\) on \(M\) factors through this commutative \(A\)
  (\cref{lem:graded_subquotient_finite_transfer}).
\end{lemma}

\begin{lemma}[Finiteness transfers along a surjective linear map]
  \label{lem:module_finite_of_surjective_mathlib}
  \lean{Module.Finite.of_surjective}
  \mathlibok
  \textit{Provided by Mathlib (\texttt{Mathlib.RingTheory.Finiteness.Defs}).}
  Let \(A\) be a ring, \(M, N\) two \(A\)-modules and \(\varphi : M \to N\) a
  surjective \(A\)-linear map. If \(M\) is a finite \(A\)-module, then so is \(N\).
  This is the finiteness-transfer step at the heart of the ambient finiteness
  transfer (\cref{lem:graded_subquotient_finite_transfer}) and of the
  numerator/denominator finiteness lemmas
  (\cref{lem:graded_polyQuot_finite_of_le_numerator},
  \cref{lem:graded_polyQuot_finite_of_le_denominator}).
\end{lemma}

\begin{lemma}[Finiteness descends along an injective linear map into a Noetherian module]
  \label{lem:module_finite_of_injective_mathlib}
  \lean{Module.Finite.of_injective}
  \mathlibok
  \textit{Provided by Mathlib (\texttt{Mathlib.RingTheory.Noetherian.Defs}).}
  Let \(A\) be a ring, \(M, N\) two \(A\)-modules with \(N\) Noetherian (in
  particular, finite over a Noetherian ring \(A\)), and \(\varphi : M \to N\) an
  injective \(A\)-linear map. Then \(M\) is a finite \(A\)-module. This is the
  submodule-finiteness step in \cref{lem:graded_polyQuot_finite_of_le_numerator}.
\end{lemma}

\begin{lemma}[The lift of a linear map through a quotient]
  \label{lem:submodule_liftQ_mathlib}
  \lean{Submodule.liftQ}
  \mathlibok
  \textit{Provided by Mathlib (\texttt{Mathlib.LinearAlgebra.Quotient.Basic}).}
  Let \(A\) be a ring, \(M, N\) two \(A\)-modules, \(P \subseteq M\) a submodule and
  \(\varphi : M \to N\) an \(A\)-linear map with \(P \subseteq \ker\varphi\). Then
  \(\varphi\) factors uniquely through a linear map \(M/P \to N\) on the quotient.
  This realises the descent of a (\(\sigma\)-semilinear) map to the subquotient
  quotient in the ambient finiteness transfer
  (\cref{lem:graded_subquotient_finite_transfer}) and the denominator lemma
  (\cref{lem:graded_polyQuot_finite_of_le_denominator}).
\end{lemma}

\begin{lemma}[Only finitely many members of an independent family are nonzero]
  \label{lem:finite_ne_bot_of_iSupIndep_mathlib}
  \lean{Submodule.finite_ne_bot_of_iSupIndep}
  \mathlibok
  \textit{Provided by Mathlib
  (\texttt{Mathlib.LinearAlgebra.Dimension.Finite}).}
  Let \(\kappa\) be a field (more generally a division ring) and \(V\) a
  finite-dimensional \(\kappa\)-vector space. If \((p_i)_{i \in \iota}\) is an
  \(\mathrm{iSupIndep}\) family of subspaces of \(V\), then the set
  \(\{\,i : p_i \ne \bot\,\}\) is finite. This is the finiteness input that turns
  degreewise independence into eventual vanishing in the base case
  (\cref{lem:graded_subquotient_base_eventuallyZero}).
\end{lemma}

\begin{lemma}[The polynomial ring in finitely many variables is Noetherian]
  \label{lem:mvPolynomial_isNoetherianRing_fin_mathlib}
  \lean{MvPolynomial.isNoetherianRing_fin}
  \mathlibok
  \textit{Provided by Mathlib
  (\texttt{Mathlib.RingTheory.Polynomial.Basic}).}
  Let \(R\) be a Noetherian commutative ring and \(r \in \mathbb{N}\). Then the
  multivariate polynomial ring \(\operatorname{MvPolynomial}(\{0, \dots, r-1\}, R)\)
  is a Noetherian ring (Hilbert's basis theorem in finitely many variables). With
  \(R = \kappa\) a field this makes \(\kappa[t_0, \dots, t_{r-1}]\) Noetherian, the
  base ring over which the subquotient finiteness is measured
  (\cref{lem:graded_polyQuot_finite_of_le_numerator}).
\end{lemma}

\begin{lemma}[A finite module over a Noetherian ring is Noetherian]
  \label{lem:isNoetherian_of_isNoetherianRing_of_finite_mathlib}
  \lean{isNoetherian_of_isNoetherianRing_of_finite}
  \mathlibok
  \textit{Provided by Mathlib (\texttt{Mathlib.RingTheory.Noetherian.Defs}).}
  Let \(A\) be a Noetherian ring and \(M\) a finite \(A\)-module. Then \(M\) is a
  Noetherian \(A\)-module. Combined with
  \cref{lem:mvPolynomial_isNoetherianRing_fin_mathlib} this makes a finite
  \(\kappa[t_0, \dots, t_{r-1}]\)-module Noetherian, so its submodules are again
  finite (\cref{lem:graded_polyQuot_finite_of_le_numerator}).
\end{lemma}

\subsubsection*{The internal grading of a homogeneous submodule (split $G_1$)}

\begin{lemma}\leanok
[$G_1$ (independence) -- the family $n \mapsto M_n \cap p$ is independent]
  \label{lem:graded_homogeneousSubmodule_iSupIndep}
  \lean{AlgebraicGeometry.GradedModule.homogeneousSubmodule_inf_iSupIndep}
  \uses{lem:submodule_isHomogeneous_mathlib}
  Let \(M = \bigoplus_n M_n\) be an internally graded \(\kappa\)-module and let
  \(p \subseteq M\) be a homogeneous submodule
  (\cref{lem:submodule_isHomogeneous_mathlib}). Then the family
  \(n \mapsto M_n \cap p\) of \(\kappa\)-subspaces of \(M\) is independent: for each
  \(n\),
  \[
    (M_n \cap p) \;\cap\; \Bigl(\bigsqcup_{m \ne n} (M_m \cap p)\Bigr) \;=\; 0 .
  \]
  The whole statement lives in the ambient module \(M\): the subspaces
  \(M_n \cap p\) are submodules of \(M\), never of a derived carrier \(\smash{\hat p}\).
\end{lemma}

\begin{proof}
  \uses{lem:submodule_isHomogeneous_mathlib}
  \leanok
  The components \(M_n\) of the internal decomposition of \(M\) are independent. The
  subfamily \(n \mapsto M_n \cap p\) consists of submodules with \(M_n \cap p \le M_n\),
  so independence is inherited: a relation among the \(M_n \cap p\) is in particular a
  relation among the \(M_n\), forcing each term to vanish.
\end{proof}

\begin{lemma}\leanok
[$G_1$ (supremum) -- $\bigsqcup_n (M_n \cap p) = p$ for $p$ homogeneous]
  \label{lem:graded_homogeneousSubmodule_iSup_eq}
  \lean{AlgebraicGeometry.GradedModule.homogeneousSubmodule_iSup_inf_eq}
  \uses{lem:submodule_isHomogeneous_mathlib}
  Let \(p \subseteq M\) be a homogeneous submodule
  (\cref{lem:submodule_isHomogeneous_mathlib}). Then
  \[
    \bigsqcup_{n} (M_n \cap p) \;=\; p .
  \]
  Together with the independence of
  \cref{lem:graded_homogeneousSubmodule_iSupIndep}, this is the ambient internal
  direct sum \(p = \bigoplus_n (M_n \cap p)\), stated entirely inside \(M\) without
  requiring a direct-sum decomposition of the subtype \(\smash{\hat p}\).
\end{lemma}

\begin{proof}
  \uses{lem:submodule_isHomogeneous_mathlib}
  \leanok
  The inclusion \(\bigsqcup_n (M_n \cap p) \le p\) is clear, since each
  \(M_n \cap p \le p\). For the reverse inclusion, take \(m \in p\) and write
  \(m = \sum_n \mathrm{proj}_n(m)\) with \(\mathrm{proj}_n(m) \in M_n\) using the
  internal decomposition of \(M\). Homogeneity of \(p\) gives
  \(\mathrm{proj}_n(m) \in p\), hence \(\mathrm{proj}_n(m) \in M_n \cap p\); thus
  \(m\) is a finite sum of elements of the \(M_n \cap p\), i.e.
  \(m \in \bigsqcup_n (M_n \cap p)\).
\end{proof}

\subsubsection*{The ambient subquotient class and its Hilbert function}

\begin{definition}\leanok
[Subquotient Hilbert data]
  \label{def:graded_subquotientHilb}
  \lean{AlgebraicGeometry.GradedModule.subquotientHilb}
  \lean{AlgebraicGeometry.GradedModule.SubquotientDatum}
  \lean{AlgebraicGeometry.GradedModule.SubquotientDatum.hilb}
  \uses{lem:graded_homogeneousSubmodule_iSupIndep,
    lem:graded_homogeneousSubmodule_iSup_eq}
  Fix a graded \(\kappa\)-module \(\mathcal{M}\), i.e.\ \(M\) together with an
  internal decomposition \(M = \bigoplus_n \mathcal{M}_n\) whose components
  \(\mathcal{M}_n\) are finite-dimensional over \(\kappa\). A \emph{subquotient
  datum of length \(r\)} consists of:
  \begin{itemize}
    \item a pair of homogeneous submodules \(N' \le N \subseteq M\); and
    \item a family \(t : \{0, \dots, r-1\} \to (M \to_\kappa M)\) of pairwise
      commuting \(\kappa\)-linear endomorphisms, each raising degree by one
      (\(t_i(\mathcal{M}_n) \subseteq \mathcal{M}_{n+1}\)), each preserving \(N\)
      and \(N'\) (\(t_i N \subseteq N\), \(t_i N' \subseteq N'\)); together with
    \item a finiteness condition: the subquotient \(N/N'\) is a finite module over the
      polynomial ring \(\kappa[t_0, \dots, t_{r-1}]\) acting through the \(t_i\) (via the
      ambient \(t\)-stable carriers \(\mathrm{polySubmodule}\,N\),
      \(\mathrm{polySubmodule}\,N'\) of \cref{def:graded_polySubmodule}).
  \end{itemize}
  The pair \((N, N')\) represents the subquotient \(N/N'\). Its \emph{ambient
  Hilbert function} is
  \[
    \mathrm{hilb}(n) \;=\; \Bigl(\dim_\kappa (N \cap \mathcal{M}_n)
      \;-\; \dim_\kappa (N' \cap \mathcal{M}_n)\Bigr) \;\in\; \mathbb{Q},
  \]
  the difference of the (finite) \(\kappa\)-dimensions of the ambient intersections,
  computed in \(\mathbb{Z}\) and cast to \(\mathbb{Q}\). Since \(N' \le N\) one has
  \(N' \cap \mathcal{M}_n \le N \cap \mathcal{M}_n\), so
  \(\mathrm{hilb}(n) = \dim_\kappa\bigl((N \cap \mathcal{M}_n)/(N' \cap \mathcal{M}_n)\bigr)\)
  is the dimension of the degree-\(n\) piece of \(N/N'\). The split \(G_1\)
  (\cref{lem:graded_homogeneousSubmodule_iSupIndep},
  \cref{lem:graded_homogeneousSubmodule_iSup_eq}) guarantees that these ambient
  intersections assemble the homogeneous submodules \(N, N'\), so the data depends
  only on the ambient families \(n \mapsto N \cap \mathcal{M}_n\),
  \(n \mapsto N' \cap \mathcal{M}_n\).
\end{definition}

\subsubsection*{Ambient homogeneity calculus}

The kernel/cokernel closure and the degreewise difference identity rest on a small
calculus of how a degree-raising endomorphism interacts with the internal grading and
with the lattice of homogeneous submodules. Throughout, \(M = \bigoplus_n \mathcal{M}_n\)
is the fixed internally graded \(\kappa\)-module and \(x : M \to M\) is a
\(\kappa\)-linear endomorphism raising degree by one. None of these statements forms a
derived carrier: every object is a submodule of the fixed \(M\).

\begin{definition}\leanok
[Degree-raising endomorphism]
  \label{def:graded_raisesDegree}
  \lean{AlgebraicGeometry.GradedModule.RaisesDegree}
  Let \(M = \bigoplus_n \mathcal{M}_n\) be an internally graded \(\kappa\)-module. A
  \(\kappa\)-linear endomorphism \(x : M \to M\) \emph{raises degree by one} when it
  carries each graded piece into the next: \(x(\mathcal{M}_n) \subseteq
  \mathcal{M}_{n+1}\) for every \(n\). This is the grading-ring-free form of
  ``multiplication by a degree-one element'': it abstracts each of the \(r\)
  commuting degree-one generators \(x \in R_1\) acting on \(M\) into a single
  endomorphism, with no reference to the grading ring \(R\).
\end{definition}

\begin{lemma}\leanok
[Membership form of degree-raising]
  \label{lem:graded_raisesDegree_mem}
  \lean{AlgebraicGeometry.GradedModule.RaisesDegree.mem}
  \uses{def:graded_raisesDegree}
  If \(x\) raises degree by one (\cref{def:graded_raisesDegree}) and \(m \in
  \mathcal{M}_n\), then \(x m \in \mathcal{M}_{n+1}\).
\end{lemma}

\begin{proof}
  \uses{def:graded_raisesDegree}
  \leanok
  Immediate from the definition: \(x m \in x(\mathcal{M}_n) \subseteq
  \mathcal{M}_{n+1}\).
\end{proof}

\begin{lemma}\leanok
[Degree-shift commutation]
  \label{lem:graded_decompose_raisesDegree}
  \lean{AlgebraicGeometry.GradedModule.decompose_raisesDegree}
  \uses{def:graded_raisesDegree, lem:graded_raisesDegree_mem}
  Let \(x\) raise degree by one. Then for every \(m \in M\) and every \(i\) the
  degree-\((i+1)\) homogeneous component of \(x m\) equals \(x\) applied to the
  degree-\(i\) component of \(m\):
  \[
    \mathrm{proj}_{i+1}(x m) \;=\; x\bigl(\mathrm{proj}_i(m)\bigr) .
  \]
  This load-bearing commutation says a degree-raising endomorphism shifts the
  internal decomposition by exactly one slot; it is what makes preimages and images
  of homogeneous submodules under \(x\) homogeneous.
\end{lemma}

\begin{proof}
  \uses{def:graded_raisesDegree, lem:graded_raisesDegree_mem}
  \leanok
  Decompose \(m = \sum_j \mathrm{proj}_j(m)\) into its finitely many homogeneous
  components and apply \(x\): \(x m = \sum_j x(\mathrm{proj}_j(m))\) with
  \(x(\mathrm{proj}_j(m)) \in \mathcal{M}_{j+1}\) by
  \cref{lem:graded_raisesDegree_mem}. Reading off the degree-\((i+1)\) component, all
  terms with \(j+1 \ne i+1\) drop out, leaving \(x(\mathrm{proj}_i(m))\).
\end{proof}

\begin{lemma}\leanok
[Degree-raising kills the bottom]
  \label{lem:graded_decompose_raisesDegree_zero}
  \lean{AlgebraicGeometry.GradedModule.decompose_raisesDegree_zero}
  \uses{def:graded_raisesDegree}
  Let \(x\) raise degree by one. Then the degree-\(0\) component of \(x m\) vanishes
  for every \(m\): \(\mathrm{proj}_0(x m) = 0\).
\end{lemma}

\begin{proof}
  \uses{def:graded_raisesDegree}
  \leanok
  Writing \(x m = \sum_j x(\mathrm{proj}_j(m))\) with each summand in
  \(\mathcal{M}_{j+1}\), no summand has degree \(0\) (every \(j+1 \ge 1\)), so the
  degree-\(0\) component is \(0\).
\end{proof}

\begin{lemma}\leanok
[Preimage of a homogeneous submodule is homogeneous]
  \label{lem:graded_comap_isHomogeneous}
  \lean{AlgebraicGeometry.GradedModule.comap_isHomogeneous}
  \uses{lem:submodule_isHomogeneous_mathlib, lem:graded_decompose_raisesDegree}
  Let \(x\) raise degree by one and let \(N' \subseteq M\) be a homogeneous submodule
  (\cref{lem:submodule_isHomogeneous_mathlib}). Then the preimage \(x^{-1}(N') =
  \{m : x m \in N'\}\) is a homogeneous submodule of \(M\).
\end{lemma}

\begin{proof}
  \uses{lem:submodule_isHomogeneous_mathlib, lem:graded_decompose_raisesDegree}
  \leanok
  Let \(m \in x^{-1}(N')\) and fix a degree \(i\). By the degree-shift commutation
  (\cref{lem:graded_decompose_raisesDegree}), \(x(\mathrm{proj}_i(m)) =
  \mathrm{proj}_{i+1}(x m)\), which lies in \(N'\) because \(x m \in N'\) and \(N'\)
  is homogeneous. Hence \(\mathrm{proj}_i(m) \in x^{-1}(N')\), so \(x^{-1}(N')\) is
  homogeneous.
\end{proof}

\begin{lemma}\leanok
[Image of a homogeneous submodule is homogeneous]
  \label{lem:graded_map_isHomogeneous}
  \lean{AlgebraicGeometry.GradedModule.map_isHomogeneous}
  \uses{lem:submodule_isHomogeneous_mathlib, lem:graded_decompose_raisesDegree,
    lem:graded_decompose_raisesDegree_zero}
  Let \(x\) raise degree by one and let \(N \subseteq M\) be a homogeneous submodule.
  Then the image \(x \cdot N = \{x m : m \in N\}\) is a homogeneous submodule of
  \(M\), with degree-\((n+1)\) piece \(x(N \cap \mathcal{M}_n)\) and zero
  degree-\(0\) piece.
\end{lemma}

\begin{proof}
  \uses{lem:submodule_isHomogeneous_mathlib, lem:graded_decompose_raisesDegree,
    lem:graded_decompose_raisesDegree_zero}
    \leanok
  Take \(x m \in x \cdot N\) with \(m \in N\) and fix a degree. In degree \(0\) the
  component \(\mathrm{proj}_0(x m) = 0 \in x \cdot N\)
  (\cref{lem:graded_decompose_raisesDegree_zero}). In degree \(i+1\),
  \(\mathrm{proj}_{i+1}(x m) = x(\mathrm{proj}_i(m))\)
  (\cref{lem:graded_decompose_raisesDegree}); since \(N\) is homogeneous,
  \(\mathrm{proj}_i(m) \in N\), so this component lies in \(x \cdot N\). Thus every
  homogeneous component of \(x m\) lies in \(x \cdot N\), which is therefore
  homogeneous.
\end{proof}

\begin{lemma}\leanok
[Intersection of homogeneous submodules is homogeneous]
  \label{lem:graded_inf_isHomogeneous}
  \lean{AlgebraicGeometry.GradedModule.inf_isHomogeneous}
  \uses{lem:submodule_isHomogeneous_mathlib}
  If \(p, q \subseteq M\) are homogeneous submodules
  (\cref{lem:submodule_isHomogeneous_mathlib}) then so is their intersection \(p
  \cap q\). Mathlib provides no lattice-closure lemmas for homogeneous submodules, so
  this is recorded here.
\end{lemma}

\begin{proof}
  \uses{lem:submodule_isHomogeneous_mathlib}
  \leanok
  For \(m \in p \cap q\) and any degree \(i\), homogeneity of \(p\) and of \(q\)
  gives \(\mathrm{proj}_i(m) \in p\) and \(\mathrm{proj}_i(m) \in q\), hence
  \(\mathrm{proj}_i(m) \in p \cap q\).
\end{proof}

\begin{lemma}\leanok
[Sum of homogeneous submodules is homogeneous]
  \label{lem:graded_sup_isHomogeneous}
  \lean{AlgebraicGeometry.GradedModule.sup_isHomogeneous}
  \uses{lem:submodule_isHomogeneous_mathlib}
  If \(p, q \subseteq M\) are homogeneous submodules then so is their sum \(p + q\).
\end{lemma}

\begin{proof}
  \uses{lem:submodule_isHomogeneous_mathlib}
  \leanok
  Write an element of \(p + q\) as \(a + b\) with \(a \in p\), \(b \in q\). For any
  degree \(i\), \(\mathrm{proj}_i(a+b) = \mathrm{proj}_i(a) + \mathrm{proj}_i(b)\)
  with \(\mathrm{proj}_i(a) \in p\) and \(\mathrm{proj}_i(b) \in q\) by homogeneity,
  so \(\mathrm{proj}_i(a+b) \in p + q\).
\end{proof}

\begin{lemma}\leanok
[Ambient image identity]
  \label{lem:graded_map_inf_degree_eq}
  \lean{AlgebraicGeometry.GradedModule.map_inf_degree_eq}
  \uses{lem:graded_decompose_raisesDegree, lem:graded_map_isHomogeneous}
  Let \(x\) raise degree by one and let \(N \subseteq M\) be homogeneous. Then the
  degree-\((n+1)\) part of the image \(x \cdot N\) is exactly \(x\) of the
  degree-\(n\) part of \(N\):
  \[
    (x \cdot N) \cap \mathcal{M}_{n+1} \;=\; x \cdot (N \cap \mathcal{M}_n) .
  \]
\end{lemma}

\begin{proof}
  \uses{lem:graded_decompose_raisesDegree, lem:graded_map_isHomogeneous}
  \leanok
  \(\supseteq\): for \(m \in N \cap \mathcal{M}_n\), \(x m \in x \cdot N\) and \(x m
  \in \mathcal{M}_{n+1}\) since \(x\) raises degree. \(\subseteq\): an element \(x m
  \in x \cdot N\) lying in \(\mathcal{M}_{n+1}\) equals its own degree-\((n+1)\)
  component \(\mathrm{proj}_{n+1}(x m) = x(\mathrm{proj}_n(m))\)
  (\cref{lem:graded_decompose_raisesDegree}); homogeneity of \(N\) puts
  \(\mathrm{proj}_n(m) \in N \cap \mathcal{M}_n\), so \(x m \in x \cdot (N \cap
  \mathcal{M}_n)\). The image \(x \cdot N\) is homogeneous by
  \cref{lem:graded_map_isHomogeneous}.
\end{proof}

\begin{lemma}\leanok
[Ambient distributive law]
  \label{lem:graded_sup_inf_degree_eq}
  \lean{AlgebraicGeometry.GradedModule.sup_inf_degree_eq}
  \uses{lem:graded_inf_isHomogeneous, lem:graded_sup_isHomogeneous,
    lem:submodule_isHomogeneous_mathlib}
  For homogeneous submodules \(P, Q \subseteq M\) and any degree \(k\), intersecting
  the sum with a graded piece distributes:
  \[
    (P + Q) \cap \mathcal{M}_k \;=\; (P \cap \mathcal{M}_k) + (Q \cap \mathcal{M}_k) .
  \]
\end{lemma}

\begin{proof}
  \uses{lem:graded_inf_isHomogeneous, lem:graded_sup_isHomogeneous,
    lem:submodule_isHomogeneous_mathlib}
    \leanok
  \(\supseteq\) is the general lattice inequality. For \(\subseteq\), take \(z = p +
  q \in (P + Q) \cap \mathcal{M}_k\) with \(p \in P\), \(q \in Q\). Since \(z \in
  \mathcal{M}_k\), \(z = \mathrm{proj}_k(z) = \mathrm{proj}_k(p) +
  \mathrm{proj}_k(q)\); homogeneity of \(P\) and \(Q\)
  (\cref{lem:submodule_isHomogeneous_mathlib}) puts \(\mathrm{proj}_k(p) \in P \cap
  \mathcal{M}_k\) and \(\mathrm{proj}_k(q) \in Q \cap \mathcal{M}_k\), so \(z \in (P
  \cap \mathcal{M}_k) + (Q \cap \mathcal{M}_k)\). The intersection and sum stay in
  the homogeneous lattice by \cref{lem:graded_inf_isHomogeneous} and
  \cref{lem:graded_sup_isHomogeneous}.
\end{proof}

\subsubsection*{Closure under kernel and cokernel}

The kernel/cokernel closure of the subquotient class is realised by eight ambient
component facts: for a degree-raising endomorphism \(x\) and a homogeneous pair
\(N' \le N\), the kernel pair is \((N \cap x^{-1}(N'),\, N')\) and the cokernel pair is
\((N,\, N' + x \cdot N)\). The lemmas below record that both new pairs are homogeneous,
nest correctly, are annihilated by \(x\), and are preserved by any endomorphism \(t\)
commuting with \(x\) that preserves the original pair.

\begin{lemma}\leanok
[Homogeneity of the kernel pair's lower module]
  \label{lem:graded_ker_isHomogeneous}
  \lean{AlgebraicGeometry.GradedModule.ker_isHomogeneous}
  \uses{lem:graded_inf_isHomogeneous, lem:graded_comap_isHomogeneous}
  Let \(x\) raise degree by one (\cref{def:graded_raisesDegree}) and let \(N, N'\) be
  homogeneous submodules of \(M\). Then the kernel pair's lower module
  \(N \cap x^{-1}(N')\) is homogeneous. Indeed it is the intersection
  (\cref{lem:graded_inf_isHomogeneous}) of the homogeneous \(N\) with the homogeneous
  preimage \(x^{-1}(N')\) (\cref{lem:graded_comap_isHomogeneous}).
\end{lemma}

\begin{lemma}\leanok
[Homogeneity of the cokernel pair's upper module]
  \label{lem:graded_coker_isHomogeneous}
  \lean{AlgebraicGeometry.GradedModule.coker_isHomogeneous}
  \uses{lem:graded_sup_isHomogeneous, lem:graded_map_isHomogeneous}
  Under the same hypotheses, the cokernel pair's upper module \(N' + x \cdot N\) is
  homogeneous: it is the sum (\cref{lem:graded_sup_isHomogeneous}) of the homogeneous
  \(N'\) with the homogeneous image \(x \cdot N\)
  (\cref{lem:graded_map_isHomogeneous}).
\end{lemma}

\begin{lemma}\leanok
[The kernel pair nests]
  \label{lem:graded_ker_le}
  \lean{AlgebraicGeometry.GradedModule.ker_le}
  If \(N' \le N\) and \(x\) preserves \(N'\) (\(x \cdot N' \subseteq N'\)), then
  \(N' \le N \cap x^{-1}(N')\). Indeed \(N' \le N\) by hypothesis, and \(x \cdot N'
  \subseteq N'\) is exactly \(N' \subseteq x^{-1}(N')\); so \(N'\) is below the
  intersection. Hence the kernel pair \((N \cap x^{-1}(N'),\, N')\) is a valid
  subquotient pair (\cref{def:graded_subquotientHilb}).
\end{lemma}

\begin{lemma}\leanok
[The cokernel pair nests]
  \label{lem:graded_coker_le}
  \lean{AlgebraicGeometry.GradedModule.coker_le}
  If \(N' \le N\) and \(x\) preserves \(N\) (\(x \cdot N \subseteq N\)), then
  \(N' + x \cdot N \le N\): both summands lie in \(N\). Hence the cokernel pair
  \((N,\, N' + x \cdot N)\) is a valid subquotient pair
  (\cref{def:graded_subquotientHilb}).
\end{lemma}

\begin{lemma}\leanok
[$x$ annihilates the kernel subquotient]
  \label{lem:graded_ker_annihilate}
  \lean{AlgebraicGeometry.GradedModule.ker_annihilate}
  The endomorphism \(x\) carries the kernel pair's lower module into \(N'\):
  \(x \cdot (N \cap x^{-1}(N')) \subseteq N'\). Indeed \(N \cap x^{-1}(N') \subseteq
  x^{-1}(N')\), and applying \(x\) to a preimage of \(N'\) lands in \(N'\). Thus \(x\)
  acts as zero on the kernel subquotient.
\end{lemma}

\begin{lemma}\leanok
[$x$ annihilates the cokernel subquotient]
  \label{lem:graded_coker_annihilate}
  \lean{AlgebraicGeometry.GradedModule.coker_annihilate}
  The endomorphism \(x\) carries \(N\) into the cokernel pair's upper module:
  \(x \cdot N \subseteq N' + x \cdot N\), trivially (the right summand). Thus \(x\)
  acts as zero on the cokernel subquotient \(N/(N' + x \cdot N)\).
\end{lemma}

\begin{lemma}\leanok
[Commuting endomorphisms preserve the preimage]
  \label{lem:graded_comap_map_le_of_commute}
  \lean{AlgebraicGeometry.GradedModule.comap_map_le_of_commute}
  Let \(t\) commute with \(x\) and preserve \(N'\) (\(t \cdot N' \subseteq N'\)). Then
  \(t\) preserves the preimage: \(t \cdot x^{-1}(N') \subseteq x^{-1}(N')\). Indeed for
  \(m\) with \(x m \in N'\) one has \(x(t m) = t(x m) \in t \cdot N' \subseteq N'\), so
  \(t m \in x^{-1}(N')\). This is what shows the remaining \(t_i\) preserve the kernel
  pair's lower module \(N \cap x^{-1}(N')\).
\end{lemma}

\begin{lemma}\leanok
[Commuting endomorphisms preserve the image]
  \label{lem:graded_map_map_le_of_commute}
  \lean{AlgebraicGeometry.GradedModule.map_map_le_of_commute}
  Let \(t\) commute with \(x\) and preserve \(N\) (\(t \cdot N \subseteq N\)). Then
  \(t\) preserves the image: \(t \cdot (x \cdot N) \subseteq x \cdot N\). Indeed for
  \(m \in N\) one has \(t(x m) = x(t m)\) with \(t m \in N\), so \(t(x m) \in x \cdot
  N\). This is what shows the remaining \(t_i\) preserve the cokernel pair's upper
  module \(N' + x \cdot N\).
\end{lemma}

\begin{lemma}\leanok
[Kernel/cokernel closure of the subquotient class]
  \label{lem:graded_subquotient_ker_coker}
  \lean{AlgebraicGeometry.GradedModule.SubquotientDatum.ker}
  % This is a narrative grouping lemma: its content is realised in Lean by the eight
  % ambient component facts cited in \uses{} below (no single Lean declaration),
  % together with the SubquotientDatum.ker/.coker constructors that bundle them into a
  % length-(r-1) datum. The `\lean` pin names the `.ker` constructor — the single decl most
  % representative of the bundled closure — so this dependency-bearing grouping node (which is
  % `\uses`'d by the keystone rationality induction) stays wired into the graph rather than
  % orphaning the component facts it `\uses`. This is a deliberate duplicate of the
  % `def:graded_subquotientDatum_ker` pin.
  \uses{lem:graded_ker_isHomogeneous, lem:graded_coker_isHomogeneous,
    lem:graded_ker_le, lem:graded_coker_le, lem:graded_ker_annihilate,
    lem:graded_coker_annihilate, lem:graded_comap_map_le_of_commute,
    lem:graded_map_map_le_of_commute,
    def:graded_subquotientHilb, lem:graded_homogeneousSubmodule_iSupIndep,
    lem:graded_homogeneousSubmodule_iSup_eq, lem:ideal_homogeneous_span_mathlib,
    lem:isHomogeneousElem_graded_smul_mathlib, lem:graded_comap_isHomogeneous,
    lem:graded_map_isHomogeneous, lem:graded_inf_isHomogeneous,
    lem:graded_sup_isHomogeneous}
  Let \((N, N')\) be a subquotient datum (\cref{def:graded_subquotientHilb}) with
  endomorphisms \(t_0, \dots, t_{r-1}\), and set \(x = t_{r-1}\). Form the two new
  pairs of homogeneous submodules of \(M\)
  \[
    K \;=\; \bigl(\,N \cap x^{-1}(N'),\ N'\,\bigr),
    \qquad
    C \;=\; \bigl(\,N,\ N' + x \cdot N\,\bigr),
  \]
  where \(x^{-1}(N') = \{\,m \in M : x \cdot m \in N'\,\}\) is the preimage and
  \(x \cdot N\) is the image of \(N\) under \(x\). Then:
  \begin{enumerate}
    \item \(N \cap x^{-1}(N')\) and \(N' + x \cdot N\) are homogeneous submodules of
      \(M\), with \(N' \le N \cap x^{-1}(N')\) and \(N' + x \cdot N \le N\), so both
      \(K\) and \(C\) are subquotient pairs;
    \item the remaining endomorphisms \(t_0, \dots, t_{r-2}\) pairwise commute,
      raise degree by one, and preserve both members of \(K\) and of \(C\); and
    \item \(x\) annihilates both subquotients: \(x \cdot (N \cap x^{-1}(N')) \subseteq N'\)
      and \(x \cdot N \subseteq N' + x \cdot N\).
  \end{enumerate}
  Consequently \(K\) and \(C\) are again subquotient data, now of length \(r-1\)
  (the endomorphisms \(t_0, \dots, t_{r-2}\)), and \(K = \ker x\) and \(C =
  \operatorname{coker} x\) on the subquotient \(N/N'\): \(K\) represents
  \(\ker\bigl(x : N/N' \to N/N'\bigr)\) and \(C\) represents
  \(\operatorname{coker}\bigl(x : N/N' \to N/N'\bigr)\). No quotient or subtype
  carrier is constructed; \(K\) and \(C\) are ambient pairs of homogeneous
  submodules of the fixed \(M\).
\end{lemma}

\begin{proof}
  \uses{lem:graded_ker_isHomogeneous, lem:graded_coker_isHomogeneous,
    lem:graded_ker_le, lem:graded_coker_le, lem:graded_ker_annihilate,
    lem:graded_coker_annihilate, lem:graded_comap_map_le_of_commute,
    lem:graded_map_map_le_of_commute,
    lem:graded_homogeneousSubmodule_iSupIndep,
    lem:graded_homogeneousSubmodule_iSup_eq, lem:ideal_homogeneous_span_mathlib,
    lem:isHomogeneousElem_graded_smul_mathlib, lem:graded_comap_isHomogeneous,
    lem:graded_map_isHomogeneous, lem:graded_inf_isHomogeneous,
    lem:graded_sup_isHomogeneous}
    \leanok
  \emph{Homogeneity.} The preimage \(x^{-1}(N')\) is homogeneous: if
  \(m = \sum_n \mathrm{proj}_n(m)\) with \(x \cdot m \in N'\), then since \(x\) raises
  degree by one (\cref{lem:isHomogeneousElem_graded_smul_mathlib}) the components
  \(x \cdot \mathrm{proj}_n(m) \in \mathcal{M}_{n+1}\) lie in distinct degrees, and
  homogeneity of \(N'\) (split \(G_1\),
  \cref{lem:graded_homogeneousSubmodule_iSup_eq}) forces each
  \(x \cdot \mathrm{proj}_n(m) \in N'\), i.e.\ \(\mathrm{proj}_n(m) \in x^{-1}(N')\).
  Thus \(N \cap x^{-1}(N')\) is an intersection of homogeneous submodules, hence
  homogeneous. The image \(x \cdot N\) is homogeneous with degree-\((n+1)\) piece
  \(x \cdot (N \cap \mathcal{M}_n)\) by
  \cref{lem:isHomogeneousElem_graded_smul_mathlib}, and a sum of homogeneous
  submodules is homogeneous (\cref{lem:ideal_homogeneous_span_mathlib} applied at the
  module level), so \(N' + x \cdot N\) is homogeneous. The inclusions
  \(N' \le N \cap x^{-1}(N')\) (as \(x \cdot N' \subseteq N'\)) and
  \(N' + x \cdot N \le N\) (as \(N' \le N\) and \(x \cdot N \subseteq N\)) are
  immediate. In the language of the ambient homogeneity calculus, this paragraph is
  exactly: \(x^{-1}(N')\) is homogeneous (\cref{lem:graded_comap_isHomogeneous}),
  \(x \cdot N\) is homogeneous (\cref{lem:graded_map_isHomogeneous}), and the
  homogeneous lattice is closed under intersection
  (\cref{lem:graded_inf_isHomogeneous}) and sum
  (\cref{lem:graded_sup_isHomogeneous}).

  \emph{Preservation.} Each \(t_i\) with \(i \le r-2\) commutes with \(x\) and
  preserves \(N, N'\); hence it preserves \(x^{-1}(N')\) (if \(x t_i m = t_i x m \in
  N'\) whenever \(x m \in N'\)), so it preserves \(N \cap x^{-1}(N')\), and it
  preserves \(x \cdot N\) (as \(t_i x N = x t_i N \subseteq x N\)), so it preserves
  \(N' + x \cdot N\). Degree-raising and commutativity are inherited.

  \emph{Annihilation by \(x\).} For \(m \in N \cap x^{-1}(N')\) one has
  \(x \cdot m \in N'\) by definition, so \(x\) carries the first member of \(K\) into
  its second; and \(x \cdot N \subseteq N' + x \cdot N\) trivially. Thus \(x\) acts as
  zero on \(N/N'\) restricted to \(K\) and on the cokernel quotient \(C\).

  \emph{Identification with \(\ker x\) and \(\operatorname{coker} x\).} On the
  subquotient \(N/N'\), the class of \(m \in N\) is killed by \(x\) iff
  \(x \cdot m \in N'\), i.e.\ iff \(m \in N \cap x^{-1}(N')\); so \(\ker(x : N/N' \to
  N/N')\) is exactly \((N \cap x^{-1}(N'))/N'\), represented by \(K\). The image of
  \(x\) on \(N/N'\) is \((N' + x \cdot N)/N'\), so the cokernel
  \((N/N')/\operatorname{im} x = N/(N' + x \cdot N)\) is represented by \(C\).
\end{proof}

\subsubsection*{The subquotient degreewise difference identity ($D_6$)}

\begin{lemma}

\leanok
[preimage dimension equals ambient-intersection dimension]
  \label{lem:graded_finrank_comap_subtype}
  \lean{AlgebraicGeometry.GradedModule.finrank_comap_subtype}
  For submodules \(p, S\) of a \(\kappa\)-vector space \(M\), the dimension of the
  preimage of \(S\) under the inclusion \(p \hookrightarrow M\) equals the dimension of
  the ambient intersection: \(\dim_\kappa(\iota_p^{-1} S) = \dim_\kappa(p \cap S)\),
  where \(\iota_p : p \to M\) is the canonical injection. Project-local helper consumed
  by the two kernel-dimension computations of \cref{lem:graded_subquotient_degreewise_diff}.
\end{lemma}
\begin{proof}

\leanok
  The inclusion \(\iota_p\) is injective, so it restricts to a \(\kappa\)-linear
  isomorphism from \(\iota_p^{-1} S\) onto its image \(\iota_p(\iota_p^{-1} S) = p \cap S\)
  (image of a comap under an injective map). An isomorphism preserves dimension, giving
  the claimed equality.
\end{proof}

\begin{lemma}\leanok
[$D_6$ -- subquotient degreewise difference]
  \label{lem:graded_subquotient_degreewise_diff}
  \lean{AlgebraicGeometry.GradedModule.subquotient_degreewise_diff}
  \uses{def:graded_subquotientHilb, lem:graded_subquotient_ker_coker,
    lem:graded_degreewise_finrank_diff, lem:graded_homogeneousSubmodule_iSupIndep,
    lem:graded_homogeneousSubmodule_iSup_eq, lem:graded_finrank_comap_subtype,
    lem:isHomogeneousElem_graded_smul_mathlib, lem:graded_map_inf_degree_eq,
    lem:graded_sup_inf_degree_eq}
  % NOTE: The proof factors the two kernel-dimension computations through the
  %   helper \cref{lem:graded_finrank_comap_subtype} (preimage dimension = ambient
  %   intersection dimension).
  % SOURCE: [Stacks Project], "Commutative Algebra", \S "Noetherian graded
  % rings", Proposition \texttt{proposition-graded-hilbert-polynomial} (Tag
  % 00K1), proof (read from references/hilbert-serre-algebra.tex, L13939-L13947).
  \textit{Source: [Stacks Project], Tag 00K1 (degreewise short exact sequence).}
  Let \((N, N')\) be a subquotient datum (\cref{def:graded_subquotientHilb}) with
  Hilbert function \(h_M\), and let \(x = t_{r-1}\). Let \(h_K\) and \(h_C\) be the
  Hilbert functions of the kernel and cokernel subquotients
  \(K = (N \cap x^{-1}(N'),\, N')\) and \(C = (N,\, N' + x \cdot N)\) of
  \cref{lem:graded_subquotient_ker_coker}. Then for every \(n\)
  \[
    h_M(n+1) - h_M(n) \;=\; h_C(n+1) - h_K(n).
  \]
  Writing \(Q_n = (N \cap \mathcal{M}_n)/(N' \cap \mathcal{M}_n)\) for the
  degree-\(n\) piece of \(N/N'\) (so \(\dim_\kappa Q_n = h_M(n)\)), multiplication by
  \(x\) is a \(\kappa\)-linear map \(x : Q_n \to Q_{n+1}\) with
  \(\ker(x : Q_n \to Q_{n+1})\) the degree-\(n\) piece of \(K\) and
  \(Q_{n+1}/\operatorname{im} x\) the degree-\((n+1)\) piece of \(C\); the identity is
  then the degreewise rank--nullity difference \(D_5\)
  (\cref{lem:graded_degreewise_finrank_diff}). Specialising to \((N, N') = (\top,
  \bot)\) recovers the classical \(0 \to \mathcal{M}_n \xrightarrow{x}
  \mathcal{M}_{n+1} \to \mathcal{M}_{n+1}/x\mathcal{M}_n \to 0\).
\end{lemma}

% SOURCE QUOTE PROOF: [Stacks Project], Tag 00K1, proof, inductive step (read
% from references/hilbert-serre-algebra.tex, L13939-L13947).
% "Let $\overline{M} = M/xM$. Note that the map $x : M \to M$
%  is {\it not} a map of graded $S$-modules, since it does
%  not map $M_d$ into $M_d$. Namely, for each $d$ we have the
%  following short exact sequence
%  $$
%  0 \to M_d \xrightarrow{x} M_{d + 1} \to \overline{M}_{d + 1} \to 0
%  $$
%  This proves that $[M_{d + 1}] - [M_d] = [\overline{M}_{d + 1}]$.
%  Hence we win by Lemma \ref{lemma-numerical-polynomial}."
\begin{proof}
  \uses{lem:graded_subquotient_ker_coker, lem:graded_degreewise_finrank_diff,
    lem:graded_homogeneousSubmodule_iSup_eq,
    lem:isHomogeneousElem_graded_smul_mathlib, lem:graded_map_inf_degree_eq,
    lem:graded_sup_inf_degree_eq}
    \leanok
  Since \(N' \le N\) are homogeneous, the split \(G_1\)
  (\cref{lem:graded_homogeneousSubmodule_iSup_eq}) identifies the degree-\(n\) piece
  of \(N/N'\) with \(Q_n = (N \cap \mathcal{M}_n)/(N' \cap \mathcal{M}_n)\), a
  finite-dimensional \(\kappa\)-vector space of dimension \(h_M(n)\). Because \(x\)
  raises degree by one and preserves \(N\) and \(N'\)
  (\cref{lem:isHomogeneousElem_graded_smul_mathlib}), it induces a \(\kappa\)-linear
  map \(x : Q_n \to Q_{n+1}\).

  \emph{Kernel.} A class \([m] \in Q_n\) (with \(m \in N \cap \mathcal{M}_n\)) lies in
  \(\ker(x : Q_n \to Q_{n+1})\) iff \(x \cdot m \in N' \cap \mathcal{M}_{n+1}\), i.e.\
  iff \(m \in (N \cap x^{-1}(N')) \cap \mathcal{M}_n\). Hence \(\ker(x : Q_n \to
  Q_{n+1}) = \bigl((N \cap x^{-1}(N')) \cap \mathcal{M}_n\bigr)/(N' \cap \mathcal{M}_n)\)
  is exactly the degree-\(n\) piece of \(K\), of dimension \(h_K(n)\)
  (\cref{lem:graded_subquotient_ker_coker}).

  \emph{Cokernel.} The image of \(x : Q_n \to Q_{n+1}\) is
  \(\bigl(x \cdot (N \cap \mathcal{M}_n) + (N' \cap \mathcal{M}_{n+1})\bigr)/(N' \cap
  \mathcal{M}_{n+1})\). Using the ambient identity \((N' + x \cdot N) \cap
  \mathcal{M}_{n+1} = (N' \cap \mathcal{M}_{n+1}) + x \cdot (N \cap \mathcal{M}_n)\)
  -- which is the ambient distributive law (\cref{lem:graded_sup_inf_degree_eq})
  applied to \(P = N'\), \(Q = x \cdot N\), combined with the ambient image identity
  \((x \cdot N) \cap \mathcal{M}_{n+1} = x \cdot (N \cap \mathcal{M}_n)\)
  (\cref{lem:graded_map_inf_degree_eq}) -- the cokernel \(Q_{n+1}/\operatorname{im} x = (N \cap
  \mathcal{M}_{n+1})/\bigl((N' + x \cdot N) \cap \mathcal{M}_{n+1}\bigr)\) is the
  degree-\((n+1)\) piece of \(C\), of dimension \(h_C(n+1)\).

  \emph{Difference.} Apply the degreewise rank--nullity difference \(D_5\)
  (\cref{lem:graded_degreewise_finrank_diff}) to \(\varphi = (x : Q_n \to Q_{n+1})\),
  with \(V = Q_n\), \(W = Q_{n+1}\):
  \[
    h_M(n+1) - h_M(n)
      = \dim_\kappa Q_{n+1} - \dim_\kappa Q_n
      = \dim_\kappa(Q_{n+1}/\operatorname{im}\varphi) - \dim_\kappa(\ker\varphi)
      = h_C(n+1) - h_K(n).
  \]
  This is the \(\mathrm{hdiff}\) hypothesis consumed by \cref{lem:ratHilb_ofDiffEq}.
\end{proof}

\subsubsection*{The free polynomial-ring module structure}

The finiteness condition of a subquotient datum (\cref{def:graded_subquotientHilb}) and
its transfer down one generator (\cref{lem:graded_subquotient_finite_transfer}) are
phrased over the \emph{free} polynomial ring \(\kappa[t_0, \dots, t_{r-1}] =
\operatorname{MvPolynomial}(\{0, \dots, r-1\}, \kappa)\). The following blocks build the
action of that ring on the ambient module \(M\) from the \(r\) pairwise-commuting
degree-raising endomorphisms \(t_i\), keeping every carrier an ambient submodule of
\(M\). The free polynomial ring -- rather than the subalgebra
\(\mathrm{adjoin}_\kappa\{t_0, \dots, t_{r-1}\} \subseteq \operatorname{End}_\kappa(M)\)
-- is used precisely so the inductive transfer has the ring surjection
\(\kappa[t_0, \dots, t_{r-1}] \twoheadrightarrow \kappa[t_0, \dots, t_{r-2}]\) available
(\cref{lem:graded_subquotient_finite_transfer}).

\begin{definition}\leanok
[Polynomial evaluation ring homomorphism]
  \label{def:graded_polyEndHom}
  \lean{AlgebraicGeometry.GradedModule.polyEndHom}
  \uses{lem:adjoinCommRingOfComm_mathlib}
  Let \(t : \{0, \dots, r-1\} \to \operatorname{End}_\kappa(M)\) be a family of
  pairwise-commuting \(\kappa\)-linear endomorphisms. The \emph{polynomial evaluation
  ring homomorphism}
  \[
    \mathrm{polyEndHom}(t) : \kappa[t_0, \dots, t_{r-1}] \longrightarrow
      \operatorname{End}_\kappa(M)
  \]
  sends each variable \(X_i\) to \(t_i\). It is constructed by evaluating into the
  \emph{commutative} \(\kappa\)-subalgebra \(A = \mathrm{adjoin}_\kappa(\{t_i\})
  \subseteq \operatorname{End}_\kappa(M)\) -- commutative by
  \cref{lem:adjoinCommRingOfComm_mathlib} because the generators commute -- and then
  including \(A \hookrightarrow \operatorname{End}_\kappa(M)\); evaluation into the
  noncommutative \(\operatorname{End}_\kappa(M)\) directly is not an algebra
  homomorphism, which is why the commutative \(A\) is interposed.
\end{definition}

\begin{lemma}\leanok
[Evaluation sends $X_i$ to $t_i$]
  \label{lem:graded_polyEndHom_X}
  \lean{AlgebraicGeometry.GradedModule.polyEndHom_X}
  \uses{def:graded_polyEndHom}
  \(\mathrm{polyEndHom}(t)(X_i) = t_i\) for each \(i\). Immediate from the construction
  (\cref{def:graded_polyEndHom}), as evaluation sends \(X_i\) to the chosen image
  \(t_i\).
\end{lemma}

\begin{lemma}\leanok
[Evaluation sends a constant to a scalar]
  \label{lem:graded_polyEndHom_C}
  \lean{AlgebraicGeometry.GradedModule.polyEndHom_C}
  \uses{def:graded_polyEndHom}
  \(\mathrm{polyEndHom}(t)(C\,c) = c \cdot \mathrm{id}_M\) for each \(c \in \kappa\):
  the evaluation homomorphism (\cref{def:graded_polyEndHom}) is \(\kappa\)-linear and
  carries the constant polynomial \(C\,c\) to the scalar endomorphism \(c \cdot 1\).
\end{lemma}

\begin{definition}\leanok
[Polynomial-ring module structure on $M$]
  \label{def:graded_polyModule}
  \lean{AlgebraicGeometry.GradedModule.polyModule}
  \uses{def:graded_polyEndHom}
  The \emph{polynomial-module structure} is the \(\kappa[t_0, \dots, t_{r-1}]\)-module
  structure on \(M\) obtained by restricting scalars along
  \(\mathrm{polyEndHom}(t)\) (\cref{def:graded_polyEndHom}): a polynomial \(p\) acts on
  \(m \in M\) by \(p \cdot m = \mathrm{polyEndHom}(t)(p)(m)\). This is the module over the
  free polynomial ring on which the ambient subquotient induction's finiteness condition
  is measured.
\end{definition}

\begin{lemma}\leanok
[$X_i$ acts as $t_i$]
  \label{lem:graded_polyModule_X_smul}
  \lean{AlgebraicGeometry.GradedModule.polyModule_X_smul}
  \uses{def:graded_polyModule, lem:graded_polyEndHom_X}
  In the polynomial-module structure (\cref{def:graded_polyModule}), \(X_i \cdot m =
  t_i(m)\) for all \(m \in M\). Immediate from
  \(\mathrm{polyEndHom}(t)(X_i) = t_i\) (\cref{lem:graded_polyEndHom_X}).
\end{lemma}

\begin{lemma}\leanok
[A constant acts as a scalar]
  \label{lem:graded_polyModule_C_smul}
  \lean{AlgebraicGeometry.GradedModule.polyModule_C_smul}
  \uses{def:graded_polyModule, lem:graded_polyEndHom_C}
  In the polynomial-module structure (\cref{def:graded_polyModule}), \((C\,c) \cdot m =
  c \cdot m\) for all \(c \in \kappa\), \(m \in M\). Immediate from
  \(\mathrm{polyEndHom}(t)(C\,c) = c \cdot \mathrm{id}_M\)
  (\cref{lem:graded_polyEndHom_C}).
\end{lemma}


\begin{definition}\leanok
[Polynomial-ring submodule from a $t$-stable submodule]
  \label{def:graded_polySubmodule}
  \lean{AlgebraicGeometry.GradedModule.polySubmodule}
  \uses{def:graded_polyModule, lem:graded_polyModule_X_smul,
    lem:graded_polyModule_C_smul}
  Let \(N \subseteq M\) be a \(\kappa\)-submodule that is stable under each \(t_i\)
  (\(t_i \cdot N \subseteq N\)). Then \(N\), with the \emph{same} ambient carrier, is a
  \(\kappa[t_0, \dots, t_{r-1}]\)-submodule of \(M\) for the polynomial-module structure
  (\cref{def:graded_polyModule}). Stability under the action of an arbitrary polynomial
  follows by induction on monomials: constants act as scalars
  (\cref{lem:graded_polyModule_C_smul}) and multiplication by \(X_i\) acts as \(t_i\)
  (\cref{lem:graded_polyModule_X_smul}), both of which preserve \(N\).
\end{definition}


\subsubsection*{Finiteness-transfer infrastructure}

The ambient finiteness transfer (\cref{lem:graded_subquotient_finite_transfer})
factors the polynomial action of the larger ring through a surjection onto the
smaller one. The following blocks build that surjection and the semilinearity
identity it rests on, and record the two elementary finiteness manipulations
(shrinking a numerator, enlarging a denominator) used to feed the kernel and
cokernel constructors.

\begin{definition}

\leanok
[The last-variable surjection of polynomial rings]
  \label{lem:graded_lastVarAlgHom}
  \lean{AlgebraicGeometry.GradedModule.lastVarAlgHom}
  \lean{AlgebraicGeometry.GradedModule.lastVarAlgHom_X_castSucc}
  \lean{AlgebraicGeometry.GradedModule.lastVarAlgHom_X_last}
  \lean{AlgebraicGeometry.GradedModule.lastVarAlgHom_C}
  \lean{AlgebraicGeometry.GradedModule.lastVarAlgHom_rename_castSucc}
  \lean{AlgebraicGeometry.GradedModule.lastVarAlgHom_surjective}
  \lean{AlgebraicGeometry.GradedModule.lastVarAlgHom_ringHomSurjective}
  \uses{lem:fg_restrictScalars_of_surjective_mathlib}
  For a field \(\kappa\) and \(r \in \mathbb{N}\), the \emph{last-variable
  homomorphism} is the \(\kappa\)-algebra map
  \[
    \mathrm{lastVarAlgHom} : \kappa[t_0, \dots, t_r]
      = \operatorname{MvPolynomial}(\{0,\dots,r\}, \kappa)
      \;\longrightarrow\;
      \operatorname{MvPolynomial}(\{0,\dots,r-1\}, \kappa) = \kappa[t_0, \dots, t_{r-1}]
  \]
  obtained by evaluation: it sends the last variable \(X_r\) to \(0\) and each
  earlier variable \(X_i\) (for \(0 \le i < r\)) to \(X_i\), and fixes constants.
  It is a left inverse of the algebra map
  \(\kappa[t_0, \dots, t_{r-1}] \to \kappa[t_0, \dots, t_r]\)
  induced by the canonical inclusion \(\{0, \dots, r-1\} \hookrightarrow \{0, \dots, r\}\)
  on variable indices, and is therefore surjective. This is the concrete
  ring surjection \(\kappa[t_0, \dots, t_r] \twoheadrightarrow
  \kappa[t_0, \dots, t_{r-1}]\) along which finiteness is transferred
  (\cref{lem:fg_restrictScalars_of_surjective_mathlib}).
\end{definition}

\begin{lemma}

\leanok
[A $t$-stable submodule is closed under the polynomial action]
  \label{lem:graded_polyEndHom_mem_of_stable}
  \lean{AlgebraicGeometry.GradedModule.polyEndHom_mem_of_stable}
  \uses{def:graded_polyEndHom, lem:graded_polyEndHom_X, lem:graded_polyEndHom_C}
  Let \(t_0, \dots, t_{r-1}\) be pairwise-commuting \(\kappa\)-linear endomorphisms
  of \(M\) and let \(P' \subseteq M\) be a \(\kappa\)-submodule stable under each
  \(t_i\) (\(t_i \cdot P' \subseteq P'\)). Then \(P'\) is closed under the action of
  every polynomial: for all \(p \in \kappa[t_0, \dots, t_{r-1}]\) and \(m \in P'\),
  \(\mathrm{polyEndHom}(t)(p)(m) \in P'\) (\cref{def:graded_polyEndHom}). The proof
  is induction on \(p\): a constant acts as a scalar (\cref{lem:graded_polyEndHom_C}),
  multiplication by a variable acts as the stabilising \(t_i\)
  (\cref{lem:graded_polyEndHom_X}), and sums are handled additively.
\end{lemma}

\begin{lemma}

\leanok
[Mod-$P'$ semilinearity of the polynomial action]
  \label{lem:graded_polyEndHom_lastVar_sub_mem}
  \lean{AlgebraicGeometry.GradedModule.polyEndHom_lastVar_sub_mem}
  \uses{def:graded_polyEndHom, lem:graded_polyEndHom_X, lem:graded_polyEndHom_C,
    lem:graded_polyEndHom_mem_of_stable, lem:graded_lastVarAlgHom}
  Let \(t_0, \dots, t_r\) be pairwise-commuting \(\kappa\)-linear endomorphisms of
  \(M\), and let \(P' \le P\) be \(\kappa\)-submodules each stable under every
  \(t_i\), with the last endomorphism \(x = t_r\) carrying \(P\) into \(P'\)
  (\(x \cdot P \subseteq P'\)). Then for every polynomial
  \(s \in \kappa[t_0, \dots, t_r]\) and every \(m \in P\),
  \[
    \mathrm{polyEndHom}(t)(s)(m)
      \;-\; \mathrm{polyEndHom}(t_0, \dots, t_{r-1})
        \bigl(\mathrm{lastVarAlgHom}(s)\bigr)(m)
      \;\in\; P' .
  \]
  In words: acting by \(s\) through the full family \(t\) agrees, modulo \(P'\), with
  first projecting away the last variable via \(\mathrm{lastVarAlgHom}\)
  (\cref{lem:graded_lastVarAlgHom}) and acting through the reduced family
  \((t_0, \dots, t_{r-1})\). This is the semilinearity identity that makes the
  \(\sigma\)-semilinear quotient map of \cref{lem:graded_subquotient_finite_transfer}
  well-defined. The proof is induction on \(s\): the case of an earlier variable
  \(s = X_i\) (\(i < r\)) is the inductive hypothesis after \(\mathrm{lastVarAlgHom}\)
  fixes it; the case of the last variable \(s = X_r\) annihilates the reduced term,
  and the full term lands in \(P'\) because \(x = t_r\) carries \(P\) into the
  \(t\)-stable \(P'\) (\cref{lem:graded_polyEndHom_mem_of_stable}).
\end{lemma}

\begin{lemma}

\leanok
[Shrinking the numerator preserves finiteness]
  \label{lem:graded_polyQuot_finite_of_le_numerator}
  \lean{AlgebraicGeometry.GradedModule.polyQuot_finite_of_le_numerator}
  \uses{def:graded_polyModule, def:graded_polySubmodule,
    lem:module_finite_of_injective_mathlib,
    lem:mvPolynomial_isNoetherianRing_fin_mathlib,
    lem:isNoetherian_of_isNoetherianRing_of_finite_mathlib}
  Let \(N_1 \le N_2\) and \(P'\) be \(\kappa\)-submodules of \(M\), each stable under
  the commuting family \(t_0, \dots, t_{r-1}\), and suppose the subquotient
  \(N_2/P'\) is finite over \(\kappa[t_0, \dots, t_{r-1}]\)
  (\cref{def:graded_polyModule}, \cref{def:graded_polySubmodule}). Then \(N_1/P'\) is
  also finite over \(\kappa[t_0, \dots, t_{r-1}]\). Indeed the inclusion
  \(N_1 \subseteq N_2\) induces an injective \(\kappa[t_0, \dots, t_{r-1}]\)-linear
  map \(N_1/P' \hookrightarrow N_2/P'\); since
  \(\kappa[t_0, \dots, t_{r-1}]\) is Noetherian
  (\cref{lem:mvPolynomial_isNoetherianRing_fin_mathlib}), the finite module
  \(N_2/P'\) is Noetherian
  (\cref{lem:isNoetherian_of_isNoetherianRing_of_finite_mathlib}), so its submodule
  \(N_1/P'\) is finite (\cref{lem:module_finite_of_injective_mathlib}). This supplies
  the finiteness witness of the kernel constructor
  (\cref{def:graded_subquotientDatum_ker}).
\end{lemma}

\begin{lemma}

\leanok
[Enlarging the denominator preserves finiteness]
  \label{lem:graded_polyQuot_finite_of_le_denominator}
  \lean{AlgebraicGeometry.GradedModule.polyQuot_finite_of_le_denominator}
  \uses{def:graded_polyModule, def:graded_polySubmodule,
    lem:module_finite_of_surjective_mathlib, lem:submodule_liftQ_mathlib}
  Let \(N\) and \(P_1 \le P_2\) be \(\kappa\)-submodules of \(M\), each stable under
  the commuting family \(t_0, \dots, t_{r-1}\), and suppose the subquotient
  \(N/P_1\) is finite over \(\kappa[t_0, \dots, t_{r-1}]\)
  (\cref{def:graded_polyModule}, \cref{def:graded_polySubmodule}). Then \(N/P_2\) is
  also finite over \(\kappa[t_0, \dots, t_{r-1}]\). Indeed the inclusion
  \(P_1 \subseteq P_2\) gives a surjection \(N/P_1 \twoheadrightarrow N/P_2\)
  (\cref{lem:submodule_liftQ_mathlib}), and finiteness transfers along it
  (\cref{lem:module_finite_of_surjective_mathlib}). This supplies the
  finiteness witness of the cokernel constructor
  (\cref{def:graded_subquotientDatum_coker}).
\end{lemma}

\subsubsection*{Ambient finiteness transfer}

\begin{lemma}\leanok
[Abstract single-pair finiteness transfer down one variable]
  \label{lem:graded_subquotient_finite_transfer}
  \lean{AlgebraicGeometry.GradedModule.subquotient_finite_transfer}
  \uses{def:graded_polyModule, def:graded_polySubmodule,
    lem:graded_lastVarAlgHom, lem:graded_polyEndHom_lastVar_sub_mem,
    lem:module_finite_of_surjective_mathlib, lem:submodule_liftQ_mathlib,
    lem:fg_restrictScalars_of_surjective_mathlib}
  Let \(t_0, \dots, t_r\) be pairwise-commuting \(\kappa\)-linear endomorphisms of
  \(M\), and let \(P' \le P\) be \(\kappa\)-submodules of \(M\), each stable under
  every \(t_i\). Suppose the last endomorphism \(x = t_r\) carries \(P\) into \(P'\)
  (\(x \cdot P \subseteq P'\)), and that the subquotient \(P/P'\) is finite over the
  free polynomial ring \(\kappa[t_0, \dots, t_r]\) for the polynomial-module structure
  of \cref{def:graded_polyModule} (with carriers the ambient \(t\)-stable submodules
  \(\mathrm{polySubmodule}\,P\), \(\mathrm{polySubmodule}\,P'\),
  \cref{def:graded_polySubmodule}). Then \(P/P'\) is finite over
  \(\kappa[t_0, \dots, t_{r-1}]\) for the reduced family \((t_0, \dots, t_{r-1})\).

  This is the abstract single-pair transfer: it does not mention the kernel/cokernel
  pairs, but is the engine applied to each of them by the kernel and cokernel
  constructors (\cref{def:graded_subquotientDatum_ker},
  \cref{def:graded_subquotientDatum_coker}). Because the last variable
  annihilates the subquotient, the \(\kappa[t_0, \dots, t_r]\)-action factors through
  the last-variable surjection \(\kappa[t_0, \dots, t_r] \twoheadrightarrow
  \kappa[t_0, \dots, t_{r-1}]\), which is what makes the transfer drop one variable.
\end{lemma}

\begin{proof}
  \uses{def:graded_polyModule, def:graded_polySubmodule,
    lem:graded_lastVarAlgHom, lem:graded_polyEndHom_lastVar_sub_mem,
    lem:module_finite_of_surjective_mathlib, lem:submodule_liftQ_mathlib}
    \leanok
  Write \(\sigma = \mathrm{lastVarAlgHom} : \kappa[t_0, \dots, t_r]
  \twoheadrightarrow \kappa[t_0, \dots, t_{r-1}]\) for the surjective ring
  homomorphism sending the last variable to \(0\) and each earlier variable \(X_i\)
  (\(i < r\)) to \(X_i\) (\cref{lem:graded_lastVarAlgHom}). Consider the map on numerators
  \[
    g : \mathrm{polySubmodule}\,P \;\longrightarrow\;
      \bigl(\mathrm{polySubmodule}\,P\bigr)\big/
        \bigl(\mathrm{polySubmodule}\,P'\bigr), \qquad y \longmapsto [\,y\,],
  \]
  where the source carries the \(\kappa[t_0, \dots, t_r]\)-module structure and the
  target the reduced \(\kappa[t_0, \dots, t_{r-1}]\)-module structure (both with the
  same ambient carrier \(P\), \cref{def:graded_polySubmodule}). The mod-\(P'\)
  semilinearity identity (\cref{lem:graded_polyEndHom_lastVar_sub_mem}) -- acting by
  \(s\) through the full family and acting by \(\sigma(s)\) through the reduced family
  agree modulo \(P'\), precisely because \(x = t_r\) carries \(P\) into \(P'\) --
  says exactly that \(g\) is \(\sigma\)-semilinear. The map \(g\) is surjective (it is
  the quotient projection on the common carrier), and it kills the denominator
  \(\mathrm{polySubmodule}\,P'\), so it descends through that submodule
  (\cref{lem:submodule_liftQ_mathlib}) to a surjective
  \(\kappa[t_0, \dots, t_{r-1}]\)-linear map
  \[
    \bigl(\mathrm{polySubmodule}\,P\bigr)\big/\bigl(\mathrm{polySubmodule}\,P'\bigr)
      \;\twoheadrightarrow\;
    \bigl(\mathrm{polySubmodule}\,P\bigr)\big/\bigl(\mathrm{polySubmodule}\,P'\bigr)
  \]
  from the \(\kappa[t_0, \dots, t_r]\)-finite subquotient onto the reduced one.
  Finiteness transfers along this surjection
  (\cref{lem:module_finite_of_surjective_mathlib}), giving finiteness of \(P/P'\)
  over \(\kappa[t_0, \dots, t_{r-1}]\).

  \emph{Why the free polynomial ring.} The smaller ring must be the \emph{free}
  polynomial ring \(\kappa[t_0, \dots, t_{r-1}]\), not the subalgebra
  \(\mathrm{adjoin}_\kappa\{t_0, \dots, t_{r-1}\} \subseteq
  \operatorname{End}_\kappa(M)\): relations among the endomorphisms inside
  \(\operatorname{End}_\kappa(M)\) would obstruct the polynomial-ring surjection
  \(\sigma\) that the scalar transfer requires
  (\cref{lem:fg_restrictScalars_of_surjective_mathlib}). Keeping free polynomial
  rings at every step keeps \(\sigma\) -- and hence the transfer -- available
  throughout. Everything is finite generation of ambient submodules of the fixed
  \(M\) and their quotients.
\end{proof}

\subsubsection*{Subquotient constructors}

The kernel/cokernel closure (\cref{lem:graded_subquotient_ker_coker}) and the
abstract finiteness transfer (\cref{lem:graded_subquotient_finite_transfer}) are
bundled into two constructors that take a length-\((r+1)\) subquotient datum to the
length-\(r\) kernel and cokernel data consumed by the induction. The two stability
lemmas below first record that the new pairs are stable under the \emph{whole}
endomorphism family (not merely the reduced one), which is what lets the over-\(\kappa[t_0,\dots,t_r]\)
finiteness inputs be stated before the variable is dropped.

\begin{lemma}

\leanok
[The kernel pair's lower module is stable under the full family]
  \label{lem:graded_ker_stable_full}
  \lean{AlgebraicGeometry.GradedModule.ker_stable_full}
  \uses{lem:graded_comap_map_le_of_commute}
  Let \((N, N')\) be a length-\((r+1)\) subquotient datum
  (\cref{def:graded_subquotientHilb}) with endomorphisms \(t_0, \dots, t_r\) and
  \(x = t_r\). Then the kernel pair's lower module \(N \cap x^{-1}(N')\) is stable
  under \emph{every} \(t_i\) (\(0 \le i \le r\)), not only the reduced family: each
  \(t_i\) preserves \(N\) and preserves the preimage \(x^{-1}(N')\) because it
  commutes with \(x\) and preserves \(N'\)
  (\cref{lem:graded_comap_map_le_of_commute}). This full stability is what allows the
  kernel subquotient's finiteness to be measured over \(\kappa[t_0, \dots, t_r]\)
  before the transfer drops the last variable.
\end{lemma}

\begin{lemma}

\leanok
[The cokernel pair's upper module is stable under the full family]
  \label{lem:graded_coker_stable_full}
  \lean{AlgebraicGeometry.GradedModule.coker_stable_full}
  \uses{lem:graded_map_map_le_of_commute}
  Under the same hypotheses, the cokernel pair's upper module \(N' + x \cdot N\) is
  stable under every \(t_i\) (\(0 \le i \le r\)): each \(t_i\) preserves \(N'\) and
  preserves the image \(x \cdot N\) because it commutes with \(x\) and preserves
  \(N\) (\cref{lem:graded_map_map_le_of_commute}), and a map preserves a sum of
  preserved submodules. This full stability supports the over-\(\kappa[t_0, \dots,
  t_r]\) finiteness input of the cokernel constructor.
\end{lemma}

\begin{definition}

\leanok
[Kernel constructor of a subquotient datum]
  \label{def:graded_subquotientDatum_ker}
  \lean{AlgebraicGeometry.GradedModule.SubquotientDatum.ker}
  \uses{def:graded_subquotientHilb, lem:graded_subquotient_ker_coker,
    lem:graded_ker_isHomogeneous, lem:graded_ker_le, lem:graded_ker_annihilate,
    lem:graded_comap_map_le_of_commute, lem:graded_ker_stable_full,
    lem:graded_subquotient_finite_transfer,
    lem:graded_polyQuot_finite_of_le_numerator}
  From a length-\((r+1)\) subquotient datum \((N, N')\) with endomorphisms
  \(t_0, \dots, t_r\) and \(x = t_r\) (\cref{def:graded_subquotientHilb}), the
  \emph{kernel constructor} produces the length-\(r\) datum
  \[
    \ker D \;=\; \bigl(\,N \cap x^{-1}(N'),\ N'\,\bigr)
  \]
  on the reduced family \((t_0, \dots, t_{r-1})\). Its
  non-finiteness fields are supplied by the ambient kernel calculus: homogeneity of
  \(N \cap x^{-1}(N')\) (\cref{lem:graded_ker_isHomogeneous}), the nesting
  \(N' \le N \cap x^{-1}(N')\) (\cref{lem:graded_ker_le}), and stability under the
  reduced family (\cref{lem:graded_ker_stable_full}). The finiteness witness
  is obtained by first shrinking the numerator from \(N\) to
  \(N \cap x^{-1}(N')\) (\cref{lem:graded_polyQuot_finite_of_le_numerator}) to get
  \(\kappa[t_0, \dots, t_r]\)-finiteness of the kernel subquotient, then transferring
  down one variable (\cref{lem:graded_subquotient_finite_transfer}), valid because
  \(x\) annihilates the kernel subquotient (\cref{lem:graded_ker_annihilate}). This
  realises the kernel half of \cref{lem:graded_subquotient_ker_coker} as a genuine
  length-\(r\) datum.
\end{definition}

\begin{definition}

\leanok
[Cokernel constructor of a subquotient datum]
  \label{def:graded_subquotientDatum_coker}
  \lean{AlgebraicGeometry.GradedModule.SubquotientDatum.coker}
  \uses{def:graded_subquotientHilb, lem:graded_subquotient_ker_coker,
    lem:graded_coker_isHomogeneous, lem:graded_coker_le, lem:graded_coker_annihilate,
    lem:graded_map_map_le_of_commute, lem:graded_coker_stable_full,
    lem:graded_subquotient_finite_transfer,
    lem:graded_polyQuot_finite_of_le_denominator}
  From the same length-\((r+1)\) datum, the \emph{cokernel constructor} produces the
  length-\(r\) datum
  \[
    \operatorname{coker} D \;=\; \bigl(\,N,\ N' + x \cdot N\,\bigr)
  \]
  on the reduced family. Its non-finiteness fields come from the ambient cokernel
  calculus: homogeneity of \(N' + x \cdot N\) (\cref{lem:graded_coker_isHomogeneous}),
  the nesting \(N' + x \cdot N \le N\) (\cref{lem:graded_coker_le}), and stability
  under the reduced family (\cref{lem:graded_coker_stable_full}). The finiteness field
  is obtained by enlarging the denominator from \(N'\) to \(N' + x \cdot N\)
  (\cref{lem:graded_polyQuot_finite_of_le_denominator}) to get
  \(\kappa[t_0, \dots, t_r]\)-finiteness of the cokernel subquotient, then
  transferring down one variable (\cref{lem:graded_subquotient_finite_transfer}),
  valid because \(x\) annihilates the cokernel subquotient
  (\cref{lem:graded_coker_annihilate}). This realises the cokernel half of
  \cref{lem:graded_subquotient_ker_coker} as a genuine length-\(r\) datum.
\end{definition}

\subsubsection*{The degreewise difference identity}

\begin{lemma}\leanok
[$D_5$ -- degreewise rank--nullity difference]
  \label{lem:graded_degreewise_finrank_diff}
  \lean{AlgebraicGeometry.GradedModule.degreewise_finrank_diff}
  \uses{lem:finrank_range_add_finrank_ker_mathlib,
    lem:finrank_ses_additive_mathlib}
  Let \(\varphi : V \to W\) be a \(\kappa\)-linear map between finite-dimensional
  \(\kappa\)-vector spaces. Then
  \[
    \dim_\kappa W - \dim_\kappa V
      \;=\; \dim_\kappa(W/\operatorname{range}\varphi) - \dim_\kappa(\ker\varphi).
  \]
  Applied degreewise with \(\varphi = (x : \mathcal{M}_n \to \mathcal{M}_{n+1})\),
  \(V = \mathcal{M}_n\), \(W = \mathcal{M}_{n+1}\) -- so that \(\ker\varphi\) is the
  degree-\(n\) piece \(K_n\) of the kernel subquotient and
  \(W/\operatorname{range}\varphi = \mathcal{M}_{n+1}/x \mathcal{M}_n\) the
  degree-\((n+1)\) piece \(C_{n+1}\) of the cokernel subquotient
  (\cref{lem:graded_subquotient_degreewise_diff}, via the split \(G_1\)
  \cref{lem:graded_homogeneousSubmodule_iSup_eq} and the degree shift
  \cref{lem:isHomogeneousElem_graded_smul_mathlib}) -- this gives, for every \(n\),
  \[
    \dim_\kappa \mathcal{M}_{n+1} - \dim_\kappa \mathcal{M}_n
      \;=\; \dim_\kappa C_{n+1} - \dim_\kappa K_n,
  \]
  which is the difference identity \(\mathrm{hdiff}\) consumed by
  \cref{lem:ratHilb_ofDiffEq}. No graded structure is used in proving the
  underlying vector-space identity -- it is pure rank--nullity.
\end{lemma}

\begin{proof}
  \uses{lem:finrank_range_add_finrank_ker_mathlib,
    lem:finrank_ses_additive_mathlib}
    \leanok
  By rank--nullity (\cref{lem:finrank_range_add_finrank_ker_mathlib}),
  \(\dim_\kappa V = \dim_\kappa(\ker\varphi) + \dim_\kappa(\operatorname{range}
  \varphi)\), so \(\dim_\kappa(\operatorname{range}\varphi) = \dim_\kappa V -
  \dim_\kappa(\ker\varphi)\). By additivity on the short exact sequence
  \(0 \to \operatorname{range}\varphi \to W \to W/\operatorname{range}\varphi \to 0\)
  (\cref{lem:finrank_ses_additive_mathlib}),
  \(\dim_\kappa(W/\operatorname{range}\varphi) = \dim_\kappa W -
  \dim_\kappa(\operatorname{range}\varphi) = \dim_\kappa W - \dim_\kappa V +
  \dim_\kappa(\ker\varphi)\). Rearranging gives
  \(\dim_\kappa W - \dim_\kappa V = \dim_\kappa(W/\operatorname{range}\varphi) -
  \dim_\kappa(\ker\varphi)\). The degreewise application identifies
  \(\ker(x : \mathcal{M}_n \to \mathcal{M}_{n+1})\) and
  \(\mathcal{M}_{n+1}/\operatorname{range}(x : \mathcal{M}_n \to \mathcal{M}_{n+1})\)
  with the degree-\(n\) piece of the kernel subquotient and the degree-\((n+1)\)
  piece of the cokernel subquotient (\cref{lem:graded_subquotient_degreewise_diff}),
  the latter because the degree-\((n+1)\) graded piece of the homogeneous submodule
  \(x\mathcal{M}\) is exactly \(x \mathcal{M}_n\)
  (\cref{lem:isHomogeneousElem_graded_smul_mathlib}).
\end{proof}

\subsubsection*{The base case ($r = 0$)}

\begin{lemma}\leanok
[Base-case finiteness over the empty polynomial ring]
  \label{lem:graded_finiteDimensional_of_mvPolynomial_isEmpty_finite}
  \lean{AlgebraicGeometry.GradedModule.finiteDimensional_of_mvPolynomial_isEmpty_finite}
  Let \(Q\) be a \(\kappa\)-vector space that is also a module over
  \(\operatorname{MvPolynomial}(\sigma, \kappa)\) for an \emph{empty} index set
  \(\sigma\) -- in particular \(\sigma = \{0, \dots, r-1\}\) with \(r = 0\) -- in a way
  compatible with the \(\kappa\)-action (scalar tower). If \(Q\) is module-finite over
  \(\operatorname{MvPolynomial}(\sigma, \kappa)\), then \(Q\) is finite-dimensional over
  \(\kappa\). Indeed for empty \(\sigma\) the polynomial ring is
  \(\operatorname{MvPolynomial}(\sigma, \kappa) \cong \kappa\), a finite \(\kappa\)-module,
  so finiteness of \(Q\) over it transfers to finiteness over \(\kappa\) along the scalar
  tower. This is the length-zero input to the ambient subquotient induction
  (\cref{lem:graded_subquotient_isRatHilb}): when \(r = 0\) the subquotient \(N/N'\) is a
  finite-dimensional \(\kappa\)-vector space, hence \(\mathrm{hilb}\) is eventually zero.
\end{lemma}

\begin{lemma}

\leanok
[independence of images modulo the kernel]
  \label{lem:graded_iSupIndep_map_of_mem_ker_sup}
  \lean{AlgebraicGeometry.GradedModule.iSupIndep_map_of_mem_ker_sup}
  Let \(\pi : A \to Q\) be a \(\kappa\)-linear map and \((g_i)_{i \in \iota}\) a family of
  subspaces of \(A\) that is \emph{independent modulo \(\ker\pi\)}: for each \(i\), every
  \(a \in g_i\) that also lies in \(\ker\pi + \sum_{j \ne i} g_j\) already lies in
  \(\ker\pi\). Then the image family \((\pi(g_i))_{i \in \iota}\) is independent
  (\(\mathrm{iSupIndep}\)) in \(Q\). Ring-agnostic lattice core of the base-case
  independence step below; applied to the quotient map \(N \twoheadrightarrow N/N'\).
\end{lemma}
\begin{proof}

\leanok
  By the disjointness criterion for independence it suffices to show, for each \(i\),
  that \(\pi(g_i) \cap \sum_{j \ne i} \pi(g_j) = 0\). Take \(q\) in that intersection;
  write \(q = \pi(a)\) with \(a \in g_i\), and (pushing the sum through \(\pi\), since
  \(\pi(\sum_{j\ne i} g_j) = \sum_{j\ne i}\pi(g_j)\)) \(q = \pi(b)\) with
  \(b \in \sum_{j \ne i} g_j\). Then \(a - b \in \ker\pi\), so
  \(a = (a-b) + b \in \ker\pi + \sum_{j \ne i} g_j\). By the hypothesis applied at \(i\),
  \(a \in \ker\pi\), whence \(q = \pi(a) = 0\).
\end{proof}

\begin{lemma}

\leanok
[Base case: the length-zero subquotient Hilbert function vanishes eventually]
  \label{lem:graded_subquotient_base_eventuallyZero}
  \lean{AlgebraicGeometry.GradedModule.subquotient_base_eventuallyZero}
  \uses{def:graded_subquotientHilb,
    lem:graded_finiteDimensional_of_mvPolynomial_isEmpty_finite,
    lem:graded_homogeneousSubmodule_iSupIndep, lem:graded_iSupIndep_map_of_mem_ker_sup,
    lem:finite_ne_bot_of_iSupIndep_mathlib, lem:ratHilb_ofEventuallyZero}
  % SOURCE: [Stacks Project], "Commutative Algebra", \S "Noetherian graded
  % rings", Proposition \texttt{proposition-graded-hilbert-polynomial} (Tag
  % 00K1), proof, base case (read from references/hilbert-serre-algebra.tex,
  % L13908-L13910).
  % SOURCE QUOTE:
  % "We prove this by induction on the minimal number of
  %  generators of $S_1$. If this number is $0$, then
  %  $M_n = 0$ for all $n \gg 0$ and the result holds."
  \textit{Source: [Stacks Project], Tag 00K1 (base case of the induction).}
  Let \((N, N')\) be a subquotient datum of length \(0\)
  (\cref{def:graded_subquotientHilb}): a homogeneous pair \(N' \le N\) of the fixed
  ambient graded module \(M = \bigoplus_n \mathcal{M}_n\) with \emph{no} endomorphisms,
  whose finiteness witness says the subquotient \(Q_0 := N/N'\) is module-finite over
  \(\operatorname{MvPolynomial}(\varnothing, \kappa)\). Then the ambient Hilbert
  function vanishes eventually: there is \(K \in \mathbb{N}\) with
  \(\mathrm{hilb}(n) = 0\) for all \(n > K\). Consequently
  \(\mathrm{IsRatHilb}(\mathrm{hilb}, 0)\) (\cref{lem:ratHilb_ofEventuallyZero}),
  which is the base case of the ambient subquotient induction
  (\cref{lem:graded_subquotient_isRatHilb}).
\end{lemma}

\begin{proof}
  \uses{lem:graded_finiteDimensional_of_mvPolynomial_isEmpty_finite,
    lem:graded_homogeneousSubmodule_iSupIndep,
    lem:finite_ne_bot_of_iSupIndep_mathlib, lem:ratHilb_ofEventuallyZero}
    \leanok
  Since the index set of variables is empty,
  \(\operatorname{MvPolynomial}(\varnothing, \kappa) \cong \kappa\), so the finiteness
  witness makes the subquotient \(Q_0 = N/N'\) a \emph{finite-dimensional}
  \(\kappa\)-vector space
  (\cref{lem:graded_finiteDimensional_of_mvPolynomial_isEmpty_finite}). For each
  degree \(n\), composing the inclusion of the homogeneous piece
  \(N \cap \mathcal{M}_n \hookrightarrow N\) with the quotient projection
  \(N \twoheadrightarrow Q_0\) gives a \(\kappa\)-linear map
  \(\psi_n : N \cap \mathcal{M}_n \to Q_0\), \(m \mapsto [m]\), whose image is the
  class of the degree-\(n\) homogeneous piece inside the subquotient.

  \emph{Step 1: independence of the images.} The family
  \(\bigl(\operatorname{range}(\psi_n)\bigr)_{n}\) is \(\mathrm{iSupIndep}\) in
  \(Q_0\). The reason is the ambient direct-sum grading of \(M\): the homogeneous
  pieces \(\mathcal{M}_n\) form an independent family
  (\cref{lem:graded_homogeneousSubmodule_iSupIndep}). An element of
  \(\bigsqcup_{j \ne n} \operatorname{range}(\psi_j)\) is a finite sum of classes of
  homogeneous elements of degrees \(j \ne n\); since \(N'\) is homogeneous, the
  degree-\(n\) homogeneous component of any representative of such a class lies in
  \(N'\), so the class has zero degree-\(n\) component in \(Q_0\). A nonzero class in
  \(\operatorname{range}(\psi_n)\), by contrast, is represented by a homogeneous
  element of degree \(n\) not in \(N'\), so has nonzero degree-\(n\) component.
  Comparing degree-\(n\) components forces the intersection
  \(\operatorname{range}(\psi_n) \cap \bigsqcup_{j \ne n} \operatorname{range}(\psi_j)
  = 0\); independence follows.
  % NOTE: ROUTE (b). Prove the iSupIndep by direct membership analysis in
  %   ⨆_{j ≠ n} range(ψ j), reading off the degree-n homogeneous component
  %   directly inside Q₀'s fixed κ-module structure.
  %   No outgoing linear map out of Q₀ is constructed, so no scalar ring other than κ
  %   is ever introduced.
  % NOTE: ROUTE (a). Building a κ-linear detector
  %   Φ : Q₀ → M/N', [m] ↦ [π_n m] (degree-n projection) via descent through the
  %   subquotient fails: the descended map carries the scalar action of Q₀'s base
  %   ring, but the target M/N' carries only a κ-module structure -- a scalar-ring
  %   incompatibility. Route (b) avoids this by working entirely within Q₀.

  \emph{Step 2: finiteness of the support.} Because \(Q_0\) is finite-dimensional and
  the ranges \(\operatorname{range}(\psi_n)\) are independent, only finitely many of
  them are nonzero: the set \(\{\,n : \operatorname{range}(\psi_n) \ne 0\,\}\) is
  finite (\cref{lem:finite_ne_bot_of_iSupIndep_mathlib}). In particular it is bounded
  above, by some \(K \in \mathbb{N}\).

  \emph{Step 3: eventual vanishing.} For \(n > K\) we have
  \(\operatorname{range}(\psi_n) = 0\), i.e.\ every homogeneous element of
  \(N \cap \mathcal{M}_n\) has class \(0\) in \(Q_0 = N/N'\); equivalently
  \(N \cap \mathcal{M}_n \le N'\), hence \(N \cap \mathcal{M}_n = N' \cap \mathcal{M}_n\)
  and
  \[
    \mathrm{hilb}(n) \;=\; \dim_\kappa(N \cap \mathcal{M}_n)
      - \dim_\kappa(N' \cap \mathcal{M}_n) \;=\; 0 .
  \]
  Thus \(\mathrm{hilb}\) is eventually zero, and
  \cref{lem:ratHilb_ofEventuallyZero} gives \(\mathrm{IsRatHilb}(\mathrm{hilb}, 0)\).
\end{proof}

\subsubsection*{The ambient subquotient induction}

\begin{lemma}\leanok
[Rationality of the ambient subquotient Hilbert function]
  \label{lem:graded_subquotient_isRatHilb}
  \lean{AlgebraicGeometry.GradedModule.subquotient_hilbertSeries_rational}
  \uses{def:graded_subquotientHilb, lem:graded_subquotient_ker_coker,
    def:graded_subquotientDatum_ker, def:graded_subquotientDatum_coker,
    lem:graded_subquotient_degreewise_diff,
    lem:graded_subquotient_finite_transfer, lem:ratHilb_ofEventuallyZero,
    lem:ratHilb_ofDiffEq, lem:ratHilb_bump,
    lem:graded_subquotient_base_eventuallyZero,
    lem:graded_finiteDimensional_of_mvPolynomial_isEmpty_finite}
  % SOURCE: [Stacks Project], "Commutative Algebra", \S "Noetherian graded
  % rings", Proposition \texttt{proposition-graded-hilbert-polynomial} (Tag
  % 00K1), proof (read from references/hilbert-serre-algebra.tex, L13908-L13914).
  \textit{Source: [Stacks Project], Tag 00K1 (induction on the number of
  degree-one generators).}
  Let \((N, N')\) be a subquotient datum of length \(r\)
  (\cref{def:graded_subquotientHilb}) with ambient Hilbert function
  \(\mathrm{hilb}\). Then \(\mathrm{IsRatHilb}(\mathrm{hilb}, r)\)
  (\cref{def:ratHilb}): the difference \(n \mapsto \dim_\kappa(N \cap \mathcal{M}_n)
  - \dim_\kappa(N' \cap \mathcal{M}_n)\) is, for \(n \gg 0\), the coefficient
  sequence of a rational series \(p \cdot (1 - X)^{-r}\) with \(p \in
  \mathbb{Q}[X]\). The induction is on the length \(r\), and every object that
  appears is an ambient pair of homogeneous submodules of the fixed graded module
  \(M\).
\end{lemma}

% SOURCE QUOTE PROOF: [Stacks Project], Tag 00K1, proof, induction set-up (read
% from references/hilbert-serre-algebra.tex, L13908-L13914).
% "We prove this by induction on the minimal number of
%  generators of $S_1$. If this number is $0$, then
%  $M_n = 0$ for all $n \gg 0$ and the result holds.
%  To prove the induction step, let $x\in S_1$
%  be one of a minimal set of generators, such that
%  the induction hypothesis applies to the
%  graded ring $S/(x)$."
\begin{proof}
  \uses{lem:graded_subquotient_ker_coker,
    def:graded_subquotientDatum_ker, def:graded_subquotientDatum_coker,
    lem:graded_subquotient_degreewise_diff,
    lem:graded_subquotient_finite_transfer, lem:ratHilb_ofEventuallyZero,
    lem:ratHilb_ofDiffEq, lem:ratHilb_bump,
    lem:graded_subquotient_base_eventuallyZero,
    lem:graded_finiteDimensional_of_mvPolynomial_isEmpty_finite}
    \leanok
  We induct on the length \(r\).

  \emph{Base case \(r = 0\).} A length-zero datum has a finiteness witness making the
  subquotient \(N/N'\) a finite-dimensional \(\kappa\)-vector space
  (\cref{lem:graded_finiteDimensional_of_mvPolynomial_isEmpty_finite}); its degreewise
  images form an independent family inside that finite-dimensional space, so only
  finitely many are nonzero and \(\mathrm{hilb}(n) = 0\) for \(n \gg 0\)
  (\cref{lem:graded_subquotient_base_eventuallyZero}). By
  \cref{lem:ratHilb_ofEventuallyZero} an eventually-zero function is rational of order
  \(0\), giving \(\mathrm{IsRatHilb}(\mathrm{hilb}, 0)\).

  \emph{Inductive step \(r \ge 1\).} Set \(x = t_{r-1}\). By the kernel/cokernel
  closure (\cref{lem:graded_subquotient_ker_coker}) the kernel subquotient
  \(K = (N \cap x^{-1}(N'),\, N')\) and the cokernel subquotient
  \(C = (N,\, N' + x \cdot N)\) are subquotient data of length \(r-1\), assembled by
  the kernel and cokernel constructors (\cref{def:graded_subquotientDatum_ker},
  \cref{def:graded_subquotientDatum_coker}) whose finiteness witnesses come from the
  abstract finiteness transfer (\cref{lem:graded_subquotient_finite_transfer}). The
  induction hypothesis applied
  to \(K\) and \(C\) gives \(\mathrm{IsRatHilb}(h_K, r-1)\) and
  \(\mathrm{IsRatHilb}(h_C, r-1)\) (\cref{def:ratHilb}). By the degreewise
  difference identity (\cref{lem:graded_subquotient_degreewise_diff}),
  \[
    \mathrm{hilb}(n+1) - \mathrm{hilb}(n) \;=\; h_C(n+1) - h_K(n)
    \qquad \text{for all } n .
  \]
  Feeding \(\mathrm{IsRatHilb}(h_C, r-1)\), \(\mathrm{IsRatHilb}(h_K, r-1)\) and this
  difference identity into the inductive-step engine
  (\cref{lem:ratHilb_ofDiffEq}) yields \(\mathrm{IsRatHilb}(\mathrm{hilb}, r)\). (The
  common order \(r-1\) for \(h_C\) and \(h_K\) is arranged, if necessary, by raising
  the lower one with \cref{lem:ratHilb_bump}.) This completes the induction.
\end{proof}

\section{Support and freeness predicates}
\label{sec:quot_predicates}

The Quot functor (\cref{def:quot_functor}) and the Grassmannian
(\cref{def:grassmannian_scheme}) classify quotient sheaves subject to two
geometric conditions: a coherent quotient must have \emph{schematic support
proper over the base}, and the Grassmannian's quotients must be \emph{locally
free of a fixed rank}. This section introduces both predicates.

\begin{definition}

\leanok
[Annihilator ideal sheaf of a sheaf of modules]
  \label{def:modules_annihilator}
  \lean{AlgebraicGeometry.Scheme.Modules.annihilator}
  % NOTE: The definition carries no \uses: the Lean construction is
  % `IdealSheafData.ofIdeals (fun U => Module.annihilator Γ(X,U) Γ(F,U))`, which is well-defined
  % unconditionally and does NOT depend on the localization bridges at definition time. The two
  % facts lem:annihilator_localization_eq_map and lem:qcoh_section_localization_basicOpen are needed
  % only for the downstream CHARACTERIZATION (the reverse inclusion / equality on affine opens, the
  % "annihilator forward direction"), which is not yet a blueprint block.
  % NOTE: The definition itself is formalized and axiom-clean; the inclusion direction
  % \(\mathrm{Ann}(\mathcal{F})(U) \subseteq \mathrm{Ann}_{\mathcal{O}_X(U)}(\mathcal{F}(U))\)
  % (\cref{lem:modules_annihilator_ideal_le}) is closed. The reverse inclusion, which identifies
  % the annihilator sheaf with the module-annihilator on affine opens, remains bridge-gated on
  % \cref{lem:qcoh_section_localization_basicOpen}.
  % SOURCE: [Nitsure], \S 1, sub-section "Stratification by Hilbert
  % Polynomials" (read from
  % references/nitsure-hilbert-quot-src/nitsure-hilbert-quot.tex, L468-L471).
  % SOURCE QUOTE:
  % "Let $X\to S$ be a finite type morphism of noetherian schemes,
  %  and let $L$ be a line bundle on $X$. Let $\F$
  %  be any coherent sheaf on
  %  $X$ whose schematic support is proper over $S$."
  \textit{Source: [Nitsure], \S 1 (FGA Explained Ch.~5).}
  Let \(X\) be a scheme and let \(\mathcal{F}\) be a sheaf of
  \(\mathcal{O}_X\)-modules that is \emph{quasi-coherent and of finite type}
  (over a locally noetherian base, so that on each affine open \(U\) the section
  module \(\mathcal{F}(U)\) is a finitely generated \(\mathcal{O}_X(U)\)-module).
  The \emph{annihilator ideal sheaf} of \(\mathcal{F}\) is the ideal sheaf
  \(\mathrm{Ann}(\mathcal{F})\) defined on affine opens: for each affine open
  \(U \subseteq X\),
  \[
    \mathrm{Ann}(\mathcal{F})(U)
    \;=\; \mathrm{Ann}_{\mathcal{O}_X(U)}\bigl(\mathcal{F}(U)\bigr)
    \;\subseteq\; \mathcal{O}_X(U),
  \]
  the annihilator of the \(\mathcal{O}_X(U)\)-module \(\mathcal{F}(U)\). These
  local data patch to a coherent ideal sheaf on \(X\) whose vanishing locus is
  the set-theoretic support of \(\mathcal{F}\).

  \emph{Well-definedness on basic opens.} The coherence condition for an ideal
  sheaf requires that, for an affine open \(U\) and \(f \in \mathcal{O}_X(U)\),
  the section \(\mathrm{Ann}(\mathcal{F})(U)\) restricts correctly to the basic
  open \(D(f) \subseteq U\), i.e. that
  \(\mathrm{Ann}_{\mathcal{O}_X(D(f))}(\mathcal{F}(D(f)))\) is the extension of
  \(\mathrm{Ann}_{\mathcal{O}_X(U)}(\mathcal{F}(U))\) along the localization map
  \(\mathcal{O}_X(U) \to \mathcal{O}_X(D(f))\). This is exactly where the two
  hypotheses enter. The bridge
  \cref{lem:qcoh_section_localization_basicOpen} supplies, for a quasi-coherent
  finite-type \(\mathcal{F}\), that \(\mathcal{O}_X(D(f)) = (\mathcal{O}_X(U))_f\)
  and that the restriction \(\mathcal{F}(U) \to \mathcal{F}(D(f))\) exhibits
  \(\mathcal{F}(D(f))\) as the localized module \((\mathrm{powers}\,f)^{-1}
  \mathcal{F}(U)\); the algebra engine
  \cref{lem:annihilator_localization_eq_map} then yields, using that
  \(\mathcal{F}(U)\) is finitely generated, the required identity
  \(\mathrm{Ann}_{(\mathcal{O}_X(U))_f}\bigl((\mathrm{powers}\,f)^{-1}
  \mathcal{F}(U)\bigr) = \mathrm{Ann}_{\mathcal{O}_X(U)}(\mathcal{F}(U)) \cdot
  (\mathcal{O}_X(U))_f\). Thus the local annihilators are compatible with basic-open
  restriction, which is the patching datum that makes \(\mathrm{Ann}(\mathcal{F})\)
  a well-defined ideal sheaf. The finite-type hypothesis is genuinely required:
  without finite generation of \(\mathcal{F}(U)\) the annihilator need not commute
  with localization, and the local data fail to patch.

  \emph{Note.} Mathlib does not carry an annihilator for sheaves of modules on a scheme.
\end{definition}

\begin{lemma}\leanok
[Annihilator ideal bounded by the module annihilator on affine opens]
  \label{lem:modules_annihilator_ideal_le}
  \lean{AlgebraicGeometry.Scheme.Modules.annihilator_ideal_le}
  \uses{def:modules_annihilator}
  Let \(X\) be a scheme, \(\mathcal{F}\) a sheaf of \(\mathcal{O}_X\)-modules and
  \(U \subseteq X\) an affine open. Then the section over \(U\) of the annihilator
  ideal sheaf \(\mathrm{Ann}(\mathcal{F})\) (\cref{def:modules_annihilator}) is
  contained in the module-annihilator of the sections of \(\mathcal{F}\) over
  \(U\):
  \[
    \mathrm{Ann}(\mathcal{F})(U)
    \;\subseteq\; \mathrm{Ann}_{\mathcal{O}_X(U)}\bigl(\mathcal{F}(U)\bigr).
  \]
  This is the inclusion direction available directly from the definition of
  \(\mathrm{Ann}(\mathcal{F})\); the reverse inclusion is the basic-open coherence
  statement gated on \cref{lem:qcoh_section_localization_basicOpen}. Proved directly
  in Lean. It supports the well-definedness of the schematic support
  (\cref{def:schematic_support}).
\end{lemma}


\begin{lemma}[Sections on a basic open localize the affine ring]
  \label{lem:isLocalization_basicOpen_mathlib}
  \lean{AlgebraicGeometry.IsAffineOpen.isLocalization_basicOpen}
  \mathlibok
  \textit{Provided by Mathlib
  (\texttt{Mathlib.AlgebraicGeometry.AffineScheme}).}
  Let \(X\) be a scheme, \(U \subseteq X\) an affine open, and
  \(f \in \Gamma(X, U)\). Then the restriction map
  \(\Gamma(X, U) \to \Gamma\bigl(X,\, D(f)\bigr)\) on the basic open
  \(D(f) = X.\mathrm{basicOpen}\,f\) exhibits \(\Gamma(X, D(f))\) as the
  localization of \(\Gamma(X, U)\) away from \(f\); that is, it is
  \(\mathrm{IsLocalization.Away}\,f\), so
  \(\Gamma(X, D(f)) = (\Gamma(X, U))_f\) as a \(\Gamma(X,U)\)-algebra.
\end{lemma}


% --- gap-2 supporting infrastructure (the named pieces of the gap-2 transport above) ---


% --- Pullback of a quasi-coherent sheaf along an open immersion is quasi-coherent ---


\begin{lemma}[Sheaf Mayer--Vietoris: the pullback cone over a two-open cover is a limit]
  \label{lem:isLimitPullbackCone_mathlib}
  \lean{TopCat.Sheaf.isLimitPullbackCone}
  \mathlibok
  \textit{Provided by Mathlib
  (\texttt{Mathlib.Topology.Sheaves.SheafCondition.PairwiseIntersections}).}
  For a sheaf \(F\) on a topological space and two opens \(U, V\), the square relating
  \(F(U \cup V)\), \(F(U)\), \(F(V)\), \(F(U \cap V)\) is a pullback (limit). This is the generic
  module Mayer--Vietoris used in the inductive step of the descent.
\end{lemma}


\begin{lemma}[Separatedness: a section is determined by its restrictions to a cover]
  \label{lem:eq_of_locally_eq_mathlib}
  \lean{TopCat.Sheaf.eq_of_locally_eq'}
  \mathlibok
  \textit{Provided by Mathlib
  (\texttt{Mathlib.Topology.Sheaves.SheafCondition.UniqueGluing}).}
  Let \(F\) be a sheaf on a topological space, \(V\) an open, and \(\{U_i\}\) a family of opens with
  \(\bigsqcup_i U_i = V\). If two sections \(s, t \in F(V)\) have equal restrictions
  \(s|_{U_i} = t|_{U_i}\) for every \(i\), then \(s = t\). This separatedness over a cover is the
  uniqueness half of the sheaf condition; it discharges the \(\mathrm{exists\_of\_eq}\) obligation
  (a global section vanishing on \(D(f)\) is killed by a power of \(f\)) and the final
  ``conclude on \(D(f)\)'' step of the cover-form descent
  (\cref{lem:section_localization_descent_of_cover}).
\end{lemma}

% NOTE: G1-core is now derived from gap1, not by a direct 3-field compact-open induction. The proof
% is the forward leg of the equivalence \cref{lem:isIso_fromTildeΓ_iff_isLocalizedModule_restrict}.
% The LEMMA STATEMENT is cited to Stacks 01HA above; the proof is a project-internal reduction (no
% transcribed Stacks proof), so there is no `% SOURCE QUOTE PROOF`.

% --- Quasi-coherence datum and the gap1 sub-build ---


% --- Mathlib dependency anchors for the gap1 slice-transport route ---

\begin{lemma}[Pullback of sheaves of modules along a morphism of schemes]
  \label{lem:modules_pullback_mathlib}
  \lean{AlgebraicGeometry.Scheme.Modules.pullback}
  \mathlibok
  \textit{Provided by Mathlib (\texttt{Mathlib.AlgebraicGeometry.Modules.Sheaf}).}
  For a morphism of schemes \(f : X \to Y\), Mathlib supplies the pullback functor
  \(f^{*} : Y.\mathrm{Modules} \to X.\mathrm{Modules}\) on sheaves of modules, left adjoint to
  pushforward, hence colimit-preserving, and carrying the unit sheaf to the unit sheaf. The project
  already uses this functor in the definition of local freeness. It is the geometric restriction
  through which a presentation is transported to an open subscheme.
\end{lemma}

\begin{lemma}[Restriction of sheaves of modules along an open immersion]
  \label{lem:modules_restrictFunctor_mathlib}
  \lean{AlgebraicGeometry.Scheme.Modules.restrictFunctor}
  \mathlibok
  \textit{Provided by Mathlib (\texttt{Mathlib.AlgebraicGeometry.Modules.Sheaf}).}
  For an open immersion \(f : X \to Y\), Mathlib supplies the restriction functor
  \(f^{\mathrm{res}} : Y.\mathrm{Modules} \to X.\mathrm{Modules}\), with
  \(\Gamma(f^{\mathrm{res}} M, U) = \Gamma(M, f(U))\) definitionally; it is a left adjoint (colimit
  preserving) and carries the unit sheaf to the unit sheaf. This is the geometric
  ``restrict to open \(U\)'' functor on \((\Spec R).\mathrm{Modules}\).
\end{lemma}

\begin{lemma}[Restriction is isomorphic to pullback]
  \label{lem:modules_restrictFunctorIsoPullback_mathlib}
  \lean{AlgebraicGeometry.Scheme.Modules.restrictFunctorIsoPullback}
  \mathlibok
  \textit{Provided by Mathlib (\texttt{Mathlib.AlgebraicGeometry.Modules.Sheaf}).}
  For an open immersion \(f\), the restriction functor (\cref{lem:modules_restrictFunctor_mathlib})
  is naturally isomorphic to the pullback functor (\cref{lem:modules_pullback_mathlib}):
  \(f^{\mathrm{res}} \cong f^{*}\). This identifies the restriction functor with the
  universal-property pullback, so a presentation transported along one transports along the other.
\end{lemma}


% --- The finite basic-open cover (topological precursor) ---


% --- The over-site/open-subspace sheaf equivalence (bridge C, topological layer) ---

% The blocks in this sub-section record the project-local infrastructure that fills the explicit
% TODO of Mathlib's `Topology/Sheaves/Over.lean`: the equivalence of sheaf categories induced by
% the site equivalence `Opens.overEquivalence U`. Together with the three Mathlib dependency
% anchors below they form the topos-theoretic (purely topological) layer of bridge C
% (\cref{lem:over_restrict_iso}); the geometric identification of the structure sheaves is the
% layer above them.

\begin{lemma}[Slice of opens equivalent to opens of the open subspace]
  \label{lem:opens_overEquivalence_mathlib}
  \lean{TopologicalSpace.Opens.overEquivalence}
  \mathlibok
  \textit{Provided by Mathlib (\texttt{Mathlib.Topology.Sheaves.Over}).}
  For a topological space \(X\) and an open \(U \subseteq X\), the slice category of opens of
  \(X\) over \(U\) (the opens of \(X\) contained in \(U\), with inclusions) is equivalent to the
  opens of the subspace \(U\):
  \((\mathrm{Opens}\,X)_{/U} \simeq \mathrm{Opens}(U)\). Mathlib supplies this underlying
  equivalence of sites; it leaves the lift to sheaf categories as an explicit TODO.
\end{lemma}

\begin{lemma}[Sheaf categories of equivalent sites are equivalent]
  \label{lem:equivalence_sheafCongr_mathlib}
  \lean{CategoryTheory.Equivalence.sheafCongr}
  \mathlibok
  \textit{Provided by Mathlib (\texttt{Mathlib.CategoryTheory.Sites.Equivalence}).}
  Let \(e : C \simeq D\) be an equivalence of categories carrying Grothendieck topologies \(J\)
  on \(C\) and \(K\) on \(D\), and suppose \(e.\mathrm{inverse}\) is a dense subsite for
  \((K, J)\). Then for any target category \(A\) the sheaf categories are equivalent:
  \(\mathrm{Sheaf}(J, A) \simeq \mathrm{Sheaf}(K, A)\).
\end{lemma}

\begin{lemma}[Module sheaves transport across an equivalence of ringed sites]
  \label{lem:pushforwardPushforwardEquivalence_mathlib}
  \lean{SheafOfModules.pushforwardPushforwardEquivalence}
  \mathlibok
  \textit{Provided by Mathlib
  (\texttt{Mathlib.Algebra.Category.ModuleCat.Sheaf.PushforwardContinuous}).}
  Let \(e : C \simeq D\) be an equivalence of sites with both \(e.\mathrm{functor}\) and
  \(e.\mathrm{inverse}\) continuous, and let \(S\) (a sheaf of rings on \(J\)) and \(R\) (a sheaf
  of rings on \(K\)) be related by mutually inverse comparison maps across the continuous
  pushforwards. Then the categories of sheaves of modules are equivalent,
  \(\mathrm{SheafOfModules}\,R \simeq \mathrm{SheafOfModules}\,S\), compatibly with the
  underlying ring-sheaf identification. This is the module-level lift used in step~3 of
  \cref{lem:over_restrict_iso}.
\end{lemma}

\begin{lemma}\leanok
[The over-equivalence functor is cocontinuous]
  \label{lem:overEquivalence_functor_isCocontinuous}
  \lean{AlgebraicGeometry.overEquivalence_functor_isCocontinuous}
  \uses{lem:opens_overEquivalence_mathlib}
  \textit{Topological layer of bridge C; fills the \texttt{Topology/Sheaves/Over.lean} TODO.}
  The functor of \(\mathrm{Opens.overEquivalence}\,U\) (\cref{lem:opens_overEquivalence_mathlib})
  is cocontinuous (cover-lifting) from the \(U\)-slice \(J_X|_U\) of the opens Grothendieck
  topology of \(X\) to the opens Grothendieck topology \(J_U\) of the subspace \(U\).
\end{lemma}
\begin{proof}
  \uses{lem:opens_overEquivalence_mathlib}
  A sieve covers in the slice topology \(J_X|_U\) iff its image sieve covers in \(J_X\), and
  covering for the opens topology is pointwise: a sieve covers iff its members' opens cover the
  base set-theoretically. Given a point \(x\) of a slice object and a cover member in \(J_U\)
  containing the image of \(x\), transport it across the open embedding \(U \hookrightarrow X\),
  which is injective: the set-theoretic image of an open \(V \subseteq U\) is an open of \(X\)
  contained in the base, a neighbourhood of \(x\) lying in the image sieve. This witnesses the
  cover-lifting condition.
\end{proof}

\begin{lemma}\leanok
[The over-equivalence inverse is cocontinuous]
  \label{lem:overEquivalence_inverse_isCocontinuous}
  \lean{AlgebraicGeometry.overEquivalence_inverse_isCocontinuous}
  \uses{lem:opens_overEquivalence_mathlib}
  \textit{Topological layer of bridge C; fills the \texttt{Topology/Sheaves/Over.lean} TODO.}
  Symmetrically, the inverse of \(\mathrm{Opens.overEquivalence}\,U\) is cocontinuous from the
  opens topology \(J_U\) of the subspace \(U\) to the \(U\)-slice \(J_X|_U\).
\end{lemma}
\begin{proof}
  \uses{lem:opens_overEquivalence_mathlib}
  As in \cref{lem:overEquivalence_functor_isCocontinuous}, reduce to the pointwise covering
  criterion and the over-sieve image criterion; given a point and a slice-cover member, pull the
  member back along the embedding \(U \hookrightarrow X\) (preimage of an open of \(X\) is an open
  of the subspace) to obtain a cover member of \(J_U\) witnessing the lift.
\end{proof}

\begin{lemma}\leanok
[The over-equivalence inverse is a dense subsite]
  \label{lem:overEquivalence_inverse_isDenseSubsite}
  \lean{AlgebraicGeometry.overEquivalence_inverse_isDenseSubsite}
  \uses{lem:overEquivalence_functor_isCocontinuous, lem:overEquivalence_inverse_isCocontinuous}
  \textit{Topological layer of bridge C; the dense-subsite witness assembled from the two
  cocontinuities.}
  The inverse of \(\mathrm{Opens.overEquivalence}\,U\) is a dense subsite for \((J_U, J_X|_U)\).
  This is formal: an equivalence both of whose legs are cocontinuous
  (\cref{lem:overEquivalence_functor_isCocontinuous},
  \cref{lem:overEquivalence_inverse_isCocontinuous}) yields a dense subsite on its inverse.
\end{lemma}

\begin{lemma}\leanok
[The over-equivalence functor is continuous]
  \label{lem:overEquivalence_functor_isContinuous}
  \lean{AlgebraicGeometry.overEquivalence_functor_isContinuous}
  \uses{lem:overEquivalence_inverse_isCocontinuous}
  \textit{Topological layer of bridge C; needed for the module-level lift of step~3.}
  The functor of \(\mathrm{Opens.overEquivalence}\,U\) is continuous (cover-preserving) for
  \((J_X|_U, J_U)\). It is obtained from the cocontinuity of the inverse
  (\cref{lem:overEquivalence_inverse_isCocontinuous}) through the equivalence adjunction: the
  left adjoint of a cocontinuous functor is continuous.
\end{lemma}

\begin{lemma}\leanok
[The over-equivalence inverse is continuous]
  \label{lem:overEquivalence_inverse_isContinuous}
  \lean{AlgebraicGeometry.overEquivalence_inverse_isContinuous}
  \uses{lem:overEquivalence_functor_isCocontinuous}
  \textit{Topological layer of bridge C; needed for the module-level lift of step~3.}
  Symmetrically, the inverse of \(\mathrm{Opens.overEquivalence}\,U\) is continuous for
  \((J_U, J_X|_U)\), obtained from the cocontinuity of the functor
  (\cref{lem:overEquivalence_functor_isCocontinuous}) through the equivalence adjunction.
\end{lemma}

\begin{definition}\leanok
[The over-site/open-subspace sheaf equivalence]
  \label{def:overEquivalence_sheafCongr}
  \lean{AlgebraicGeometry.overEquivalence_sheafCongr}
  \uses{lem:opens_overEquivalence_mathlib, lem:equivalence_sheafCongr_mathlib,
    lem:overEquivalence_inverse_isDenseSubsite}
  \textit{Keystone of the topological layer of bridge C.}
  For any category \(A\), the site equivalence \(\mathrm{Opens.overEquivalence}\,U\)
  (\cref{lem:opens_overEquivalence_mathlib}) lifts to an equivalence of sheaf categories
  \[
    \mathrm{Sheaf}(J_X|_U,\, A) \;\simeq\; \mathrm{Sheaf}(J_U,\, A),
  \]
  obtained by feeding the dense-subsite witness
  (\cref{lem:overEquivalence_inverse_isDenseSubsite}) to the general sheaf-congruence
  construction (\cref{lem:equivalence_sheafCongr_mathlib}). This is exactly the equivalence left
  as a TODO in Mathlib's \texttt{Topology/Sheaves/Over.lean}, and it is the object across which
  \cref{lem:over_restrict_iso} transports the structure sheaf and the module \(M\).
\end{definition}

% --- (C, step 3) The module-level equivalence and its functor identification ---

% The two blocks below are the module-category layer assembled on top of the topological
% sheaf-congruence def:overEquivalence_sheafCongr: the step-3 equivalence of categories of sheaves of
% modules (def:over_restrict_equiv) and the step-4 functor identification (lem:over_restrict_functor_iso)
% whose object component is the slice-touching bridge lem:over_restrict_iso below. They are
% project-bespoke infrastructure (no external source).


% --- (C) The single slice-touching bridge ---

\begin{lemma}\leanok
[The Grothendieck-slice restriction equals the geometric restriction (gap1, C)]
  \label{lem:over_restrict_iso}
  \lean{AlgebraicGeometry.Scheme.Modules.overRestrictIso}
  \uses{lem:modules_restrictFunctor_mathlib, lem:modules_restrictFunctorIsoPullback_mathlib,
    lem:isLocalization_basicOpen_mathlib, def:overEquivalence_sheafCongr,
    lem:overEquivalence_functor_isContinuous, lem:overEquivalence_inverse_isContinuous,
    lem:opens_overEquivalence_mathlib, lem:pushforwardPushforwardEquivalence_mathlib}
  % NOTE: RESOLVED and axiom-clean. AlgebraicGeometry.Scheme.Modules.overRestrictIso
  % now exists (#print axioms = {propext, Classical.choice, Quot.sound}); gap1 bridge C
  % (steps 1->4) is fully closed. The step-2 geometric ring-sheaf identification collapsed to
  % `rfl` (U.toScheme.ringCatSheaf is definitionally the transport of the sliced sheaf, via
  % toScheme_presheaf_obj/map). The statement is routed through the step-3 equivalence functor:
  % overRestrictIso : (overRestrictEquiv U).functor.obj (M.over U) ≅ (restrictFunctor U.ι).obj M.
  \textit{The one lemma that must touch the Grothendieck slice site; everything downstream of it is
  geometric.}
  Let \(X\) be a scheme and \(U \subseteq X\) an open. The \emph{abstract Grothendieck-slice}
  restriction \(M.\mathrm{over}\,U\) — an object of the category of sheaves of modules over the
  sliced structure sheaf \(\mathcal{O}_X.\mathrm{over}\,U\) on the slice site \(J.\mathrm{over}\,U\) —
  is canonically isomorphic to the \emph{geometric} restriction
  \((\mathrm{restrictFunctor}\,U.\iota).\mathrm{obj}\,M\)
  (\cref{lem:modules_restrictFunctor_mathlib}), equivalently to the pullback \(U.\iota^{*} M\)
  (\cref{lem:modules_restrictFunctorIsoPullback_mathlib}), as a sheaf of modules on the small Zariski
  site of the open subscheme \(U.\mathrm{toScheme}\). The isomorphism is induced by the equivalence
  of sites \(J.\mathrm{over}\,U \simeq \mathrm{Opens}(U.\mathrm{toScheme})\) coming from the opens
  functor of the open immersion \(U.\iota\), under which the sliced sheaf of rings
  \(\mathcal{O}_X.\mathrm{over}\,U\) corresponds to the structure sheaf of \(U.\mathrm{toScheme}\).
  For \(X = \Spec R\) and \(U = D(r)\) a basic open, the target category is
  \((\Spec R_r).\mathrm{Modules}\) under the affine identification \(D(r) \cong \Spec R_r\)
  (\cref{lem:isLocalization_basicOpen_mathlib}).
\end{lemma}
\begin{proof}
  \uses{def:overEquivalence_sheafCongr, lem:overEquivalence_functor_isContinuous,
    lem:overEquivalence_inverse_isContinuous, lem:opens_overEquivalence_mathlib,
    lem:pushforwardPushforwardEquivalence_mathlib, lem:modules_restrictFunctor_mathlib,
    lem:modules_restrictFunctorIsoPullback_mathlib, lem:isLocalization_basicOpen_mathlib}
  The isomorphism is built in four steps, ascending from the purely topological layer to the
  module-level statement.

  \emph{Step 1 (topological layer --- done).} The opens functor of the open immersion
  \(U.\iota\) realizes the slice of the opens site of \(X\) over \(U\) as the opens site of the
  subspace \(U\); on sheaf categories this is the equivalence
  \(\mathrm{Sheaf}(J_X|_U,\,A) \simeq \mathrm{Sheaf}(J_U,\,A)\) of
  \cref{def:overEquivalence_sheafCongr}, for every target category \(A\). This is the
  topos-theoretic layer, assembled from the cocontinuity/continuity of both legs of
  \(\mathrm{Opens.overEquivalence}\,U\) (\cref{lem:opens_overEquivalence_mathlib}).

  \emph{Step 2 (geometric ring-sheaf identification --- the current obstacle).} Identify the
  sliced structure sheaf \(\mathcal{O}_X.\mathrm{over}\,U\) with the structure sheaf
  \(\mathcal{O}_{U}\) of the open subscheme \(U.\mathrm{toScheme}\), transported across the
  \(\mathrm{RingCat}\)-valued instance of \cref{def:overEquivalence_sheafCongr}; this is an
  isomorphism of sheaves of rings. The bridge is the factorization of the opens functor of the
  immersion,
  \[
    U.\iota.\mathrm{opensFunctor}
      \;=\; (\mathrm{Opens.overEquivalence}\,U).\mathrm{inverse}
        \;\circ\; \mathrm{Over.forget}\,U ,
  \]
  under which the structure-sheaf comparison reduces to the standard open-immersion
  structure-sheaf identifications; recall that \(\mathrm{restrictFunctor}\,U.\iota\)
  (\cref{lem:modules_restrictFunctor_mathlib}) is precisely pushforward along
  \(U.\iota.\mathrm{opensFunctor}\). This ring-sheaf identification is the remaining geometric
  obstacle.

  \emph{Step 3 (lift to modules).} Given the step-2 identification of the underlying ring sheaves
  and the continuity of both legs (\cref{lem:overEquivalence_functor_isContinuous},
  \cref{lem:overEquivalence_inverse_isContinuous}), lift the sheaf-of-rings comparison to an
  equivalence of categories of sheaves of modules via
  \(\mathrm{pushforwardPushforwardEquivalence}\)
  (\cref{lem:pushforwardPushforwardEquivalence_mathlib}). Under this equivalence the sliced
  module \(M.\mathrm{over}\,U\) corresponds to the geometric restriction
  \((\mathrm{restrictFunctor}\,U.\iota).\mathrm{obj}\,M\).

  \emph{Step 4 (compose).} Compose with
  \(\mathrm{restrictFunctorIsoPullback}\) (\cref{lem:modules_restrictFunctorIsoPullback_mathlib})
  to land the isomorphism \(M.\mathrm{over}\,U \cong M.\mathrm{restrict}\,U.\iota \cong
  U.\iota^{*}M\). The two sides live a priori in different module categories, so the literal Lean
  statement of \(\mathrm{overRestrictIso}\) is phrased through the step-3 equivalence functor; for
  \(X = \Spec R\) and \(U = D(r)\) the target is \((\Spec R_r).\mathrm{Modules}\) under
  \(D(r) \cong \Spec R_r\) (\cref{lem:isLocalization_basicOpen_mathlib}).
\end{proof}

% --- (C, step 4') The pullback form of the slice-to-geometric isomorphism ---

\begin{lemma}

\leanok
[The slice-to-geometric isomorphism in pullback form (gap1, C, step 4')]
  \label{lem:over_restrict_pullback_iso}
  \lean{AlgebraicGeometry.Scheme.Modules.overRestrictPullbackIso}
  \uses{lem:over_restrict_iso, lem:modules_restrictFunctorIsoPullback_mathlib,
    lem:modules_pullback_mathlib}
  \textit{The pullback packaging of the slice-touching bridge; the form consumed by the P1 transport.}
  With \(X\) and \(U \subseteq X\) as above and \(M\) a sheaf of modules on \(X\), the transport of the
  abstract Grothendieck-slice restriction \(M.\mathrm{over}\,U\) under the step-3 equivalence functor
  (\cref{def:over_restrict_equiv}) is canonically isomorphic to the inverse-image (pullback) of \(M\)
  along the open immersion \(U.\iota\):
  \[
    (\mathrm{overRestrictEquiv}\,U).\mathrm{functor}.\mathrm{obj}\,(M.\mathrm{over}\,U)
      \;\cong\; (U.\iota^{*})\,M
  \]
  (\cref{lem:modules_pullback_mathlib}). This is the pullback packaging of the slice-touching bridge
  \cref{lem:over_restrict_iso}; it is exactly the form in which a presentation of
  \(M.\mathrm{over}\,U\) is transported into a presentation of the geometric pullback \(U.\iota^{*}M\),
  and it is the ingredient consumed by the per-element local-tilde transport
  \cref{lem:isIso_fromTildeΓ_basicOpen_of_quasicoherent}.
\end{lemma}
\begin{proof}
  \uses{lem:over_restrict_iso, lem:modules_restrictFunctorIsoPullback_mathlib}
  Compose the slice-to-geometric isomorphism \cref{lem:over_restrict_iso}, which identifies the
  transported slice restriction with the geometric restriction
  \((\mathrm{restrictFunctor}\,U.\iota).\mathrm{obj}\,M\), with the component at \(M\) of the natural
  isomorphism \(\mathrm{restrictFunctorIsoPullback}\,U.\iota\)
  (\cref{lem:modules_restrictFunctorIsoPullback_mathlib}) identifying the geometric restriction with
  the pullback \(U.\iota^{*}M\).
\end{proof}

% --- (P1 prep) Slice→geometric Presentation-transport machinery ---
% These three project-bespoke helpers transport a presentation across the step-3
% slice equivalence into a presentation of the geometric pullback; they feed the
% P1 keystone (lem:isIso_fromTildeΓ_basicOpen_of_quasicoherent). No external source.


% --- (P1 prep) Affine-open presentation transport: slice → geometrize → affine-identify ---
% Six project-bespoke Scheme.Modules helpers (no external source — Mathlib presentation/pullback
% bookkeeping): they carry the global presentation of a cover member down to an arbitrary open W,
% then transport it across the affine identification W ≅ Spec Γ(Z,W). Together they assemble the
% general keystone lem:isIso_fromTildeΓ_presentationPullback, of which the basic-open P1 keystone
% lem:isIso_fromTildeΓ_basicOpen_of_quasicoherent is the affine instance W = (q.X i).ι ⁻¹ᵁ D(r).


% --- (P1) Per-element local-tilde transport ---


% --- (D) The section-localization descent: the gap1 keystone ---

% NOTE: the lemma STATEMENT is cited to Stacks lemma-invert-f-sections / Hartshorne II.5.3; the proof
% below is the project-internal finite-cover descent (gluing + separatedness + flatness of
% localization) that deliberately does NOT route through the global affine equivalence (which is gap1
% itself), so there is no transcribed `% SOURCE QUOTE PROOF`.

% --- (D, engines) The private sheaf-gluing sub-steps of the cover-form descent (coverage) ---


% --- (D, basic-open cover form) The form gap1 actually consumes ---


% --- (D, section transport) The Hfr ingredient: pullback-along-open-immersion section comparison ---


% NOTE: How lem:pullback_gamma_top_iso produces the gated Hfr (and thereby the named-form descent and
% gap1). For M : (Spec R).Modules quasi-coherent with datum q and a cover member q.X i ⊇ D(r), the P1
% transport (\cref{lem:isIso_fromTildeΓ_basicOpen_of_quasicoherent}) gives IsIso fromTildeΓ for the
% module M'_transported obtained by pulling M back along ι(q.X i), then along the basic-open ι, then
% across IsAffineOpen.isoSpec.inv. Chaining the section comparison \cref{lem:pullback_gamma_top_iso}
% through those same three pullbacks (and the pullback-form bridge
% \cref{lem:over_restrict_pullback_iso}) identifies Γ(M'_transported, ⊤) ≅ Γ(M, D(r)) and, at the
% f-locus open, Γ(M'_transported, D(f')) ≅ Γ(M, D(f) ⊓ D(r)), intertwining the restriction maps.
% Together with IsIso fromTildeΓ and \cref{lem:isLocalizedModule_restrict_of_isIso_fromTildeΓ} this
% yields IsLocalizedModule(powers f) on Γ(M, D(r)) → Γ(M, D(f) ⊓ D(r)) — exactly Hfr at U = D(r) (and
% at the overlaps U = D(r) ⊓ D(r')). Feeding Hfr to \cref{lem:section_localization_descent_of_cover}
% at the quasi-coherence cover \cref{lem:exists_finite_basicOpen_cover_le_quasicoherentData} closes
% the named form \cref{lem:section_localization_descent}, and gap1
% \cref{lem:qcoh_affine_isIso_fromTildeΓ} follows via
% \cref{lem:isIso_fromTildeΓ_of_isLocalizedModule_restrict}.


% NOTE: the remaining wall between the two bridges above and Hfr is the SEMILINEARITY of the
% section transport gammaPullbackImageIso over the open-immersion structure-sheaf ring iso. The next two
% blocks state that sub-build.


% NOTE: the lemma STATEMENT is cited to Stacks lemma-invert-f-sections / Hartshorne II.5.3; the proof
% below is a project-internal descent (finite sheaf equalizer over the cover + flatness of
% localization) that deliberately does NOT route through the global affine equivalence (which is
% gap1 itself), so there is no transcribed `% SOURCE QUOTE PROOF`.

% NOTE: the proof is a one-line assembly — feed the keystone descent (D) into the in-file
% equivalence isIso_fromTildeΓ_iff_isLocalizedModule_restrict — so it routes through no transcribed
% Stacks gluing proof; the SOURCE QUOTE above cites the lemma STATEMENT (Stacks
% lemma-quasi-coherent-affine).

\subsection{Section-transport producer for the basic-open Hfr}

This subsection builds the sole remaining ingredient of the gap1 keystone descent
\cref{lem:section_localization_descent}: the geometric \emph{producer}
\cref{lem:section_localization_hfr_basicOpen} that manufactures, for a quasi-coherent \(M\) on
\(\Spec R\) and a basic open \(D(s)\) inside a quasi-coherence cover member, the per-cover-element
localization datum \(\mathrm{Hfr}\) the basic-open cover form
\cref{lem:section_localization_descent_of_basicOpen_cover} consumes. The construction transports the
\(R_s\)-level localization living on the affine slice \(D(s) \cong \Spec R_s\) back to an
\(R\)-level statement through the combiner \cref{lem:isLocalizedModule_powers_transport}; the four
sub-lemmas below supply the geometric inputs (an iterated-pullback \(\mathrm{fromTildeΓ}\) iso, the
image/range computations, the \(\top\)-level semilinearity, and the scalar-tower instances).


\subsection{G1-assemble: from per-basic-open section localization to the
\texorpdfstring{$\widetilde{(-)}$}{tilde}-counit}

This subsection records the bridge \(\text{section-localization on every basic open} \Rightarrow
\mathrm{IsIso}\,M.\mathrm{fromTildeΓ}\): given that the section restriction localizes on every basic
open, the tilde-Gamma counit is an isomorphism. Together with its converse
\cref{lem:isLocalizedModule_restrict_of_isIso_fromTildeΓ} it assembles the equivalence
\cref{lem:isIso_fromTildeΓ_iff_isLocalizedModule_restrict}, which is exactly what makes G1-core
(\cref{lem:qcoh_affine_section_localization}) and gap1
(\cref{lem:qcoh_affine_isIso_fromTildeΓ}) interderivable -- the route by which G1-core is obtained as
a corollary of gap1. The supporting Mathlib facts are recorded as dependency anchors.

% --- Mathlib dependency anchors for the G1-assemble step ---


% --- Project-original G1-assemble nodes (Lean realizations) ---


\begin{definition}

\leanok
[Schematic support of a sheaf of modules]
  \label{def:schematic_support}
  \lean{AlgebraicGeometry.Scheme.Modules.schematicSupport}
  \uses{def:modules_annihilator, lem:modules_annihilator_ideal_le,
    def:schematic_support_immersion}
  % SOURCE: [Nitsure], \S 1, sub-section "Stratification by Hilbert
  % Polynomials" (read from
  % references/nitsure-hilbert-quot-src/nitsure-hilbert-quot.tex, L468-L471).
  % SOURCE QUOTE:
  % "Let $X\to S$ be a finite type morphism of noetherian schemes,
  %  and let $L$ be a line bundle on $X$. Let $\F$
  %  be any coherent sheaf on
  %  $X$ whose schematic support is proper over $S$."
  \textit{Source: [Nitsure], \S 1 (FGA Explained Ch.~5).}
  Let \(\mathcal{F}\) be a sheaf of modules on a scheme \(X\). The
  \emph{schematic support} of \(\mathcal{F}\) is the closed subscheme of \(X\)
  cut out by its annihilator ideal sheaf (\cref{def:modules_annihilator}):
  \[
    \mathrm{schematicSupport}(\mathcal{F})
    \;=\; V\bigl(\mathrm{Ann}(\mathcal{F})\bigr) \;\hookrightarrow\; X,
  \]
  equipped with its canonical closed immersion into \(X\). This realises
  Nitsure's ``schematic support of \(\mathcal{F}\)'' as a closed subscheme of
  \(X\).
\end{definition}

\begin{definition}\leanok
[Schematic-support immersion]
  \label{def:schematic_support_immersion}
  \lean{AlgebraicGeometry.Scheme.Modules.schematicSupportι}
  \uses{def:schematic_support, def:modules_annihilator}
  Let \(\mathcal{F}\) be a sheaf of modules on a scheme \(X\). The
  \emph{schematic-support immersion} is the canonical closed immersion
  \[
    \mathrm{schematicSupport}(\mathcal{F}) \;\hookrightarrow\; X
  \]
  of the schematic support (\cref{def:schematic_support}) into the ambient
  scheme, namely the subscheme immersion of the annihilator ideal sheaf
  \(\mathrm{Ann}(\mathcal{F})\) (\cref{def:modules_annihilator}). It is a
  quasi-compact preimmersion onto the support, and it is the morphism whose
  composite with a base morphism \(f : X \to S\) defines proper support
  (\cref{def:has_proper_support}). Proved directly in Lean.
\end{definition}

\begin{lemma}[Properness as a morphism property]
  \label{lem:isProper_mathlib}
  \lean{AlgebraicGeometry.IsProper}
  \mathlibok
  \textit{Provided by Mathlib
  (\texttt{Mathlib.AlgebraicGeometry.Morphisms.Proper}).}
  A morphism \(f : X \to S\) of schemes is \emph{proper} iff it is separated,
  universally closed, and locally of finite type. Properness is stable under base
  change: if \(f : X \to S\) is proper and \(g : T \to S\) is any morphism then
  the base change \(X \times_S T \to T\) is proper. Consequently the base-change
  clause that Nitsure invokes (``properness \ldots\ preserved by base-change'')
  applies directly to the proper-support condition below.
\end{lemma}

\begin{definition}

\leanok
[Proper support over the base]
  \label{def:has_proper_support}
  \lean{AlgebraicGeometry.Scheme.Modules.HasProperSupport}
  \uses{def:schematic_support, def:schematic_support_immersion, lem:isProper_mathlib}
  % SOURCE: [Nitsure], \S 1, sub-section "The Functors $\Hilb_{X/S}$ and
  % $\Quot_{E/X/S}$" (read from
  % references/nitsure-hilbert-quot-src/nitsure-hilbert-quot.tex, L407-L409
  % + L418-L419).
  % SOURCE QUOTE:
  % "(i) a coherent sheaf $\F$ on $X_T = X\times_ST$
  %  such that the schematic support of $\F$ is proper over $T$
  %  and $\F$ is flat over $T$, together with
  %  [...]
  %  Then as properness and flatness are preserved
  %  by base-change, and as tensor-product is right exact,"
  \textit{Source: [Nitsure], \S 1 (FGA Explained Ch.~5).}
  Let \(\mathcal{F}\) be a sheaf of modules on a scheme \(X\) and let
  \(f : X \to S\) be a morphism. We say \(\mathcal{F}\) \emph{has proper support
  over \(S\) (along \(f\))} when the composite morphism
  \[
    \mathrm{schematicSupport}(\mathcal{F}) \hookrightarrow X \xrightarrow{f} S
  \]
  is proper (\cref{lem:isProper_mathlib}). Since properness is stable under base
  change (\cref{lem:isProper_mathlib}), this condition is preserved by pull-back
  along any \(S'\to S\), which is exactly what the well-definedness of the Quot
  functor's pullback action requires.
\end{definition}

\begin{definition}

\leanok
[Locally free of rank \(d\)]
  \label{def:is_locally_free_of_rank}
  \lean{AlgebraicGeometry.SheafOfModules.IsLocallyFreeOfRank}
  % SOURCE: [Nitsure], \S 1, "Elementary Examples, Exercises (2)" and
  % "Grassmannian of a vector bundle" (read from
  % references/nitsure-hilbert-quot-src/nitsure-hilbert-quot.tex, L562-L564
  % + L584-L585).
  % SOURCE QUOTE:
  % "$\langle \F, q\rangle$ of quotients $q: \oplus^r\,{\mathcal O}_T \to \F$
  %  where $\F$ is a locally free sheaf on $T$ of rank $d$.
  %  [...]
  %  the set of all equivalence classes $\langle \F, q\rangle$ of quotients
  %  $q: E_T \to \F$ where $\F$ is a locally free sheaf on $T$ of rank $d$,"
  \textit{Source: [Nitsure], \S 1, Exercise (2) (FGA Explained Ch.~5).}
  Let \(\mathcal{M}\) be a sheaf of modules on a scheme \(X\) and let
  \(d \in \mathbb{N}\). We say \(\mathcal{M}\) is \emph{locally free of rank
  \(d\)} when there exists an open cover \(\{U_i\}\) of \(X\) together with
  isomorphisms \(\mathcal{M}|_{U_i} \cong \mathcal{O}_{U_i}^{\oplus d}\) for
  each \(i\).

  \emph{Note.} Mathlib does not provide this predicate.
\end{definition}

\section{The Quot functor}
\label{sec:quot_functor}

\begin{definition}\leanok
[Quot functor]
  \label{def:quot_functor}
  \lean{AlgebraicGeometry.Scheme.QuotFunctor}
  \uses{def:hilbert_polynomial, def:has_proper_support}

  % SOURCE: [Nitsure], \S 1, sub-section "The Functors \(\Hilb_{X/S}\) and
  % \(\Quot_{E/X/S}\)" (read from
  % references/nitsure-hilbert-quot-src/nitsure-hilbert-quot.tex, L394-L433
  % + L481-L498).
  % SOURCE QUOTE:
  % "Let \(S\) be a noetherian scheme and let \(X\to S\) be a finite type
  %  scheme over it. Let \(E\) be a coherent sheaf on \(X\).
  %  Let \(Sch_S\) denote the category of all locally noetherian schemes
  %  over \(S\). For any \(T\to S\) in \(Sch_S\), a \textbf{family of quotients of \(E\)
  %  parametrised by \(T\)} will mean
  %  a pair \((\F,q)\) consisting of
  %
  %  (i) a coherent sheaf \(\F\) on \(X_T = X\times_ST\)
  %  such that the schematic support of \(\F\) is proper over \(T\)
  %  and \(\F\) is flat over \(T\), together with
  %
  %  (ii) a surjective \({\mathcal O}_{X_T}\)-linear homomorphism of sheaves
  %  \(q:E_T  \to \F\) where \(E_T\) is the pull-back of \(E\) under
  %  the projection \(X_T \to X\).
  %
  %  Two such families \((\F,q)\) and \((\F,q)\) parametrised by \(T\) will
  %  be regarded as equivalent if \(\ker(q) = \ker(q')\)),
  %  and \(\langle \F, q \rangle\) will denote
  %  an equivalence class. Then as properness and flatness are preserved
  %  by base-change, and as tensor-product is right exact,
  %  the pull-back of \(\langle \F, q \rangle\) under an \(S\)-morphism \(T'\to T\) is
  %  well-defined, which gives
  %  a set-valued contravariant functor
  %  \(\Quot_{E/X/S} : Sch_S \to Sets\) under which
  %  \(\)T\mapsto  \{\mbox{ All } \langle \F, q \rangle \mbox{ parametrised by }T\}\(\)
  %  [...]
  %  This shows that the functor \(\Quot_{E/X/S}\) naturally decomposes
  %  as a co-product
  %  \(\)\Quot_{E/X/S} = \coprod_{{\Phi}\in \Q[\lambda]}\, \Quot_{E/X/S}^{\Phi,L}\(\)
  %  where for any polynomial \({\Phi}\in \Q[\lambda]\), the functor
  %  \(\Quot_{E/X/S}^{\Phi,L}\) associates
  %  to any \(T\) the set of all equivalence classes of
  %  families \(\langle \F,q \rangle\) such that at each \(t\in T\) the
  %  Hilbert polynomial of the restriction \(\F_t\), calculated
  %  using the pull-back of \(L\), is \({\Phi}\)."
  \textit{Source: [Nitsure], \S 1 (FGA Explained Ch.~5).}
  Let \(S\) be a noetherian scheme, \(\pi : X \to S\) a finite-type morphism,
  \(L\) a line bundle on \(X\), \(E\) a coherent sheaf on \(X\), and
  \(\Phi \in \mathbb{Q}[\lambda]\). The \emph{Quot functor with Hilbert
  polynomial \(\Phi\)}
  \[
    \Quot^{\Phi,L}_{E/X/S} : (\mathrm{Sch}/S)^{op} \longrightarrow \mathbf{Set}
  \]
  sends an \(S\)-scheme \(T \to S\) to the set of equivalence classes
  \(\langle \mathcal{F}, q\rangle\) of pairs \((\mathcal{F}, q)\) where
  \begin{enumerate}
    \item \(\mathcal{F}\) is a coherent sheaf on \(X_T = X \times_S T\) whose
      schematic support is proper over \(T\) (\cref{def:has_proper_support}) and
      which is flat over \(T\);
    \item \(q : E_T \twoheadrightarrow \mathcal{F}\) is a surjective
      \(\mathcal{O}_{X_T}\)-linear homomorphism, where \(E_T\) is the pull-back
      of \(E\) along \(X_T \to X\);
    \item for every \(t \in T\) the Hilbert polynomial of \(\mathcal{F}|_{X_t}\)
      with respect to \(L|_{X_t}\) (\cref{def:hilbert_polynomial}) equals
      \(\Phi\).
  \end{enumerate}
  Two pairs \((\mathcal{F}, q)\) and \((\mathcal{F}', q')\) are equivalent iff
  \(\ker(q) = \ker(q')\). On morphisms \(f : T' \to T\) over \(S\), the functor
  acts by pulling back the quotient along \(X_{T'} \to X_T\) (well-defined
  since tensor product is right-exact and preserves both flatness and
  properness). The Hilbert scheme is the special case
  \(\Hilb^{\Phi,L}_{X/S} := \Quot^{\Phi,L}_{\mathcal{O}_X/X/S}\).
\end{definition}

\section{The Grassmannian scheme}
\label{sec:grassmannian_scheme}

The Quot construction proceeds by embedding \(\Quot^{\Phi,L}_{E/X/S}\) as a
locally closed subfunctor of a Grassmannian over \(S\). Since Mathlib carries
no Grassmannian \emph{scheme} (only a finite-rank-quotient variety in
linear algebra), we package the construction here.

\begin{definition}\leanok
[Grassmannian scheme]
  \label{def:grassmannian_scheme}
  \lean{AlgebraicGeometry.Scheme.Grassmannian}
  \uses{def:quot_functor, def:is_locally_free_of_rank}
  % SOURCE: [Nitsure], \S 1, "Elementary Examples, Exercises~(2)" and
  % "Construction of Grassmannian" (read from
  % references/nitsure-hilbert-quot-src/nitsure-hilbert-quot.tex, L546-L599
  % + L773-L850).
  % SOURCE QUOTE:
  % "For any integers \(r\ge d\ge 1\), an explicit construction
  %  the Grassmannian scheme
  %  \(\grass(r,d)\) over \(\Z\), together with the tautological
  %  quotient \(u:\oplus^r\,{\mathcal O}_{\grass(r,d)} \to \U\) where \
  %  \(\U\) is a rank \(d\) locally free sheaf on \(\grass(r,d)\),
  %  has been given at the end of this section. [...]
  %  Show that \(\grass(r,d)\) together with the quotient
  %  \(u:\oplus^r\,{\mathcal O}_{\grass(r,d)} \to \U\) represents the
  %  contravariant functor
  %  \(\)\Grass(r,d) = \Quot_{\oplus^r\,{\mathcal O}_{\Z}/\Z/\Z}^{d, {\mathcal O}_{\Z}}\(\)
  %  from schemes to sets, which associates to any \(T\) the
  %  set of all equivalence classes
  %  \(\langle \F, q\rangle\) of quotients \(q: \oplus^r\,{\mathcal O}_T \to \F\)
  %  where \(\F\) is a locally free sheaf on \(T\) of rank \(d\).
  %  [...]
  %  Using the above show that, more generally, if \(S\) is a scheme and
  %  \(E\) is a locally free \({\mathcal O}_S\)-module of rank \(r\),
  %  the functor \(\Grass(E,d)=\Quot_{E/S/S}^{d,{\mathcal O}_S}\)
  %  on all \(S\)-schemes which
  %  by definition associates to any \(T\)
  %  the set of all equivalence classes \(\langle \F, q\rangle\) of quotients
  %  \(q: E_T \to \F\) where \(\F\) is a locally free sheaf on \(T\) of rank \(d\),
  %  is representable. The representing scheme
  %  is denoted by \(\grass(E,d)\) and is called the rank \(d\)
  %  relative Grassmannian of \(E\) over \(S\)."
  \textit{Source: [Nitsure], \S 1, Exercise (2) and "Construction of
  Grassmannian" (FGA Explained Ch.~5).}
  Let \(S\) be a noetherian scheme and let \(V\) be a locally free
  \(\mathcal{O}_S\)-module of rank \(r\). For an integer \(1 \le d \le r\), the
  \emph{Grassmannian of rank-\(d\) quotients of \(V\)} is the functor
  \[
    \mathrm{Grass}(V, d) : (\mathrm{Sch}/S)^{op} \longrightarrow \mathbf{Set},
    \qquad
    T \longmapsto \{ \langle \mathcal{F}, q\rangle :
    q : V_T \twoheadrightarrow \mathcal{F},\; \mathcal{F}
    \text{ locally free of rank } d \}
  \]
  with equivalence \(\ker(q) = \ker(q')\); the rank-\(d\) local freeness of
  \(\mathcal{F}\) is as in \cref{def:is_locally_free_of_rank}. Concretely,
  \(\mathrm{Grass}(V, d) = \Quot^{d,\mathcal{O}_S}_{V/S/S}\): the Quot functor
  (\cref{def:quot_functor}) in the special case \(X = S\), \(E = V\), constant
  Hilbert polynomial \(\Phi = d\).
\end{definition}

\begin{lemma}[Representability predicate]
  \label{lem:functor_is_representable_mathlib}
  \lean{CategoryTheory.Functor.IsRepresentable}
  \mathlibok
  \textit{Provided by Mathlib.}
  Let \(\mathcal{C}\) be a category and \(F : \mathcal{C}^{op} \to \mathbf{Set}\)
  a presheaf of sets. The predicate \(F.\mathrm{IsRepresentable}\) asserts that
  \(F\) is representable, i.e.\ that there exists an object \(Y\) of
  \(\mathcal{C}\) together with a natural bijection
  \(\mathrm{Hom}_{\mathcal{C}}(-, Y) \cong F\) (the Yoneda bijection).
\end{lemma}

\begin{theorem}\leanok
[Representability of the Grassmannian]
  \label{thm:grassmannian_representable}
  \lean{AlgebraicGeometry.Scheme.Grassmannian.representable}
  \uses{def:grassmannian_scheme, thm:relative_spec_exists, thm:relative_spec_univ,
    lem:functor_is_representable_mathlib, def:gr_glued_scheme, lem:gr_separated,
    lem:gr_proper}

  % NOTE (representability proof blocked): The proof sketch recovers the representing object
  % through the RepresentableBy Yoneda form via \cref{thm:relative_spec_univ}. However,
  % \cref{thm:relative_spec_univ} is currently established in Lean in a form weaker than a full
  % \(\mathrm{Functor.RepresentableBy}\) witness. The proof is blocked on either (a)
  % strengthening \cref{thm:relative_spec_univ} to deliver a RepresentableBy witness, or (b)
  % finding a RepresentableBy-free argument via direct gluing of chart-wise classifying morphisms.
  % This is a deferred open question.
  % NOTE: The Lean statement \lean{AlgebraicGeometry.Scheme.Grassmannian.representable} is
  % currently a weakened existence skeleton that omits smoothness, properness, relative dimension
  % \(d(r-d)\), the tautological rank-\(d\) quotient, and the Pl\"ucker embedding. The
  % \verb|\lean{}| pin above points at a declaration that under-delivers the prose statement; it
  % should be strengthened or split into a separate skeleton label.

  % SOURCE: [Nitsure], \S 1, sub-section "Construction by gluing together
  % affine patches" through "Projective embedding"
  % (read from references/nitsure-hilbert-quot-src/nitsure-hilbert-quot.tex,
  % L798-L930).
  % SOURCE QUOTE:
  % "For any integers \(r\ge d\ge 1\), the Grassmannian scheme
  %  \(\grass(r,d)\) over \(\Z\), together with the tautological
  %  quotient \(u:\oplus^r\,{\mathcal O}_{\grass(r,d)} \to \U\) where \
  %  \(\U\) is a rank \(d\) locally free sheaf on \(\grass(r,d)\),
  %  can be explicitly constructed as follows. [...]
  %  the schemes
  %  \(U^I\), as \(I\) varies over all the
  %  \({r\choose d}\) different
  %  subsets of  \(\{ 1,\ldots, r\}\) of cardinality \(d\),
  %  can be glued together by the co-cycle \((\theta_{I,J})\) to form a
  %  finite-type scheme \(\grass(r,d)\) over \(\Z\). As each \(U^I\) is
  %  isomorphic to \(\AA^{d(r-d)}_{\Z}\), it follows that
  %  \(\grass(r,d) \to \spec \Z\) is smooth of relative dimension \(d(r-d)\).
  %  [...]
  %  We now show that \(\pi : \grass(r,d) \to \spec \Z\) is proper. It is enough to
  %  verify the valuative criterion of properness for discrete valuation rings.
  %  [...]
  %  As \(\U\) is given by the transition functions \(g_{I,J}\) described above,
  %  the determinant line bundle
  %  \(\det(\U)\) is given by the transition functions
  %  \(\det(g_{I,J}) = 1/P^I_J \in GL_1(U^I_J)\).
  %  [...]
  %  We leave it to the reader to verify that the sections \(\sigma_I\) form
  %  a linear system which is base point free and separates points
  %  relative to \(\spec \Z\), and so gives an embedding of
  %  \(\grass(r,d)\) into \(\PP^m_{\Z}\) where \(m = {r\choose d}-1\).
  %  This is a closed embedding by the properness of
  %  \(\pi:\grass(r,d) \to \spec \Z\). In particular, \(\det(\U)\)
  %  is a relatively very ample line bundle on \(\grass(r,d)\) over \(\Z\)."
  \textit{Source: [Nitsure], \S 1, "Construction of Grassmannian" (FGA
  Explained Ch.~5).}
  Let \(S\) be a noetherian scheme, \(V\) a locally free \(\mathcal{O}_S\)-module
  of rank \(r\), and \(1 \le d \le r\). The functor \(\mathrm{Grass}(V,d)\) of
  \cref{def:grassmannian_scheme} is representable by a smooth projective
  \(S\)-scheme \(\mathrm{Gr}_S(V, d) \to S\) of relative dimension \(d(r-d)\),
  equipped with a tautological quotient
  \(\pi^* V \twoheadrightarrow \mathcal{U}\) of rank \(d\). The determinant
  line bundle \(\det(\mathcal{U})\) is relatively very ample, and the
  associated linear system gives a closed embedding
  \(\mathrm{Gr}_S(V, d) \hookrightarrow
  \mathbb{P}_S\!\left(\bigwedge^d V\right)\) -- the \emph{Pl\"ucker embedding}.

\end{theorem}

% SOURCE QUOTE PROOF: see the "Construction of Grassmannian" sub-section of
% [Nitsure], \S 1 (references/nitsure-hilbert-quot-src/nitsure-hilbert-quot.tex,
% L773-L930), which gives the construction in five named steps
% ("Construction by gluing together affine patches", "Separatedness",
% "Properness", "Universal quotient", "Projective embedding"). The total
% verbatim text is ~150 lines; it is reproduced piecewise in the
% SOURCE QUOTE comment of the theorem statement above.
\begin{proof}
  
  Over \(S = \Spec \mathbb{Z}\) and \(V = \mathcal{O}_S^{\oplus r}\) the
  Grassmannian \(\mathrm{Gr}(r, d)\) is constructed by gluing
  \(\binom{r}{d}\) affine charts \(U^I = \Spec\mathbb{Z}[X^I]\) indexed by
  size-\(d\) subsets \(I \subseteq \{1, \dots, r\}\). The chart \(U^I\)
  carries the matrix \(X^I\) of size \(d \times r\) whose \(I\)-th \(d \times d\)
  minor is the identity, with the remaining \(d(r-d)\) entries free
  variables (so \(U^I \cong \mathbb{A}^{d(r-d)}_{\mathbb{Z}}\)). The
  transition maps \(\theta_{I,J} : U^I \cap U^J \to U^J \cap U^I\) defined
  by \(X^J \mapsto (X^I_J)^{-1} X^I\) on the locus where \(\det X^I_J\) is
  invertible satisfy the cocycle condition
  \(\theta_{I,K} = \theta_{I,J} \circ \theta_{J,K}\). The glued
  \(\mathrm{Gr}(r,d) \to \Spec\mathbb{Z}\) is smooth of relative dimension
  \(d(r-d)\). Separatedness is checked from the diagonal
  \(\Delta_{I,J} \subset U^I \times U^J\) cut out by \(X^J_I X^I - X^J = 0\).
  Properness follows from the valuative criterion for DVRs: given
  \(\varphi : \Spec K \to \mathrm{Gr}(r,d)\) over a DVR \(R\) with fraction
  field \(K\), factoring through some chart \(U^I\) via a ring map
  \(f : \mathbb{Z}[X^I] \to K\), choose \(J\) minimising the valuation
  \(\nu(f(P^I_J))\) of the minor; on the chart \(U^J\) all entries have
  non-negative valuation, so \(f\) extends uniquely to \(\Spec R\).

  The universal quotient \(u : \mathcal{O}_{\mathrm{Gr}(r,d)}^{\oplus r}
  \twoheadrightarrow \mathcal{U}\) is given on \(U^I\) by the matrix \(X^I\)
  and glued by \(g_{I,J} = (X^I_J)^{-1} \in GL_d(U^I \cap U^J)\). The
  determinant line bundle \(\det\mathcal{U}\) has transition functions
  \(\det g_{I,J} = 1/P^I_J\); the global sections
  \(\sigma_I \in \Gamma(\mathrm{Gr}(r,d), \det\mathcal{U})\) defined by
  \(\sigma_I|_{U^J} = P^J_I\) form a base-point-free linear system separating
  points over \(\Spec\mathbb{Z}\), giving an embedding into
  \(\mathbb{P}^{\binom{r}{d}-1}_{\mathbb{Z}}\), which is closed by
  properness. Representability of \(\mathrm{Grass}(r,d)\) is then verified
  on the universal pair \((\mathrm{Gr}(r,d), u)\): a quotient
  \(q : \mathcal{O}_T^{\oplus r} \twoheadrightarrow \mathcal{F}\) of
  rank \(d\) on \(T\) is everywhere locally trivialised; on the locus where
  the columns indexed by \(I\) are a basis the matrix of \(q\) has the form
  \(X^I\), defining the classifying morphism \(T \to U^I \subset
  \mathrm{Gr}(r,d)\). Base change from \(\mathbb{Z}\) to a general noetherian
  base \(S\) with a vector bundle \(V\) is direct: trivialise \(V\) on an open
  cover and glue. For non-trivial \(V\) the same construction produces
  \(\mathrm{Gr}_S(V, d)\) with Pl\"ucker embedding into
  \(\mathbb{P}_S(\bigwedge^d V)\).
\end{proof}


%% ===========================================================================
%% UNION MERGE (2026-06-22): base-change cohomology lane preserved from this
%% project's own Picard_QuotScheme.tex (canonicalBaseChangeMap, flatBaseChange-
%% Cohomology, pullback_app_isoTensor and their support lemmas). Kept alongside
%% the subproject's graded/descent development above.
%% ===========================================================================
\section{Representability of the Quot scheme}
\label{sec:quot_representable}

\begin{theorem}\leanok
[Representability of the Quot functor]
  \label{thm:quot_representable}
  \lean{AlgebraicGeometry.Scheme.QuotScheme}
  \uses{def:quot_functor, lem:quot_boundedness, lem:quot_alpha_injective,
    lem:quot_reduction_to_pi_star_W, lem:quot_valuative_criterion,
    thm:grassmannian_representable, thm:flattening_stratification_exists,
    thm:flat_base_change_cohomology}

  % SOURCE: [Nitsure], \S 5, "Main Existence Theorems"
  % (read from references/nitsure-hilbert-quot-src/nitsure-hilbert-quot.tex,
  % L2220-L2273); cf.\ Grothendieck, FGA TDTE-IV.
  % SOURCE QUOTE:
  % "Let \(S\) be a noetherian scheme, \(\pi: X\to S\) a projective morphism,
  %  and \(L\) a relatively very ample line bundle on \(X\). Then for any
  %  coherent \({\mathcal O}_X\)-module \(E\) and any polynomial \(\Phi \in \Q[\lambda]\),
  %  the functor \(\Quot_{E/X/S}^{\Phi,L}\) is representable
  %  by a projective \(S\)-scheme \(\quot_{E/X/S}^{\Phi,L}\)."
  \textit{Source: [Nitsure], \S 5, Theorem (Grothendieck), Theorem
  (Altman--Kleiman) (FGA Explained Ch.~5); cf.\ Grothendieck, FGA TDTE-IV.}
  Let \(S\) be a noetherian scheme, \(\pi : X \to S\) a projective morphism,
  \(L\) a relatively very ample line bundle on \(X\), \(E\) a coherent
  \(\mathcal{O}_X\)-module, and \(\Phi \in \mathbb{Q}[\lambda]\). Then the
  functor \(\Quot^{\Phi,L}_{E/X/S}\) of \cref{def:quot_functor} is
  representable by a projective \(S\)-scheme \(\Quot^{\Phi,L}_{E/X/S}\),
  equipped with a universal quotient
  \[
    q^{\mathrm{univ}} :
    \pi^*_{\Quot} E \twoheadrightarrow \mathcal{F}^{\mathrm{univ}}
    \quad \text{on } X \times_S \Quot^{\Phi,L}_{E/X/S},
  \]
  flat over \(\Quot^{\Phi,L}_{E/X/S}\) with constant Hilbert polynomial
  \(\Phi\) on every fiber. The Hilbert scheme is the special case
  \(E = \mathcal{O}_X\): \(\Hilb^{\Phi,L}_{X/S} = \Quot^{\Phi,L}_{\mathcal{O}_X/X/S}\)
  represents the functor of \(T\)-flat closed subschemes
  \(Y \subset X_T\) with fiberwise Hilbert polynomial \(\Phi\).
\end{theorem}

% SOURCE QUOTE PROOF: this proof is long (~5 named sub-sections in [Nitsure]
% \S 5). Each sub-step is recorded with its own SOURCE QUOTE in the
% sub-lemmas below. The structural skeleton of the proof is restated in
% project notation in the \begin{proof} block.
\begin{proof}
  
  The proof has four steps; each step is a sub-lemma in the next
  sub-sections.
  \begin{enumerate}
    \item \textbf{(Boundedness, \cref{lem:quot_boundedness}.)} Pick the
      Castelnuovo--Mumford regularity bound \(m = m(\mathrm{rank}\,V,
      \mathrm{rank}\,W, \Phi)\) uniformly across all fibers (by Mumford's
      \(m\)-regularity theorem applied to \(E_s\) and to all coherent
      quotients of \(E_s\) with Hilbert polynomial \(\Phi\)). For \(r \ge m\)
      and any \(T \to S\) and any \(T\)-flat quotient
      \(q : E_T \twoheadrightarrow \mathcal{F}\) with Hilbert polynomial
      \(\Phi\), the higher direct images \(R^i (\pi_T)_* \mathcal{F}(r)\)
      and \(R^i (\pi_T)_* E_T(r)\) and \(R^i (\pi_T)_* \mathcal{G}(r)\)
      vanish for \(i \ge 1\) (where \(\mathcal{G} = \ker q\)), and the
      direct images \((\pi_T)_* \mathcal{F}(r)\), \((\pi_T)_* E_T(r)\),
      \((\pi_T)_* \mathcal{G}(r)\) are locally free of ranks fixed by the
      data; the canonical maps \((\pi_T)^* (\pi_T)_*(-)(r) \to (-)(r)\)
      are surjective. Cohomology-and-base-change (\cref{thm:flat_base_change_cohomology},
      Stacks tag 02KH) propagates these properties across \(T\)-base change.
    \item \textbf{(Grassmannian embedding, \cref{lem:quot_alpha_injective}.)}
      Reduce by \cref{lem:quot_reduction_to_pi_star_W} to the case
      \(X = \mathbb{P}_S(V)\), \(E = \pi^* W\) with \(V, W\) locally free
      on \(S\), \(L = \mathcal{O}_{\mathbb{P}(V)}(1)\). Fix \(r \ge m\).
      The Grassmannian embedding
      \[
        \alpha : \Quot^{\Phi,L}_{E/X/S} \longrightarrow
        \mathrm{Grass}\!\left(
          W \otimes_{\mathcal{O}_S} \mathrm{Sym}^r V,\; \Phi(r)
        \right)
      \]
      sends \(\langle \mathcal{F}, q\rangle\) on \(T\) to the locally free
      quotient \((\pi_T)_*(q(r)) : (\pi_T)_* E_T(r) \to (\pi_T)_* \mathcal{F}(r)\).
      This is injective on \(T\)-valued points: from the map \(T \to
      \mathrm{Gr}_S(W \otimes \mathrm{Sym}^r V, \Phi(r))\) one recovers
      the kernel-pull-back \((\pi_T)^* (\pi_T)_* \mathcal{G}(r) \to E_T(r)\),
      whose cokernel is \(q(r) : E_T(r) \to \mathcal{F}(r)\);
      twisting back by \(-r\) recovers \(q\).
    \item \textbf{(Locally closed in the Grassmannian via flattening
      stratification.)} Let \(T_0 = \mathrm{Gr}_S(W \otimes \mathrm{Sym}^r V,
      \Phi(r))\) with universal quotient
      \(f : (W \otimes \mathrm{Sym}^r V)_{T_0} \twoheadrightarrow \mathcal{J}\)
      and kernel \(\mathcal{K}\). Form the composite
      \(h : \pi_{T_0}^* \mathcal{K} \hookrightarrow
      \pi_{T_0}^* (W \otimes \mathrm{Sym}^r V) = \pi_{T_0}^* (\pi_{T_0})_* E_{T_0}
      \twoheadrightarrow E_{T_0}\)
      and let \(q := \mathrm{coker}(h) : E_{T_0} \to \mathcal{F}\). By the
      flattening stratification (\cref{thm:flattening_stratification_exists}
      of \cref{chap:Picard_FlatteningStratification}) applied to \(\mathcal{F}\)
      on \(X_{T_0}\), there exists a locally closed subscheme
      \(T_0^\Phi \subseteq T_0\) such that for any \(\phi : Y \to T_0\), the
      pulled-back cokernel \(\mathcal{F}_Y\) is \(Y\)-flat with fiberwise
      Hilbert polynomial \(\Phi\) iff \(\phi\) factors through \(T_0^\Phi\).
      The pair \((T_0^\Phi, \text{universal cokernel})\) represents
      \(\Quot^{\Phi,L}_{E/X/S}\).
    \item \textbf{(Closed, hence projective,
      \cref{lem:quot_valuative_criterion}.)} The functor
      \(\Quot^{\Phi,L}_{E/X/S}\) satisfies the valuative criterion of
      properness with respect to DVRs (Exercise (7) of \S 1, proved
      using that flatness over a DVR is torsion-freeness). Combined
      with the locally closed embedding of step 3 into the projective
      \(\mathrm{Gr}_S \hookrightarrow \mathbb{P}_S(\bigwedge^{\Phi(r)}
      (W \otimes \mathrm{Sym}^r V))\), this upgrades to a \emph{closed}
      embedding. Hence \(\Quot^{\Phi,L}_{E/X/S}\) is a projective \(S\)-scheme.
  \end{enumerate}
  Finally, by \cref{lem:quot_reduction_to_pi_star_W}, the general case
  (arbitrary \(X\) projective over \(S\) and arbitrary coherent \(E\)) follows
  from the case \(X = \mathbb{P}(V)\), \(E = \pi^* W\): pick a closed
  embedding \(X \subset \mathbb{P}(V)\) and write \(E\) as a coherent quotient
  of \(\pi^*(W)(\nu)\) for some vector bundle \(W\) on \(S\) and integer \(\nu\);
  the Quot scheme of the original data is then a closed subscheme of the
  Quot scheme of the universal data, by the closed-embedding clause of
  \cref{lem:quot_reduction_to_pi_star_W}.
\end{proof}

\section{Sub-lemmas of the construction}
\label{sec:quot_sub_lemmas}

\begin{lemma}[Reduction to the universal case]
  \label{lem:quot_reduction_to_pi_star_W}
  \lean{AlgebraicGeometry.TODO.quotReductionToPiStarW}
  \uses{def:quot_functor}

  % SOURCE: [Nitsure], \S 5, sub-section "Reduction to the case of
  % \(\Quot_{\pi^*W/\PP(V)/S}^{\Phi,L}\)"
  % (read from references/nitsure-hilbert-quot-src/nitsure-hilbert-quot.tex,
  % L2294-L2332).
  % SOURCE QUOTE:
  % "It is enough to prove Theorem REF in the special case
  %  that \(X = \PP(V)\) and
  %  \(E= \pi^*(W)\) where \(V\) and \(W\) are vector bundles on \(S\), as
  %  a consequence of the next lemma.
  %
  %  \begin{lemma}
  %  (i) Let \(\nu\) be any integer.
  %  Then tensoring by \(L^{\nu}\) gives an isomorphism of
  %  functors from \(\Quot_{E/X/S}^{\Phi,L}\) to
  %  \(\Quot_{E(\nu)/X/S}^{\Psi,L}\) where the polynomial
  %  \(\Psi \in \Q[\lambda]\) is defined by \(\Psi(\lambda) = \Phi(\lambda + \nu)\).
  %
  %  (ii) Let \(\phi : E \to G\) be a surjective homomorphism of
  %  coherent sheaves on \(X\). Then the corresponding natural transformation
  %  \(\Quot_{G/X/S}^{\Phi,L} \to \Quot_{E/X/S}^{\Phi,L}\)
  %  is a closed embedding.
  %  \end{lemma}"
  \textit{Source: [Nitsure], \S 5 (FGA Explained Ch.~5).}
  Let \(S\), \(X\), \(L\), \(E\), \(\Phi\) be as in \cref{def:quot_functor}.
  \begin{enumerate}
    \item (Twist-shift.) For \(\nu \in \mathbb{Z}\), tensoring by
      \(L^{\otimes \nu}\) gives an isomorphism of functors
      \(\Quot^{\Phi,L}_{E/X/S} \xrightarrow{\sim}
      \Quot^{\Psi,L}_{E(\nu)/X/S}\) with \(\Psi(\lambda) = \Phi(\lambda + \nu)\).
    \item (Surjection induces closed embedding.) Given a surjection
      \(\phi : E \twoheadrightarrow G\) of coherent sheaves on \(X\), the
      induced natural transformation
      \(\Quot^{\Phi,L}_{G/X/S} \to \Quot^{\Phi,L}_{E/X/S}\) is a closed
      embedding of functors.
  \end{enumerate}
  Consequently, if \(\Quot^{\Phi,L}_{\pi^* W / \mathbb{P}(V)/S}\) is
  representable for all \(V\), \(W\) vector bundles on \(S\), then so is
  \(\Quot^{\Phi,L}_{E/X/S}\) for every \(X \subset \mathbb{P}(V)\) closed and
  every coherent \(E\) on \(X\) that is a quotient of \(\pi^*(W)(\nu)|_X\) for
  some \(W\) and \(\nu\).
\end{lemma}

% SOURCE QUOTE PROOF: "The statement (i) is obvious. The statement (ii) just
% says that given any locally noetherian scheme \(T\) and a family
% \(\langle \F,q\rangle \in \Quot_{E/X/S}^{\Phi,L}(T)\), there exists
% a closed subscheme \(T' \subset T\) with the following universal property:
% for any locally noetherian scheme \(U\) and a morphism \(f: U\to T\),
% the pulled back homomorphism of \({\mathcal O}_{X_U}\)-modules \(q_U : E_U \to \F_U\)
% factors via the pulled back homomorphism \(\phi_U : E_U \to G_U\)
% if and only if \(U \to T\) factors via \(T'\hra T\).
% This is satisfied by taking \(T'\) to be the vanishing scheme
% for the composite homomorphism \(\ker(\phi) \hra  E \stackrel{q}{\to} \F\)
% of coherent sheaves on \(X_T\) (see Remark REF),
% which makes sense here as both \(\ker(\phi)\) and \(\F\) are coherent
% on \(X_T\) and \(\F\) is flat over \(T\)."
\begin{proof}
  
  (i) is direct from the definitions: tensoring with \(L^{\otimes \nu}\) is
  an autoequivalence of coherent sheaves on \(X\) (and on every \(X_T\)),
  preserving surjectivity, \(T\)-flatness, and properness of support, and
  shifting the Hilbert polynomial by \(\nu\) in the variable.

  (ii) Given a family \(\langle \mathcal{F}, q\rangle\) on \(T\) in
  \(\Quot^{\Phi,L}_{E/X/S}\), the locus where \(q\) factors through
  \(\phi : E \twoheadrightarrow G\) is cut out by the vanishing of the
  composite \(\ker(\phi)|_{X_T} \hookrightarrow E_T \xrightarrow{q}
  \mathcal{F}\). Because \(\mathcal{F}\) is \(T\)-flat, the schematic image of
  this composite descends to a closed subscheme \(T' \subset T\) such that
  \(f : U \to T\) factors through \(T'\) iff \(q_U\) factors through \(\phi_U\).
  Hence \(\Quot^{\Phi,L}_{G/X/S} \hookrightarrow \Quot^{\Phi,L}_{E/X/S}\)
  is a closed embedding of subfunctors.
\end{proof}

\begin{lemma}[Uniform Castelnuovo--Mumford regularity bound]
  \label{lem:quot_boundedness}
  \lean{AlgebraicGeometry.TODO.quotBoundedness}
  \uses{thm:flat_base_change_cohomology}

  % SOURCE: [Nitsure], \S 5, sub-section "Use of \(m\)-Regularity"
  % (read from references/nitsure-hilbert-quot-src/nitsure-hilbert-quot.tex,
  % L2346-L2404). Builds on [Nitsure] \S 2 (Castelnuovo--Mumford regularity)
  % and \S 3 (semi-continuity and base-change), cited internally as
  % Theorem~"Mumford" and Theorem~"Base-change for flat sheaves".
  % SOURCE QUOTE:
  % "We consider the sheaf \(E = \pi^*(W)\) on \(X=\PP(V)\) where
  %  \(V\) is a vector bundle on \(S\),
  %  and take \(L = {\mathcal O}_{\PP(V)}(1)\).
  %  For any field \(k\) and a \(k\)-valued point
  %  \(s\) of \(S\), we have an isomorphism
  %  \(\PP(V)_s \simeq \PP^n_k\) where \(n=\rank(V)-1\),
  %  and the restricted sheaf \(E_s\) on \(\PP(V)_s\) is isomorphic to
  %  \(\oplus^p{\mathcal O}_{\PP(V)_s}\) where \(p=\rank(W)\).
  %  It follows from Theorem REF
  %  that given any \(\Phi \in \Q[\lambda]\), there exists an integer
  %  \(m\) which depends only on \(\rank(V)\), \(\rank(W)\) and \(\Phi\),
  %  such that for any field \(k\) and a \(k\)-valued point \(s\) of \(S\),
  %  the sheaf \(E_s\) on \(\PP(V)_s\) is \(m\)-regular, and for any
  %  coherent quotient \(q: E_s \to \F\) on \(\PP(V)_s\) with
  %  Hilbert polynomial \(\Phi\), the sheaf \(\F\) and the kernel sheaf
  %  \(\G\subset E_s\) of \(q\) are both \(m\)-regular."
  \textit{Source: [Nitsure], \S 5, "Use of \(m\)-Regularity" (FGA Explained
  Ch.~5).}
  Let \(X = \mathbb{P}_S(V)\), \(E = \pi^* W\) with \(V\), \(W\) vector bundles on
  \(S\) of ranks \(n+1\), \(p\), \(L = \mathcal{O}_{\mathbb{P}(V)}(1)\), and
  \(\Phi \in \mathbb{Q}[\lambda]\). There exists an integer
  \(m = m(n, p, \Phi)\), depending only on \(n\), \(p\), and \(\Phi\), such that
  for every geometric point \(s\) of \(S\) and every coherent quotient
  \(q : E_s \twoheadrightarrow \mathcal{F}\) with Hilbert polynomial
  \(\Phi\), the sheaves \(E_s\), \(\mathcal{F}\), and \(\ker q\) on \(X_s\) are all
  \(m\)-regular in the sense of Castelnuovo--Mumford. Consequently, for
  every \(T \to S\) and every \(T\)-flat quotient
  \(q : E_T \twoheadrightarrow \mathcal{F}\) on \(X_T\) with Hilbert
  polynomial \(\Phi\) and every \(r \ge m\), the conclusions \textbf{(*)} and
  \textbf{(**)} of \cite[\S 5, "Use of \(m\)-Regularity"]{nitsure-hilbert-quot}
  hold: the direct images \((\pi_T)_* \mathcal{F}(r)\), \((\pi_T)_* E_T(r)\),
  \((\pi_T)_* \mathcal{G}(r)\) are locally free of fixed ranks; the
  canonical surjections \((\pi_T)^* (\pi_T)_*(-)(r) \twoheadrightarrow (-)(r)\)
  hold; and the higher direct images \(R^i (\pi_T)_* (-)(r)\) vanish for
  \(i \ge 1\).
\end{lemma}

% SOURCE QUOTE PROOF: see [Nitsure], \S 5, "Use of \(m\)-Regularity"
% (references/nitsure-hilbert-quot-src/nitsure-hilbert-quot.tex, L2346-L2404).
% The proof invokes the m-regularity theorem of [Nitsure] \S 2
% (Theorem "Mumford") to extract the uniform bound m, and the
% base-change-for-flat-sheaves theorem of [Nitsure] \S 3 (cited as
% "Base-change for flat sheaves") to propagate the conclusions across \(T\).
\begin{proof}
  
  Each geometric fiber \(E_s\) is a trivial bundle on \(\mathbb{P}^n_{\kappa(s)}\)
  of rank \(p\), so \(E_s\) is \(0\)-regular. The Castelnuovo--Mumford
  \(m\)-regularity theorem of [Nitsure] \S 2 produces, for any fixed Hilbert
  polynomial \(\Phi\), a uniform integer \(m = m(n, p, \Phi)\) above which
  every coherent quotient of \(E_s\) with Hilbert polynomial \(\Phi\) -- and
  its kernel -- is \(m\)-regular. The Castelnuovo Lemma of [Nitsure] \S 2
  converts \(m\)-regularity into vanishing of \(H^i((-) (r))\) for \(i \ge 1\)
  and \(r \ge m\), and into global generation of \(H^0((-)(r))\). To
  propagate the fiber-wise statement to families over \(T\), apply the
  flat-base-change theorem (\cref{thm:flat_base_change_cohomology}, Stacks
  tag 02KH) along \(T \to S\): cohomology of flat families commutes with
  base change, so the direct images \((\pi_T)_*(-)(r)\) are locally free of
  the rank dictated by the fiber Hilbert polynomials, and the canonical
  maps \((\pi_T)^* (\pi_T)_*(-)(r) \to (-)(r)\) inherit fiberwise
  surjectivity.
\end{proof}

\begin{lemma}[Grassmannian embedding \(\alpha\) is injective]
  \label{lem:quot_alpha_injective}
  \lean{AlgebraicGeometry.TODO.quotAlphaInjective}
  \uses{lem:quot_boundedness}

  % SOURCE: [Nitsure], \S 5, sub-section "Embedding Quot into Grassmannian"
  % (read from references/nitsure-hilbert-quot-src/nitsure-hilbert-quot.tex,
  % L2409-L2451).
  % SOURCE QUOTE:
  % "We now fix a positive integer \(r\) such that \(r\ge m\).
  %  Note that the rank of \({\pi_T}_*\F(r)\) is \(\Phi(r)\)
  %  and \(\pi_*E(r) = W \otimes_{{\mathcal O}_S}\sym^r V\).
  %  Therefore the surjective homomorphism
  %  \({\pi_T}_*E_T(r) \to {\pi_T}_*\F(r)\)
  %  defines an element of the set
  %  \(\Grass( W \otimes_{{\mathcal O}_S}\sym^r V , \Phi(r))(T)\).
  %  We thus get a morphism of functors
  %  \(\)\alpha : \Quot^{\Phi,L}_{E/X/S} \to
  %  \Grass( W \otimes_{{\mathcal O}_S}\sym^r V, \Phi(r))\(\)
  %  It associates to \(q: E_T \to \F\) the quotient
  %  \({\pi_T}_*(q(r)) : {\pi_T}_*E_T(r) \to {\pi_T}_*\F(r)\).
  %
  %  The above morphism \(\alpha\) is injective because the quotient
  %  \(q: E_T \to \F\) can be recovered from
  %  \({\pi_T}_*(q(r)) : {\pi_T}_*E_T(r) \to {\pi_T}_*\F(r)\)
  %  as follows."
  \textit{Source: [Nitsure], \S 5, "Embedding Quot into Grassmannian"
  (FGA Explained Ch.~5).}
  In the universal setup \(X = \mathbb{P}_S(V)\), \(E = \pi^* W\), fix
  \(r \ge m\) from \cref{lem:quot_boundedness}. The natural transformation
  \[
    \alpha :
    \Quot^{\Phi,L}_{E/X/S} \longrightarrow
    \mathrm{Grass}\!\left(
      W \otimes_{\mathcal{O}_S} \mathrm{Sym}^r V,\; \Phi(r)
    \right),
    \qquad
    \langle \mathcal{F}, q\rangle \longmapsto
    \langle (\pi_T)_*\mathcal{F}(r),\, (\pi_T)_*(q(r))\rangle
  \]
  is well-defined (the target is locally free of rank \(\Phi(r)\), and the
  canonical isomorphism \(\pi_* E(r) \cong W \otimes_{\mathcal{O}_S}
  \mathrm{Sym}^r V\) gives the source) and \emph{injective} on \(T\)-points:
  the data of the locally free quotient at level \(r\) determines
  \(q : E_T \to \mathcal{F}\).
\end{lemma}

\begin{proof}
  
  Suppose \(\alpha(T)\) sends both \(\langle \mathcal{F}_1, q_1\rangle\) and
  \(\langle \mathcal{F}_2, q_2\rangle\) to the same \(T\)-point of the
  Grassmannian. Then the pulled-back kernel
  \((\pi_T)^* (\pi_T)_* \mathcal{G}_i(r) \to (\pi_T)^* (\pi_T)_* E_T(r)
  = \pi_{T_0}^* (W \otimes \mathrm{Sym}^r V)_T\)
  is the same. The diagram \textbf{(**)} of
  \cref{lem:quot_boundedness} gives a right-exact sequence
  \[
    (\pi_T)^* (\pi_T)_* \mathcal{G}_i(r) \xrightarrow{h}
    E_T(r) \xrightarrow{q_i(r)} \mathcal{F}_i(r) \to 0,
  \]
  so \(q_i(r)\) is recovered as the cokernel of \(h\); twisting by \(-r\)
  recovers \(q_i\), giving \(q_1 = q_2\) and hence
  \(\langle \mathcal{F}_1, q_1\rangle = \langle \mathcal{F}_2, q_2\rangle\).
\end{proof}

\begin{lemma}[Valuative criterion of properness for \(\Quot\)]
  \label{lem:quot_valuative_criterion}
  \lean{AlgebraicGeometry.TODO.quotValuativeCriterion}
  % SOURCE: [Nitsure], \S 5, sub-section "Valuative Criterion for Properness"
  % (read from references/nitsure-hilbert-quot-src/nitsure-hilbert-quot.tex,
  % L2510-L2542). Original reference: EGA IV(2) 2.8.1, cited there.
  % SOURCE QUOTE:
  % "The functor \(\Quot_{E/X/S}^{\Phi,L}\) satisfies the following valuative
  %  criterion for properness over \(S\): given any discrete valuation ring
  %  \(R\) over \(S\) with quotient field \(K\), the restriction map
  %  \(\)\Quot_{E/X/S}^{\Phi,L}(\spec R) \to \Quot_{E/X/S}^{\Phi,L}(\spec K)\(\)
  %  is bijective. This can be seen as follows.
  %  Given any coherent quotient
  %  \(q: E_K \to \F\) on \(X_R\) which defines an element \(\langle \F, q\rangle\)
  %  of \(\Quot_{E/X/S}^{\Phi,L}(\spec K)\).
  %  Let \(\ov{\F}\) be the image of the composite homomorphism
  %  \(E_R \to j_*(E_K) \to j_*\F\) where \(j : X_K \hra X_R\) is the open
  %  inclusion. Let \(\ov{q} : E_R\to \ov{\F}\) be the induced surjection.
  %  Then we leave it to the reader to verify that
  %  \(\langle \F, q\rangle\) is an element of \(\Quot_{E/X/S}^{\Phi,L}(\spec R)\)
  %  which maps to \(\langle \F, q\rangle\), and is the unique such element.
  %  (Use the basic fact that being flat over a dvr
  %  is the same as being torsion-free.)"
  \textit{Source: [Nitsure], \S 5, "Valuative Criterion for Properness"
  (FGA Explained Ch.~5); EGA IV(2) 2.8.1.}
  The functor \(\Quot^{\Phi,L}_{E/X/S}\) satisfies the valuative criterion
  for properness over \(S\): for every discrete valuation ring \(R\) with
  fraction field \(K\) and every morphism \(\Spec R \to S\), the restriction
  map
  \[
    \Quot^{\Phi,L}_{E/X/S}(\Spec R)
    \longrightarrow
    \Quot^{\Phi,L}_{E/X/S}(\Spec K)
  \]
  is bijective.
\end{lemma}

\begin{proof}
  
  Given a \(\Spec K\)-point \(\langle \mathcal{F}_K, q_K\rangle\) on \(X_K\),
  let \(j : X_K \hookrightarrow X_R\) be the open inclusion. Define
  \(\overline{\mathcal{F}}\) as the image of the composite
  \(E_R \to j_* E_K \to j_* \mathcal{F}_K\), with induced surjection
  \(\overline{q} : E_R \to \overline{\mathcal{F}}\). The sheaf
  \(\overline{\mathcal{F}}\) is torsion-free over \(R\) (being a subsheaf of
  \(j_* \mathcal{F}_K\)), hence \(R\)-flat (flatness over a DVR
  \(=\) torsion-freeness), and restricts on the generic fiber to
  \(\mathcal{F}_K\). The pair \((\overline{\mathcal{F}}, \overline{q})\) is
  the unique extension of \(\langle \mathcal{F}_K, q_K\rangle\) to a
  \(\Spec R\)-point (uniqueness: any other extension agrees on the dense
  open \(X_K \subset X_R\) and is torsion-free, hence equals the image
  construction). The Hilbert polynomial is constant on connected
  \(\Spec R\) by flatness.
\end{proof}

\section{Cohomology and base change}
\label{sec:cohomology_base_change}

The Quot construction uses cohomology-and-base-change in two places: the
boundedness step (\cref{lem:quot_boundedness}) and the Grassmannian
embedding (\cref{lem:quot_alpha_injective}). We record the relevant
statement as a named theorem so the Lean encoding can cite it directly.

\begin{theorem}\leanok
[Flat base change of cohomology]
  \label{thm:flat_base_change_cohomology}
  \lean{AlgebraicGeometry.flatBaseChangeCohomology}
  \uses{thm:flat_base_change_higher, lem:cech_flat_base_change,
    lem:higher_direct_image_quasi_coherent}
  % SOURCE: [Stacks Project], tag 02KH (lemma-flat-base-change-cohomology)
  % (read from references/stacks-coherent.tex, L947-L970).
  % SOURCE QUOTE:
  % "Consider a cartesian diagram of schemes
  %  \(\)
  %  \xymatrix{
  %  X' \ar[d]_{f'} \ar[r]_{g'} & X \ar[d]^f \\
  %  S' \ar[r]^g & S
  %  }
  %  \(\)
  %  Let \(\mathcal{F}\) be a quasi-coherent \(\mathcal{O}_X\)-module
  %  with pullback \(\mathcal{F}' = (g')^*\mathcal{F}\).
  %  Assume that \(g\) is flat and that \(f\) is quasi-compact and quasi-separated.
  %  For any \(i \geq 0\)
  %  \begin{enumerate}
  %  \item the base change map of
  %  Cohomology, Lemma REF
  %  is an isomorphism
  %  \(\)
  %  g^*R^if_*\mathcal{F} \longrightarrow R^if'_*\mathcal{F}',
  %  \(\)
  %  \item if \(S = \Spec(A)\) and \(S' = \Spec(B)\), then
  %  \(H^i(X, \mathcal{F}) \otimes_A B = H^i(X', \mathcal{F}')\).
  %  \end{enumerate}"
  \textit{Source: [Stacks Project], tag 02KH (lemma-flat-base-change-cohomology).}
  Let
  \[
    \begin{array}{ccc}
      X' & \xrightarrow{g'} & X \\
      \downarrow f' &  & \downarrow f \\
      S' & \xrightarrow{g} & S
    \end{array}
  \]
  be a cartesian diagram of schemes with \(g\) flat and \(f\) quasi-compact and
  quasi-separated. Let \(\mathcal{F}\) be a quasi-coherent
  \(\mathcal{O}_X\)-module with pull-back
  \(\mathcal{F}' = (g')^* \mathcal{F}\). For every \(i \ge 0\):
  \begin{enumerate}
    \item the base-change map
      \(g^* R^i f_* \mathcal{F} \to R^i f'_* \mathcal{F}'\) is an isomorphism;
    \item if \(S = \Spec A\) and \(S' = \Spec B\), then \(H^i(X, \mathcal{F})
      \otimes_A B = H^i(X', \mathcal{F}')\).
  \end{enumerate}
\end{theorem}

\subsection{Project-side base-change substrate}
\label{subsec:quot_project_basechange}

The Lean encoding of \cref{thm:flat_base_change_cohomology} (and the
Quot construction's iso between pushforward and tensor-pullback for the
universal section pattern) relies on four project-side declarations
introduced in \texttt{AlgebraicJacobian/Picard/QuotScheme.lean}. These
are typed-sorry today: each carries a substantive Lean type that
encodes the intended mathematical statement, with the body deferred to
later iters (see \cref{sec:quot_lean_encoding} for the
phase-by-phase plan).

\begin{definition}\leanok
[Canonical base-change map on sheaves of modules]
  \label{def:quot_canonical_basechange_map}
  \lean{AlgebraicGeometry.canonicalBaseChangeMap}
  \uses{thm:flat_base_change_cohomology}
  Given a cartesian square of schemes
  \(\square = (g' : X' \to X,\ f' : X' \to S',\ f : X \to S,\ g : S' \to S)\)
  with \(g \circ f' = f \circ g'\) and \(g\) flat, the
  \emph{canonical base-change map} on sheaves of modules is the natural
  transformation \(g'^* f^* \mathcal{F} \longrightarrow f'^* g^* \mathcal{F}\)
  arising from the universal property of pullback. Encoded as a morphism
  of \(\mathrm{ModuleCat}\)-valued sheaf functors on the pulled-back open
  cover; this is the morphism whose iso-status discharges
  \cref{thm:flat_base_change_cohomology}(1) in the present
  finite-locally-free / quasi-coherent-of-finite-type setting.
\end{definition}

\begin{theorem}\leanok
[Stalkwise base-change isomorphism on affines]
  \label{thm:quot_canonical_basechange_app_app_isIso}
  \lean{AlgebraicGeometry.canonicalBaseChangeMap_app_app_isIso}
  \uses{def:quot_canonical_basechange_map, thm:flat_base_change_cohomology, thm:quot_canonical_basechange_app_app_isIso_of_affineCover, thm:quot_canonical_basechange_app_app_isIso_of_isAffineOpen}
  For every affine open \(V \subseteq S'\) and every affine open
  \(U \subseteq X\) lying over \(V\), the local component
  \((\text{canonicalBaseChangeMap}\ \square).\text{app}\ \mathcal{F}.\text{app}\ U\)
  of \cref{def:quot_canonical_basechange_map} is an isomorphism. Proof
  reduces (Stacks 02KH(2)) to the affine-base case: \(H^0\) commutes with
  flat extension of scalars on a quasi-compact quasi-separated affine
  morphism.
\end{theorem}

\begin{theorem}

\leanok
[Affine-base case of the stalkwise base-change isomorphism]
  \label{thm:quot_canonical_basechange_app_app_isIso_of_isAffineBase}
  \lean{AlgebraicGeometry.canonicalBaseChangeMap_app_app_isIso_of_isAffineOpen_of_isAffineBase}
  \uses{def:quot_pullback_app_isoTensor, lem:pushforward_isQuasicoherent, lem:pullback_tildeIso}
  Assume in addition that the base \(S\) is affine, that \(f\) is
  quasi-compact and quasi-separated, that \(g\) is flat, and that
  \(\mathcal{F}\) is quasi-coherent. Then for every affine open
  \(U \subseteq S'\) the section of the canonical base-change map
  (\cref{def:quot_canonical_basechange_map}) over \(U\) is an isomorphism.
  Taking \(V = S\) (affine) as the compatible base open, the statement
  reduces to the algebraic flat base-change identity: the global sections
  \(\Gamma(S', U)\) are flat over \(\Gamma(S, S)\), and the
  tensor-presentation of the pulled-back sections
  (\cref{def:quot_pullback_app_isoTensor}, applied to the pushforward
  \((\pi_f)_* \mathcal{F}\), which stays quasi-coherent by
  \cref{lem:pushforward_isQuasicoherent}) matches the base-change map under
  the Stacks~01HQ identification of \cref{lem:pullback_tildeIso}.
\end{theorem}
\begin{proof} Proved directly in Lean. \end{proof}

\begin{theorem}

\leanok
[Affine-open case of the stalkwise base-change isomorphism]
  \label{thm:quot_canonical_basechange_app_app_isIso_of_isAffineOpen}
  \lean{AlgebraicGeometry.canonicalBaseChangeMap_app_app_isIso_of_isAffineOpen}
  \uses{thm:quot_canonical_basechange_app_app_isIso_of_isAffineBase, lem:quot_pushforward_pullback_section_eq}
  With \(f\) quasi-compact and quasi-separated, \(g\) flat, and
  \(\mathcal{F}\) quasi-coherent (but \(S\) now arbitrary), the section of
  the canonical base-change map (\cref{def:quot_canonical_basechange_map})
  over an affine open \(U \subseteq S'\) is an isomorphism. The substantive
  algebraic content is the affine-base specialization
  (\cref{thm:quot_canonical_basechange_app_app_isIso_of_isAffineBase}); the
  passage from a general base \(S\) to the affine case is a base-side
  Mayer-Vietoris descent along a finite affine cover of \(S\), using the
  definitional section identity
  \cref{lem:quot_pushforward_pullback_section_eq} on each piece.
\end{theorem}
\begin{proof} Proved directly in Lean. \end{proof}

\begin{theorem}

\leanok
[Open-cover descent for the stalkwise base-change isomorphism]
  \label{thm:quot_canonical_basechange_app_app_isIso_of_affineCover}
  \lean{AlgebraicGeometry.canonicalBaseChangeMap_app_app_isIso_of_affineCover}
  \uses{def:quot_canonical_basechange_map}
  Suppose the section of the canonical base-change map
  (\cref{def:quot_canonical_basechange_map}) is an isomorphism over
  \emph{every} affine open of \(S'\). Then it is an isomorphism over every
  open \(U \subseteq S'\). This is the standard Mayer-Vietoris descent for a
  morphism of quasi-coherent sheaves on the base: cover \(U\) by affine
  opens, on each of which the section is an isomorphism by hypothesis, and
  glue along the intersections (quasi-compact because \(f\) is
  quasi-separated). It is the descent half of the stalkwise statement
  \cref{thm:quot_canonical_basechange_app_app_isIso}, whose Lean proof feeds
  the affine-open case
  \cref{thm:quot_canonical_basechange_app_app_isIso_of_isAffineOpen} as the
  affine-open hypothesis.
\end{theorem}
\begin{proof} Proved directly in Lean. \end{proof}

\begin{lemma}

\leanok
[Definitional section identity for pushforward of a pullback]
  \label{lem:quot_pushforward_pullback_section_eq}
  \lean{AlgebraicGeometry.pushforward_pullback_section_eq_pullback_section}
  \uses{def:quot_canonical_basechange_map}
  Let \(f' : X' \to S'\) and \(g' : X' \to X\) be morphisms of schemes,
  \(\mathcal{F}\) a sheaf of \(\mathcal{O}_X\)-modules, and \(U \subseteq S'\)
  an open. The sections over \(U\) of the pushforward along \(f'\) of the
  pullback along \(g'\) of \(\mathcal{F}\) coincide \emph{definitionally}
  with the sections of the pullback \(g'^* \mathcal{F}\) over the preimage
  \(f'^{-1}(U)\):
  \[
    \Gamma\bigl(U,\ f'_* (g'^* \mathcal{F})\bigr)
    \;=\;
    \Gamma\bigl(f'^{-1}(U),\ g'^* \mathcal{F}\bigr).
  \]
  This records step~(3) of the affine-open construction
  (\cref{thm:quot_canonical_basechange_app_app_isIso_of_isAffineOpen}): the
  target of the canonical base-change map
  (\cref{def:quot_canonical_basechange_map}) is read through this identity.
\end{lemma}
\begin{proof} Proved directly in Lean. \end{proof}

\begin{theorem}\leanok
[Canonical base-change map is a sheaf isomorphism]
  \label{thm:quot_canonical_basechange_isIso}
  \lean{AlgebraicGeometry.canonicalBaseChangeMap_isIso}
  \uses{thm:quot_canonical_basechange_app_app_isIso}
  Globalising the stalk-wise statement
  \cref{thm:quot_canonical_basechange_app_app_isIso}: the natural
  transformation \(g'^* f^* \mathcal{F} \to f'^* g^* \mathcal{F}\) of
  \cref{def:quot_canonical_basechange_map} is an isomorphism of
  quasi-coherent sheaves of modules. This is the sheaf-level form of
  \cref{thm:flat_base_change_cohomology}(1).
\end{theorem}

\begin{definition}\leanok
[Tensor-presentation of the pullback-of-sections functor]
  \label{def:quot_pullback_app_isoTensor}
  \lean{AlgebraicGeometry.Scheme.Modules.pullback_app_isoTensor}
  \uses{def:quot_canonical_basechange_map, thm:quot_pullback_app_isoTensor_isBaseChange}
  Given a morphism \(g : X \to Y\) of schemes, affine opens
  \(U \subseteq Y\), \(V \subseteq X\) with \(g(V) \subseteq U\), and a
  quasi-coherent sheaf of modules \(\mathcal{N}\) on \(Y\), the
  \emph{tensor-presentation iso} packages the canonical iso
  \(\Gamma(V,\, g^* \mathcal{N}) \;\cong\;
   \Gamma(V,\, \mathcal{O}_X) \otimes_{\Gamma(U,\, \mathcal{O}_Y)}
     \Gamma(U,\, \mathcal{N})\)
  as a `LinearEquiv`. This is the project-side rfl-bridge used by the
  consumer \(\text{Quot}\)-functor iso assembly: it discharges the
  pullback-app step of the Tilde-isoTop route (Stacks tag 01HQ / 01I8)
  to a base-change-of-scalars identity. The body factors as
  (Step 1, axiom-clean): the underlying linear map
  \texttt{pullback\_app\_isoTensor\_unitAtV};
  (Step 2, project-side typed sorry): the
  \texttt{IsBaseChange}-packaging step
  \texttt{pullback\_app\_isoTensor\_isBaseChange} reducing to
  Stacks 02KE bijectivity of the canonical map.
\end{definition}

\begin{definition}

\leanok
[Adjunction-unit base linear map at a section]
  \label{def:quot_pullback_app_isoTensor_unitAtV}
  \lean{AlgebraicGeometry.pullback_app_isoTensor_unitAtV}
  Given a morphism \(g : Y \to X\) of schemes, a sheaf of
  \(\mathcal{O}_X\)-modules \(\mathcal{N}\), and an open \(V \subseteq X\),
  the unit of the pullback--pushforward adjunction yields a morphism of
  \(\mathcal{O}_X\)-modules \(\mathcal{N} \to g_*(g^* \mathcal{N})\).
  Evaluating its \(V\)-component gives the \(\Gamma(X, V)\)-linear
  \emph{unit-at-section} map
  \[
    \Gamma(\mathcal{N}, V)
    \;\longrightarrow\;
    \Gamma\bigl(g_*(g^* \mathcal{N}),\ V\bigr).
  \]
  This is the axiom-clean building block (Step~1 of the Tilde-isoTop route)
  underlying the tensor-presentation iso of
  \cref{def:quot_pullback_app_isoTensor}.
\end{definition}
\begin{proof} Proved directly in Lean. \end{proof}

\begin{definition}

\leanok
[Section-level base map for the pullback]
  \label{def:quot_pullback_app_isoTensor_baseMap}
  \lean{AlgebraicGeometry.pullback_app_isoTensor_baseMap}
  \uses{def:quot_pullback_app_isoTensor_unitAtV}
  For a morphism \(g : Y \to X\), a sheaf \(\mathcal{N}\) on \(X\), and
  compatible affine opens \(U \subseteq Y\), \(V \subseteq X\) with
  \(g(U) \subseteq V\), composing the unit-at-section map
  (\cref{def:quot_pullback_app_isoTensor_unitAtV}) with the presheaf
  restriction from \(g^{-1}(V)\) down to \(U\) produces the
  \(\Gamma(X, V)\)-linear \emph{base map}
  \[
    \Gamma(\mathcal{N}, V)
    \;\longrightarrow\;
    \Gamma\bigl(g^* \mathcal{N},\ U\bigr),
  \]
  where the \(\Gamma(X, V)\)-action on the target is restriction of scalars
  along the ring map \(\Gamma(X, V) \to \Gamma(Y, U)\) induced by \(g\). This
  is the candidate linear map whose base-change property
  (\cref{thm:quot_pullback_app_isoTensor_baseMap_isBaseChange}) yields the
  tensor presentation.
\end{definition}
\begin{proof} Proved directly in Lean. \end{proof}

\begin{theorem}

\leanok
[The base map is a base change]
  \label{thm:quot_pullback_app_isoTensor_baseMap_isBaseChange}
  \lean{AlgebraicGeometry.pullback_app_isoTensor_baseMap_isBaseChange}
  \uses{def:quot_pullback_app_isoTensor_baseMap, def:pullback_app_isoTensor_sigma}
  For compatible affine opens \(U \subseteq Y\), \(V \subseteq X\) of a
  morphism \(g : Y \to X\) and a quasi-coherent sheaf \(\mathcal{N}\) on
  \(X\), the base map of \cref{def:quot_pullback_app_isoTensor_baseMap}
  exhibits \(\Gamma(g^* \mathcal{N}, U)\) as the base change of
  \(\Gamma(\mathcal{N}, V)\) along the ring map
  \(\Gamma(X, V) \to \Gamma(Y, U)\); i.e.\ the induced
  \(\Gamma(Y, U) \otimes_{\Gamma(X, V)} \Gamma(\mathcal{N}, V)
   \to \Gamma(g^* \mathcal{N}, U)\) is an isomorphism. The proof feeds the
  packaged section-level linear equivalence and its intertwining property
  (\cref{def:pullback_app_isoTensor_sigma}) into the universal
  characterization of a base change.
\end{theorem}
\begin{proof} Proved directly in Lean. \end{proof}

\begin{theorem}

\leanok
[Existence of the tensor-presentation equivalence]
  \label{thm:quot_pullback_app_isoTensor_isBaseChange}
  \lean{AlgebraicGeometry.pullback_app_isoTensor_isBaseChange}
  \uses{thm:quot_pullback_app_isoTensor_baseMap_isBaseChange}
  In the same affine setup, there exists a \(\Gamma(Y, U)\)-linear
  isomorphism
  \[
    \Gamma\bigl(g^* \mathcal{N},\ U\bigr)
    \;\cong\;
    \Gamma(Y, U) \otimes_{\Gamma(X, V)} \Gamma(\mathcal{N}, V),
  \]
  obtained from the base-change property
  (\cref{thm:quot_pullback_app_isoTensor_baseMap_isBaseChange}) by passing to
  the induced equivalence. This is the existence statement consumed by the
  tensor-presentation iso \cref{def:quot_pullback_app_isoTensor}.
\end{theorem}
\begin{proof} Proved directly in Lean. \end{proof}

\begin{definition}

\leanok
[\(\Sigma\)-pair packaging: section-level LinearEquiv + intertwining]
  \label{def:pullback_app_isoTensor_sigma}
  \lean{AlgebraicGeometry.pullback_app_isoTensor_baseMap_sectionLinearEquiv}
  % NOTE iter-283: replaced the spurious `def:quot_pullback_app_isoTensor` with
  % `def:quot_pullback_app_isoTensor_baseMap` to break a dependency cycle that
  % crashed `leanblueprint web`. In Lean (QuotScheme.lean L867) this Sigma-pair
  % is built from `tildeIso_of_isQuasicoherent_isAffineOpen`/`pullback_tildeIso`
  % and its statement references `pullback_app_isoTensor_baseMap`, NOT the
  % high-level final iso `pullback_app_isoTensor`. The final iso is built FROM
  % this Sigma-pair (def:quot_pullback_app_isoTensor -> thm:..._isBaseChange ->
  % thm:..._baseMap_isBaseChange -> def:pullback_app_isoTensor_sigma), so the
  % former edge was the back-edge closing the cycle. The corrected dependency
  % (on the lower-level baseMap) is both acyclic and accurate.
  \uses{def:quot_pullback_app_isoTensor_baseMap, def:quot_canonical_basechange_map,
        lem:pullback_tildeIso, lem:tildeIso_of_isQuasicoherent_isAffineOpen,
        lem:pullback_of_openImmersion_iso_restrict, lem:pushforward_isQuasicoherent}
  Iter-188 Lane F-introduced \(\Sigma\)-pair packaging declaration that
  bundles BOTH the section-level LinearEquiv
  \(\Gamma(V,(g^*\mathcal N)) \;\simeq_{\Gamma(Y,U)}\; \Gamma(Y,U) \otimes
    \Gamma(N,V)\) AND its intertwining property with the canonical
  base-change map into a single \(\Sigma\)-struct returned via
  \texttt{Nonempty} for \texttt{IsBaseChange.of\_equiv} consumption.
  Concretely:
  \[
    \exists\, (f : \Gamma(V,(\text{pullback}\, g).\mathrm{obj}\,
                     \mathcal N) \;\simeq_{\Gamma(Y,U)}\;
                     \Gamma(X,V) \otimes_{\Gamma(Y,U)} \Gamma(N,V)),\;
      \forall x,\; f(1\otimes x)\;=\;\text{baseMap}\,g\,\mathcal N\,e\,x.
  \]
  This packaging idiom keeps consumer signatures clean and is the project-
  side bridge into Mathlib's \texttt{IsBaseChange.of\_equiv} consumer.
  Body is gated on \cref{lem:pullback_tildeIso} (Stacks 01HQ) plus a
  6-stage Beck-Chevalley intertwining decomposition recorded in
  \cref{sec:quot_iter195_bc_substrate} and mirroring the in-file
  comment block at
  \texttt{AlgebraicJacobian/Picard/QuotScheme.lean:999-1037}:
  \begin{itemize}
    \item \textbf{Stage 1} (closed iter-195 via the \(\Sigma\)-pair
      iso-characterizing identity \texttt{\_step2\_apply} together with
      \texttt{inv\_hom\_id}): rewrites
      \((\text{step2}.\mathrm{inv}.\mathrm{app}\,\top)\bigl(\text{tilde.toOpen}\,TR\,\top\,(1 \otimes x)\bigr)\)
      as
      \(\text{baseMap}\,(\text{Spec.map}\,\varphi)\,(\widetilde{\Gamma(\mathcal N, V)})\,e_1\,(\text{tilde.toOpen}\,\Gamma(\mathcal N, V)\,\top\,x)\).
    \item \textbf{Stage 2} (uses (N1) baseMap naturality of
      \cref{lem:baseMap_pullbackComp_apply} together with the
      \(\Sigma\)-pair \texttt{\_step1\_apply} identity): transports
      \(\text{step1}.\mathrm{inv}\) across
      \(\text{baseMap}\,(\text{Spec.map}\,\varphi)\), producing
      \(\text{baseMap}\,(\text{Spec.map}\,\varphi)\,((\text{pullback}\,\mathtt{\_hV.fromSpec}).\mathrm{obj}\,\mathcal N)\,e_2\,(\text{baseMap}\,\mathtt{\_hV.fromSpec}\,\mathcal N\,e_3\,x)\).
    \item \textbf{Stage 3} (uses (N2) baseMap composition via
      \texttt{pullbackComp} of
      \cref{lem:baseMap_pullback_comp_apply}): collapses the nested
      pair of \texttt{baseMap}'s of Stage 2 into a single
      \(\text{baseMap}\,(\text{Spec.map}\,\varphi \circ \mathtt{\_hV.fromSpec})\,\mathcal N\,e_4\,x\)
      along the composite morphism.
    \item \textbf{Stage 4} (uses (N3) baseMap transport via
      \texttt{pullbackCongr} of
      \cref{lem:baseMap_pullbackCongr_apply}): transports the
      composite of Stage 3 along the propositional equality
      \(\text{Spec.map}\,\varphi \circ \mathtt{\_hV.fromSpec}
      = \mathtt{\_hU.fromSpec} \circ g\) (which holds by definition of
      \texttt{Scheme.Hom.appLE} on \(g\) with \(g(V) \subseteq U\)),
      producing
      \(\text{baseMap}\,(\mathtt{\_hU.fromSpec} \circ g)\,\mathcal N\,e_5\,x\).
    \item \textbf{Stage 5} (uses (N2) again,
      \cref{lem:baseMap_pullback_comp_apply}): re-decomposes the
      composite of Stage 4 into nested \texttt{baseMap}'s along
      \(\mathtt{\_hU.fromSpec}\) and \(g\), exposing
      \(\text{baseMap}\,g\,\mathcal N\,e\,x\) as the inner section.
    \item \textbf{Stage 6} (uses (N4) \texttt{step3} inversion of
      \cref{lem:baseMap_inv_step3_open_immersion}): inverts the outer
      \(\text{baseMap}\,\mathtt{\_hU.fromSpec}\) via \texttt{step3} —
      the transport iso of
      \cref{lem:pullback_of_openImmersion_iso_restrict} for the open
      immersion \(\mathtt{\_hU.fromSpec}\) — leaving
      \(\text{baseMap}\,g\,\mathcal N\,e\,x\), which matches the RHS
      of the Σ-pair intertwining property.
  \end{itemize}
\end{definition}

\subsection{Iter-187 analogist-licensed named substrate gaps}
\label{sec:quot_analogist_gaps}

The iter-187 \texttt{mathlib-analogist quotscheme-isbasechange-tilde}
verdict (see \texttt{analogies/quotscheme-isbasechange-tilde.md})
identified two distinct Mathlib substrate gaps at the pinned commit
\texttt{b80f227}. Both are pinned here as named project-side typed
sorries; closing either is iter-189+ sub-build work.

\begin{lemma}

\leanok
[Pullback of tilde is tilde of base change (Stacks 01HQ)]
  \label{lem:pullback_tildeIso}
  \lean{AlgebraicGeometry.pullback_tildeIso}
  \uses{def:quot_canonical_basechange_map}
  \textit{Source: Stacks~\href{https://stacks.math.columbia.edu/tag/01HQ}{01HQ}
  / \href{https://stacks.math.columbia.edu/tag/0BJ8}{0BJ8} (pullback of
  quasi-coherent modules along a Spec morphism).}
  For a ring homomorphism \(\varphi : A \to B\) and an \(A\)-module
  \(M\), there is a canonical iso between the pullback of the
  tilde-sheaf and the tilde of the base-changed module:
  \[
    (\mathtt{Scheme.Modules.pullback}\, (\mathrm{Spec.map}\, \varphi)).\mathrm{obj}\,
      (\widetilde{M})
    \;\cong\;
    \widetilde{B \otimes_A M}.
  \]
  Mathlib \texttt{b80f227} ships
  \texttt{PullbackFree.lean:122 pullbackObjFreeIso} for \emph{free} sheaves
  only — the general-module form (Stacks 01HQ) is unowned.
\end{lemma}
\begin{proof}
  \uses{def:quot_canonical_basechange_map}
  Substantive body (Stacks 01HQ): construct via the
  \texttt{tilde.adjunction} naturality square; the construction passes
  through Spec-restriction transport.

  % NOTE (iter-188 plan-phase): body ~115-200 LOC, iter-189+ sub-build.
  % The iter-188 prover does NOT target this body; the iter-188 Lane F
  % target is the consumer `pullback_app_isoTensor_baseMap_isBaseChange`
  % which is closable via `pullback_tildeIso` as a hypothesis-providing
  % helper.
\end{proof}

\subsection{qcqs section localization and Stacks 01XJ}
\label{sec:quot_qcqs_section_localization}

The T12 session (2026-07-03) closed the \Cref{lem:pushforward_isQuasicoherent}
pin. The route generalizes the affine section-localization keystone to
quasi-compact quasi-separated opens by a Mayer--Vietoris induction, then runs
the tilde--$\Gamma$ comparison backwards at each affine open of the base.

\begin{lemma}[Affine basic-open section localization for quasi-coherent modules]
  \label{lem:qcoh_section_localization_basicOpen}
  \lean{AlgebraicGeometry.Scheme.Modules.isLocalizedModule_basicOpen}
  \leanok
  \textit{Source: Stacks~\href{https://stacks.math.columbia.edu/tag/01IB}{01IB}
  / Hartshorne II.5.3 (\texttt{lemma-invert-f-sections}, affine case).}
  Let $X$ be a scheme, $\mathcal{M}$ a quasi-coherent
  $\mathcal{O}_X$-module, $U \subseteq X$ an affine open, and
  $f \in \Gamma(X, U)$. Then the section restriction
  $\Gamma(\mathcal{M}, U) \to \Gamma(\mathcal{M}, D(f))$ exhibits the target as
  the localization $\Gamma(\mathcal{M}, U)[1/f]$ over $\Gamma(X, U)$.
\end{lemma}
\begin{proof}
  \leanok
  Pull $\mathcal{M}$ back along the affine immersion
  $\operatorname{Spec} \Gamma(X, U) \to X$; the pullback is quasi-coherent
  (pullback along an open immersion preserves quasi-coherence), so its
  tilde--$\Gamma$ counit is an isomorphism, identifying it with
  $\widetilde{\Gamma(\mathcal{M}, U)}$; the sections of a tilde over a basic
  open are the localization, and the section comparison transports the
  statement back to $X$.
\end{proof}

\begin{lemma}[qcqs section localization, torsion half]
  \label{lem:qcqs_section_localization_torsion}
  \lean{AlgebraicGeometry.Scheme.Modules.exists_pow_smul_res_eq_zero_of_isCompact}
  \leanok
  \uses{lem:qcoh_section_localization_basicOpen}
  \textit{Source: Stacks~\href{https://stacks.math.columbia.edu/tag/01P0}{01P0}
  (first half).}
  Let $X$ be a scheme, $\mathcal{M}$ a quasi-coherent $\mathcal{O}_X$-module,
  $W \subseteq X$ open, $g \in \Gamma(X, W)$, and $U \le W$ a quasi-compact
  open. If $x \in \Gamma(\mathcal{M}, U)$ restricts to zero on
  $U \cap D(g)$, then $(g|_U)^n \cdot x = 0$ for some $n \in \mathbb{N}$.
\end{lemma}
\begin{proof}
  \leanok
  \uses{lem:qcoh_section_localization_basicOpen}
  Induction on the quasi-compact open $U$, adding one affine open at a time.
  For $U = \emptyset$ the section vanishes by sheaf separation over the empty
  cover. For $U = S \cup V$ with $V$ affine and the claim known on the
  quasi-compact $S$: the restriction of $x$ to $S$ is killed by a power of
  $g|_S$ by induction, and the restriction to $V$ is killed by a power of
  $g|_V$ by the injectivity half of
  \Cref{lem:qcoh_section_localization_basicOpen} at the affine $V$ (note
  $D(g|_V) = V \cap D(g)$). Taking the maximum $n$ of the two exponents,
  $(g|_U)^n \cdot x$ restricts to zero on both members of the cover
  $\{S, V\}$, hence vanishes by sheaf separation.
\end{proof}

\begin{lemma}[qcqs section localization, surjectivity half]
  \label{lem:qcqs_section_localization_surj}
  \lean{AlgebraicGeometry.Scheme.Modules.exists_res_eq_pow_smul_of_isCompact}
  \leanok
  \uses{lem:qcoh_section_localization_basicOpen, lem:qcqs_section_localization_torsion}
  \textit{Source: Stacks~\href{https://stacks.math.columbia.edu/tag/01P0}{01P0}
  (second half).}
  Let $X$, $\mathcal{M}$, $W$, $g$ be as above, with $W$ quasi-separated, and
  let $U \le W$ be quasi-compact. For every
  $y \in \Gamma(\mathcal{M}, U \cap D(g))$ there are $n \in \mathbb{N}$ and
  $x \in \Gamma(\mathcal{M}, U)$ with
  $x|_{U \cap D(g)} = (g|_{U \cap D(g)})^n \cdot y$.
\end{lemma}
\begin{proof}
  \leanok
  \uses{lem:qcoh_section_localization_basicOpen, lem:qcqs_section_localization_torsion}
  Induction on the quasi-compact $U$ as in
  \Cref{lem:qcqs_section_localization_torsion}. In the step $U = S \cup V$
  ($V$ affine), the induction hypothesis and the affine surjectivity half of
  \Cref{lem:qcoh_section_localization_basicOpen} produce candidate sections
  $x_S$ on $S$ and $x_V$ on $V$ with a common normalization exponent $n$.
  Their difference on the overlap $S \cap V$ --- quasi-compact because $W$ is
  quasi-separated --- restricts to zero on $(S \cap V) \cap D(g)$, so by the
  torsion half a further power $g^m$ kills it; after multiplying both
  candidates by $g^m$ they agree on the overlap and glue to a section $x$
  over $U$ with $x|_{U \cap D(g)} = g^{n+m} \cdot y$, checked by sheaf
  separation over the cover $\{S \cap D(g), D(g|_V)\}$ of $U \cap D(g)$.
\end{proof}

\begin{lemma}[qcqs section localization]
  \label{lem:qcqs_section_localization}
  \lean{AlgebraicGeometry.Scheme.Modules.isLocalizedModule_basicOpen_of_isCompact}
  \leanok
  \uses{lem:qcqs_section_localization_torsion, lem:qcqs_section_localization_surj}
  \textit{Source: Stacks~\href{https://stacks.math.columbia.edu/tag/01P0}{01P0}.}
  Let $X$ be a scheme, $\mathcal{M}$ a quasi-coherent $\mathcal{O}_X$-module,
  and $W \subseteq X$ a quasi-compact, quasi-separated open. For every
  $g \in \Gamma(X, W)$ the section restriction
  $\Gamma(\mathcal{M}, W) \to \Gamma(\mathcal{M}, D(g))$ exhibits the target
  as the localization $\Gamma(\mathcal{M}, W)[1/g]$ over $\Gamma(X, W)$.
\end{lemma}
\begin{proof}
  \leanok
  \uses{lem:qcqs_section_localization_torsion, lem:qcqs_section_localization_surj}
  The element $g$ restricts to a unit of $\Gamma(X, D(g))$, so every power of
  $g$ acts invertibly on $\Gamma(\mathcal{M}, D(g))$; the surjectivity and
  torsion (injectivity) conditions of the localization are
  \Cref{lem:qcqs_section_localization_surj} and
  \Cref{lem:qcqs_section_localization_torsion} at $U := W$.
\end{proof}

\begin{lemma}[Converse tilde--$\Gamma$ transport at an affine open]
  \label{lem:fromTildeGamma_pullback_fromSpec_converse}
  \lean{AlgebraicGeometry.Scheme.Modules.isIso_fromTildeΓ_pullback_fromSpec_of_isLocalizedModule}
  \leanok
  \uses{lem:qcoh_section_localization_basicOpen}
  Let $Y$ be a scheme, $\mathcal{N}$ an $\mathcal{O}_Y$-module, and
  $U \subseteq Y$ an affine open with canonical immersion
  $j : \operatorname{Spec} \Gamma(Y, U) \to Y$. If for every
  $f' \in \Gamma(Y, U)$ the restriction
  $\Gamma(\mathcal{N}, U) \to \Gamma(\mathcal{N}, D(f'))$ is a localization at
  the powers of $f'$ over $\Gamma(Y, U)$, then the tilde--$\Gamma$ counit of
  the pullback $j^* \mathcal{N}$ is an isomorphism.
\end{lemma}
\begin{proof}
  \leanok
  The sections of $j^* \mathcal{N}$ over an open of the spectrum are the
  sections of $\mathcal{N}$ over its image ($j$ is an open immersion), with
  $j''(\top) = U$ and $j''(D(f')) = D(f')$; under these identifications --
  which are semilinear over the canonical section-ring comparisons, by the
  naturality of the immersion's application isomorphisms -- the hypothesis
  says exactly that all basic-open section restrictions of $j^* \mathcal{N}$
  on the spectrum are localizations, which characterizes the invertibility of
  the tilde--$\Gamma$ counit.
\end{proof}

\begin{lemma}[Pushforward basic-open section localization]
  \label{lem:pushforward_section_localization}
  \lean{AlgebraicGeometry.Scheme.Modules.isLocalizedModule_basicOpen_pushforward}
  \leanok
  \uses{lem:qcqs_section_localization}
  Let $\pi : X \to S$ be a quasi-compact, quasi-separated morphism of
  schemes, $\mathcal{F}$ a quasi-coherent $\mathcal{O}_X$-module,
  $U \subseteq S$ an affine open, and $f' \in \Gamma(S, U)$. Then the section
  restriction
  $\Gamma(\pi_* \mathcal{F}, U) \to \Gamma(\pi_* \mathcal{F}, D(f'))$ is a
  localization at the powers of $f'$ over $\Gamma(S, U)$.
\end{lemma}
\begin{proof}
  \leanok
  \uses{lem:qcqs_section_localization}
  The preimage $W := \pi^{-1}(U)$ is quasi-compact and quasi-separated
  ($\pi$ is qcqs and $U$ is affine), and
  $\pi^{-1}(D_S(f')) = D_X(g)$ for $g := \pi^\sharp f' \in \Gamma(X, W)$. The
  sections of the pushforward are by definition the sections of
  $\mathcal{F}$ over these preimages, with the $\Gamma(S, U)$-action given
  through $\pi^\sharp$, so the statement is
  \Cref{lem:qcqs_section_localization} for $\mathcal{F}$ at $g$ over $W$,
  read through $\pi^\sharp$ (a localization over the base ring restricts to a
  localization along any ring map matching the multiplicative sets).
\end{proof}

\begin{lemma}

\leanok
[Pushforward of quasi-coherent modules is quasi-coherent (Stacks 01XJ)]
  \label{lem:pushforward_isQuasicoherent}
  \lean{AlgebraicGeometry.pushforward_isQuasicoherent}
  \textit{Source: Stacks~\href{https://stacks.math.columbia.edu/tag/01XJ}{01XJ}
  (preservation of quasi-coherence under pushforward).}
  Let \(f : X \to Y\) be a quasi-compact, quasi-separated morphism of
  schemes, and let \(\mathcal{F}\) be a quasi-coherent
  \(\mathcal{O}_X\)-module. Then \(f_* \mathcal{F}\) is a quasi-coherent
  \(\mathcal{O}_Y\)-module. Mathlib \texttt{b80f227} ships no
  \texttt{IsQuasicoherent.pushforward} instance at the
  \texttt{Scheme.Modules} layer.
\end{lemma}
\begin{proof}
  \leanok
  \uses{lem:pushforward_section_localization, lem:fromTildeGamma_pullback_fromSpec_converse}
  Quasi-coherence is local on $Y$, so it suffices to produce a presentation
  of the slice of $f_* \mathcal{F}$ over each affine open $U \subseteq Y$.
  By \Cref{lem:pushforward_section_localization} all basic-open section
  restrictions of $f_* \mathcal{F}$ at $U$ are localizations, so by
  \Cref{lem:fromTildeGamma_pullback_fromSpec_converse} the pullback of
  $f_* \mathcal{F}$ along $\operatorname{Spec} \Gamma(Y, U) \to Y$ is
  isomorphic to the tilde of its global sections. A tilde admits a global
  presentation; transporting it along
  $U \hookrightarrow Y = (\operatorname{Spec}\Gamma(Y,U) \cong U)
  \hookrightarrow Y$ yields a presentation of the geometric restriction of
  $f_* \mathcal{F}$ to $U$, hence of the slice. (The previously sketched
  ``right adjoints preserve quasi-coherence'' argument is not a proof:
  adjointness yields colimit preservation only.)
\end{proof}

\subsection{Iter-189 analogist-licensed unbundle pins}
\label{sec:quot_iter189_unbundle}

The iter-189 \texttt{mathlib-analogist lane-f-isbasechange} verdict
(see \texttt{analogies/lane-f-isbasechange.md}) returned verdict (A)
\emph{STRUCTURAL OK} on the project's use of
\texttt{IsBaseChange.of\_equiv}. The corrective was to \emph{unbundle}
the iter-188 \(\Sigma\)-pair helper's three internal Mathlib gaps into
separately named typed-sorry pins so that each can be targeted
individually (~30-50 LOC apiece) in subsequent iters. Iter-189 Lane F
prover landed the unbundle. The two new pins are recorded below
(parallel to the existing \cref{lem:pullback_tildeIso} pin, which
remains Step 2 of the 3-step chain).

\begin{lemma}

\leanok
[Quasi-coherent module on an affine open is \texttt{tilde} of its sections (Stacks 01I8)]
  \label{lem:tildeIso_of_isQuasicoherent_isAffineOpen}
  \lean{AlgebraicGeometry.tildeIso_of_isQuasicoherent_isAffineOpen}
  \uses{def:quot_canonical_basechange_map}
  \textit{Source: Stacks~\href{https://stacks.math.columbia.edu/tag/01I8}{01I8}
  (a quasi-coherent module on an affine scheme is \(\widetilde M\) for
  some module).}
  Let \(X\) be a scheme, \(\mathcal N\) a quasi-coherent
  \(\mathcal O_X\)-module, and \(V \subset X\) an affine open. Then the
  pullback of \(\mathcal N\) along \texttt{hV.fromSpec} is canonically
  isomorphic to \texttt{tilde} of its sections:
  \[
    (\mathtt{Scheme.Modules.pullback}\,\mathtt{hV.fromSpec}).\mathrm{obj}\,
      \mathcal N
    \;\cong\;
    \widetilde{\,\Gamma(\mathcal N, V)\,}
    \quad\text{(on \(\mathrm{Spec}\,\Gamma(X, V)\)).}
  \]
  Mathlib \texttt{b80f227} ships the analog for \emph{free}
  quasi-coherent modules (\texttt{isIso\_fromTildeΓ\_of\_presentation}
  partial), but the general statement
  ``\([N.IsQuasicoherent] \Rightarrow N|_V \cong \widetilde{\Gamma(N,V)}\)
  on an affine open \(V\)'' (Stacks 01I8 \emph{lemma-affine-qc-module-iff})
  is unowned at the \texttt{Scheme.Modules} layer.
\end{lemma}
\begin{proof}
  \leanok
  \uses{def:quot_canonical_basechange_map, lem:modules_pullback_mathlib,
    lem:modules_restrictFunctorIsoPullback_mathlib}
  Write \(j = \mathtt{hV.fromSpec} : \Spec\Gamma(X,V) \to X\) (an open
  immersion with image \(V\)) and \(G = j^{*}\mathcal N\). Pullback along an
  open immersion preserves quasi-coherence
  (Lean \texttt{isQuasicoherent\_pullback\_fromSpec}), so \(G\) is a
  quasi-coherent module on the affine scheme \(\Spec\Gamma(X,V)\); by the
  affine structure theorem (Stacks 01I8, Lean
  \texttt{isIso\_fromTildeΓ\_of\_isQuasicoherent}) the tilde--\(\Gamma\)
  counit \(\widetilde{\Gamma(G,\top)} \to G\) is an isomorphism.

  The canonical base map \(b : \Gamma(\mathcal N, V) \to \Gamma(G, \top)\)
  (the \(V\)-sections of the pullback--pushforward adjunction unit followed by
  restriction along \(\top \le j^{-1}V\)) is bijective. Indeed the two
  adjunctions \(\mathrm{restrict} \dashv j_{*}\) and \(j^{*} \dashv j_{*}\)
  share the right adjoint, so their units are identified through the canonical
  natural isomorphism \(\mathrm{restrict} \cong j^{*}\)
  (\cref{lem:modules_restrictFunctorIsoPullback_mathlib}). Under this
  identification, the \(V\)-component of the unit of the restriction
  adjunction is the presheaf restriction of \(\mathcal N\) along the
  \emph{equality} of opens \(j(j^{-1}V) = V\), hence bijective; the component
  of the natural isomorphism is bijective; and the final restriction is along
  the equality \(\top = j^{-1}V\). Thus \(b\) is a composite of three
  bijections.

  Since \(b\) is a bijective \(\Gamma(X,V)\)-linear map (the
  \(\Gamma(X,V)\)-action on \(\Gamma(G,\top)\) being through the canonical
  ring isomorphism \(\Gamma(X,V) \cong \Gamma(\Spec\Gamma(X,V),\top)\)),
  it induces an isomorphism
  \(\widetilde{b} : \widetilde{\Gamma(\mathcal N,V)} \cong
  \widetilde{\Gamma(G,\top)}\). The required isomorphism is the inverse of
  \(\widetilde{b}\) followed by the counit. The characterizing section
  identity holds because \(\widetilde{(-)}\)-functoriality intertwines the
  canonical maps \(M \to \Gamma(\widetilde M, U)\) (naturality of
  \texttt{toOpen}), while the counit composed with \(\mathtt{toOpen}\) at
  \(\top\) is the identity restriction.
\end{proof}

\begin{lemma}

\leanok
[Section-level transport for pullback along an affine open's
\texttt{fromSpec}]
  \label{lem:pullback_of_openImmersion_iso_restrict}
  \lean{AlgebraicGeometry.pullback_of_openImmersion_iso_restrict}
  \uses{def:quot_canonical_basechange_map}
  \textit{Source: Project-bespoke transport, formed from
  Stacks~\href{https://stacks.math.columbia.edu/tag/01HH}{01HH} +
  \texttt{tilde.isoTop}.}
  Let \(Y\) be a scheme, \(\mathcal N\) an \(\mathcal O_Y\)-module, and
  \(U \subset Y\) an affine open. Equip \(\Gamma((\mathrm{Spec}\,\Gamma(Y, U)),
  \top)\) with the \(\Gamma(Y, U)\)-algebra structure inherited from
  \texttt{Scheme.ΓSpecIso} (the canonical
  \(\Gamma(\mathrm{Spec}\,R, \top) \cong R\) for any commutative ring \(R\)),
  and equip
  \(\Gamma((\text{pullback}\,\mathtt{hU.fromSpec}).\mathrm{obj}\,\mathcal N,
  \top)\) with the \(\Gamma(Y, U)\)-module structure inherited via
  \texttt{Module.compHom}. Then there is a canonical \(\Gamma(Y, U)\)-linear
  isomorphism between this pullback-section module and the original section
  module \(\Gamma(\mathcal N, U)\):
  \[
    \Gamma((\text{pullback}\,\mathtt{hU.fromSpec}).\mathrm{obj}\,\mathcal N,
      \top)
    \;\simeq_{\Gamma(Y, U)}\;
    \Gamma(\mathcal N, U).
  \]
  The map identifies a section of the pulled-back sheaf at \(\top\) of
  \(\mathrm{Spec}\,\Gamma(Y, U)\) with the corresponding section of
  \(\mathcal N\) on the affine open \(U \subset Y\). This pin is the
  iter-190+ prover target as Step 3 of the iter-189 \(\Sigma\)-pair
  unbundle.
\end{lemma}
\begin{proof}
  \uses{def:quot_canonical_basechange_map}
  Substantive body: chain \texttt{hU.toSpecΓ\_fromSpec} (identifies
  \texttt{pullback hU.fromSpec} with \texttt{restrict\_to\_U} up to
  natural iso) with naturality of \texttt{tilde.isoTop} (extracts the
  \(\top\)-section of \texttt{tilde} as the input module), then transport
  the \(\Gamma(Y, U)\)-linear structure through the chain.

  % NOTE (iter-189 plan-phase): body ~30-50 LOC, iter-190+ sub-build.
  % This is Step 3 of the 3-pin unbundle of `_sectionLinearEquiv`, and
  % is the LIGHTEST of the 3 pins per the analogist verdict (the chain
  % `hU.toSpecΓ_fromSpec` + `tilde.isoTop` naturality is a routine but
  % multi-step compositional transport). The iter-190 Lane F prover
  % targets this pin axiom-clean.
\end{proof}

\subsection{Iter-195 Beck-Chevalley 6-stage decomposition + (N1)--(N4) substrate}
\label{sec:quot_iter195_bc_substrate}

The iter-195 plan-phase refactor \texttt{lane-f-step12-sigma-pair}
upgraded the helpers \texttt{step1} and \texttt{step2} of the
Tilde-isoTop route into \(\Sigma\)-pair components carrying
iso-characterizing identities \texttt{\_step1\_apply} and
\texttt{\_step2\_apply}. With these identities in hand, the
Beck-Chevalley intertwining of
\cref{def:pullback_app_isoTensor_sigma} unfolds into the six named
stages spelled out in the prose of that definition; mirror copy at
\texttt{AlgebraicJacobian/Picard/QuotScheme.lean:999-1037}. Stage~1
is closed axiom-clean by the iter-195 Lane F prover; Stages 2--6 are
gated on four substrate naturality lemmas (N1)--(N4) for the
project-side helper \texttt{Scheme.Modules.baseMap}, introduced here
as the iter-196+ prover targets.

The four substrate lemmas all follow a single logical pattern. The
project's \texttt{Scheme.Modules.baseMap} is defined from the unit of
\texttt{pullbackPushforwardAdjunction} (Mathlib
\texttt{Mathlib.AlgebraicGeometry.Sheaves.Modules.Pullback}), so
(N1)--(N3) record naturality squares of that adjunction unit at
triples of morphisms / module homomorphisms / propositional
equalities, and (N4) records the inversion of the \texttt{step3}
component built from \texttt{restrictFunctorIsoPullback} (Mathlib
\texttt{Mathlib.AlgebraicGeometry.OpenImmersion}). Together they
discharge the Beck-Chevalley sequence stage-by-stage. LOC estimate
\(\sim\)20--30 LOC per sub-lemma; total \(\sim\)80--120 LOC of
substrate volume. This is the iter-195 progress-critic-flagged
``volume gap'' relative to the iter-195 plan-phase estimate
(\(\sim\)10--30 LOC for the consumer, which did not budget the
substrate).

\begin{lemma}
[(N1) baseMap naturality in input sheaf]
  \label{lem:baseMap_pullbackComp_apply}
  \lean{AlgebraicGeometry.Scheme.Modules.baseMap_pullbackComp_apply}
  \uses{def:quot_canonical_basechange_map, lem:pullback_tildeIso}
  Let \(g : X \to Y\) be a morphism of schemes, let \(\mathcal N\),
  \(\mathcal N'\) be \(\mathcal O_Y\)-modules, and let
  \(\eta : \mathcal N \to \mathcal N'\) be a morphism of
  \(\mathcal O_Y\)-modules. For affine opens \(U \subseteq Y\) and
  \(V \subseteq X\) with \(g(V) \subseteq U\) and any
  \(x \in \Gamma(\mathcal N, U)\), the section-level component of
  \cref{def:quot_canonical_basechange_map} satisfies
  \[
    \bigl((\text{pullback}\,g).\mathrm{map}\,\eta\bigr).\mathrm{app}\,V\;
    \bigl(\text{baseMap}\,g\,\mathcal N\,e\,x\bigr)
    \;=\;
    \text{baseMap}\,g\,\mathcal N'\,e\,(\eta.\mathrm{app}\,U\,x),
  \]
  i.e.\ \texttt{baseMap} is natural in its module argument. Used in
  Stage~2 of the \cref{def:pullback_app_isoTensor_sigma}
  decomposition to transport \(\text{step1}.\mathrm{inv}\) (a
  morphism of \(\mathcal O_Y\)-modules) across \texttt{baseMap} of
  \(\text{Spec.map}\,\varphi\).
\end{lemma}
\begin{proof}
  \uses{def:quot_canonical_basechange_map}
  Direct from the naturality square of
  \texttt{pullbackPushforwardAdjunction.unit} at the morphism
  \(\eta : \mathcal N \to \mathcal N'\) in the input category, then
  passed to the section level via the definitional unfolding of
  \texttt{Scheme.Modules.baseMap} (which is a section-level reading
  of the unit composed with the canonical pushforward-section map).

  % NOTE (iter-196 plan-phase): body ~20-30 LOC. (N1) of the iter-195
  % 6-stage Beck-Chevalley decomposition; iter-196+ Lane F prover target.
\end{proof}

\begin{lemma}
[(N2) baseMap composition via pullbackComp]
  \label{lem:baseMap_pullback_comp_apply}
  \lean{AlgebraicGeometry.Scheme.Modules.baseMap_pullback_comp_apply}
  \uses{def:quot_canonical_basechange_map, lem:pullback_tildeIso}
  Let \(f : X \to Y\) and \(g : Y \to Z\) be morphisms of schemes and
  let \(\mathcal N\) be an \(\mathcal O_Z\)-module. The
  Mathlib-canonical iso
  \(\text{pullbackComp}\,f\,g :
    (\text{pullback}\,f) \circ (\text{pullback}\,g)
    \cong \text{pullback}\,(f \circ g)\)
  intertwines \texttt{baseMap}: for affine opens \(W \subseteq Z\),
  \(V \subseteq X\) with \(f(g(V)) \subseteq W\) and any
  \(x \in \Gamma(\mathcal N, W)\),
  \[
    ((\text{pullbackComp}\,f\,g)\,\mathcal N).\mathrm{hom}.\mathrm{app}\,V\;
    \bigl(\text{baseMap}\,f\,((\text{pullback}\,g).\mathrm{obj}\,\mathcal N)\,
       e_f\,(\text{baseMap}\,g\,\mathcal N\,e_g\,x)\bigr)
    \;=\;
    \text{baseMap}\,(f \circ g)\,\mathcal N\,e\,x.
  \]
  In words: applying \texttt{baseMap} along \(g\) then along \(f\),
  routed through the canonical composition iso, yields the
  \texttt{baseMap} along the composite \(f \circ g\). Used in Stages~3
  and~5 of the \cref{def:pullback_app_isoTensor_sigma}
  decomposition.
\end{lemma}
\begin{proof}
  \uses{def:quot_canonical_basechange_map}
  The unit of a composition of left adjoints
  \(F \circ G \dashv R \circ R'\) satisfies the canonical composition
  identity \(\eta_{FG} = R'(\eta_F) \circ \eta_G\) (cf.\ Mathlib
  \texttt{Adjunction.comp} unit decomposition). Specialised to the
  pullback-pushforward adjunction at the triple of schemes
  \(X \xrightarrow{f} Y \xrightarrow{g} Z\) and transported to the
  section level via \texttt{Scheme.Modules.baseMap}'s definition,
  this yields the stated formula.

  % NOTE (iter-196 plan-phase): body ~25-35 LOC. (N2) of the iter-195
  % 6-stage Beck-Chevalley decomposition; iter-196+ Lane F prover target.
  % Used twice in the decomposition (Stages 3 and 5).
\end{proof}

\begin{lemma}
[(N3) baseMap transport via pullbackCongr]
  \label{lem:baseMap_pullbackCongr_apply}
  \lean{AlgebraicGeometry.Scheme.Modules.baseMap_pullbackCongr_apply}
  \uses{def:quot_canonical_basechange_map, lem:pullback_tildeIso}
  Let \(f_1, f_2 : X \to Y\) be morphisms of schemes with
  \(f_1 = f_2\) as a propositional equality \(h_{eq}\), and let
  \(\mathcal N\) be an \(\mathcal O_Y\)-module. The Mathlib-canonical
  equality-transport iso
  \(\text{pullbackCongr}\,h_{eq} : \text{pullback}\,f_1
    \cong \text{pullback}\,f_2\)
  intertwines \texttt{baseMap}: for affine opens \(U \subseteq Y\),
  \(V \subseteq X\) with \(f_1(V) \subseteq U\) (equivalently
  \(f_2(V) \subseteq U\) via \(h_{eq}\)) and any
  \(x \in \Gamma(\mathcal N, U)\),
  \[
    ((\text{pullbackCongr}\,h_{eq})\,\mathcal N).\mathrm{inv}.\mathrm{app}\,V\;
    \bigl(\text{baseMap}\,f_2\,\mathcal N\,e_2\,x\bigr)
    \;=\;
    \text{baseMap}\,f_1\,\mathcal N\,e_1\,x.
  \]
  I.e.\ \texttt{baseMap} is invariant under the equality-transport iso
  of the morphism argument. Used in Stage~4 of the
  \cref{def:pullback_app_isoTensor_sigma} decomposition to re-associate
  the composite
  \(\text{Spec.map}\,\varphi \circ \mathtt{\_hV.fromSpec}
    = \mathtt{\_hU.fromSpec} \circ g\) under the Σ-pair packaging.
\end{lemma}
\begin{proof}
  \uses{def:quot_canonical_basechange_map}
  By propositional-equality elimination: write
  \(\text{pullbackCongr}\,h_{eq}\) as
  \(\text{eqToIso}\,(\text{congrArg}\,\text{pullback}\,h_{eq})\) and
  unfold \texttt{eqToIso\_app} to reduce to \texttt{Eq.mpr}; the
  definition of \texttt{baseMap} depends on the morphism argument only
  through \texttt{pullbackPushforwardAdjunction.unit}, which is itself
  definitionally invariant under \texttt{Eq.mpr} along
  \(h_{eq}\).

  % NOTE (iter-196 plan-phase): body ~15-25 LOC; the easiest of the
  % four — almost pure equality transport. (N3) of the iter-195 6-stage
  % Beck-Chevalley decomposition; iter-196+ Lane F prover target.
\end{proof}

\begin{lemma}
[(N4) step3 inversion identity for an open immersion]
  \label{lem:baseMap_inv_step3_open_immersion}
  \lean{AlgebraicGeometry.Scheme.Modules.baseMap_inv_step3_open_immersion}
  \uses{def:quot_canonical_basechange_map,
        lem:pullback_of_openImmersion_iso_restrict}
  Let \(Y\) be a scheme, \(U \subseteq Y\) an affine open with
  canonical open immersion
  \(\mathtt{\_hU.fromSpec} : \text{Spec}\,\Gamma(Y, U) \to Y\), and
  \(\mathcal N\) an \(\mathcal O_Y\)-module. Then the iter-189
  Spec-restriction transport iso \texttt{step3} (constructed in
  \cref{lem:pullback_of_openImmersion_iso_restrict}) is the inverse
  of \texttt{baseMap} along \(\mathtt{\_hU.fromSpec}\) at the
  \(\top\)-section:
  \[
    \text{step3}\bigl(
      \text{baseMap}\,\mathtt{\_hU.fromSpec}\,
        ((\text{pullback}\,g).\mathrm{obj}\,\mathcal N)\,
        \text{le\_top}'\,y
    \bigr) \;=\; y
    \qquad\text{for every } y \in \Gamma(\mathcal N, U).
  \]
  This identifies the inverse of \texttt{step3} with the natural
  \texttt{baseMap} action of the open immersion
  \(\mathtt{\_hU.fromSpec}\). Used in Stage~6 of the
  \cref{def:pullback_app_isoTensor_sigma} decomposition, closing the
  intertwining identity.
\end{lemma}
\begin{proof}
  \uses{def:quot_canonical_basechange_map,
        lem:pullback_of_openImmersion_iso_restrict}
  The transport iso \texttt{step3} is built (in
  \cref{lem:pullback_of_openImmersion_iso_restrict}) from
  \texttt{restrictFunctorIsoPullback} (Mathlib
  \texttt{AlgebraicGeometry.OpenImmersion}), which exhibits, for an
  open immersion \(j\), a natural iso between \(j^*\) and the
  restriction functor \(\text{restrict}\,j\). At the \(\top\)-section
  of \(\text{Spec}\,\Gamma(Y, U)\), the inverse of this natural iso is
  exactly \texttt{baseMap}\,\(\mathtt{\_hU.fromSpec}\), because the
  adjunction unit for an open-immersion pullback acts at \(\top\) as
  the canonical restriction-to-section map.

  % NOTE (iter-196 plan-phase): body ~20-30 LOC. (N4) of the iter-195
  % 6-stage Beck-Chevalley decomposition; iter-196+ Lane F prover target.
\end{proof}

\section{Lean encoding}
\label{sec:quot_lean_encoding}

The Lean target file is the new
\texttt{AlgebraicJacobian/Picard/QuotScheme.lean}. The construction
proceeds in three phases.

\paragraph{Phase 1: Hilbert polynomial (\cref{def:hilbert_polynomial}).}
On a finite-type morphism \(\pi : X \to S\) with \(S\) noetherian and a
coherent \(\mathcal{F}\) on \(X\) with proper support over \(S\), the per-fiber
Hilbert polynomial is encoded as a function \(s \mapsto
\Phi_{\mathcal{F},s} \in \mathbb{Q}[\lambda]\). Mathlib provides
\texttt{CategoryTheory.PolynomialModule} and Euler-characteristic
infrastructure on coherent sheaves on noetherian schemes; the polynomial
property of \(m \mapsto \chi(\mathcal{F} \otimes L^{\otimes m})\) requires
Snapper's Lemma, which is a project-side sub-build on top of the
graded-Euler-characteristic infrastructure.

\paragraph{Phase 2: Grassmannian scheme (\cref{def:grassmannian_scheme},
\cref{thm:grassmannian_representable}).} The Grassmannian
\(\mathrm{Gr}_S(V, d)\) is encoded as the relative gluing of \(\binom{r}{d}\)
affine charts \(U^I\) (each isomorphic to a relative affine space
\(\mathbb{A}^{d(r-d)}_S\)) via the Pl\"ucker cocycle \(\theta_{I,J}\). The
gluing backbone is \texttt{AlgebraicGeometry.Scheme.GlueData}. The
universal quotient sheaf \(\mathcal{U}\) is glued by the transition cocycle
\(g_{I,J} = (X^I_J)^{-1}\), and the determinant line bundle
\(\det \mathcal{U}\) gives the Pl\"ucker closed embedding into
\(\mathbb{P}_S(\bigwedge^d V)\). The representability statement uses
\texttt{CategoryTheory.Functor.RepresentableBy} (as in
\cref{chap:Picard_RelativeSpec}) on the universal pair
\((\mathrm{Gr}_S(V, d), \mathcal{U})\).

\paragraph{Phase 3: Quot scheme (\cref{thm:quot_representable}).}
The Quot functor is constructed as a locally closed subfunctor of the
Grassmannian
\(\mathrm{Grass}(W \otimes_{\mathcal{O}_S} \mathrm{Sym}^r V, \Phi(r))\), via:
\begin{enumerate}
  \item the uniform \(m\)-regularity bound from
    \cref{lem:quot_boundedness} (input: Castelnuovo--Mumford theorem,
    \texttt{AlgebraicGeometry.CastelnuovoMumford} -- project-side sub-build);
  \item the Grassmannian embedding \(\alpha\) from
    \cref{lem:quot_alpha_injective};
  \item the flattening stratification (\cref{thm:flattening_stratification_exists}
    of \cref{chap:Picard_FlatteningStratification}, A.2.a) applied to the
    universal quotient on the Grassmannian, producing the locally closed
    stratum \(T_0^\Phi\);
  \item the valuative-criterion upgrade
    (\cref{lem:quot_valuative_criterion}) to closed embedding, hence
    projective \(S\)-scheme.
\end{enumerate}
The reduction \cref{lem:quot_reduction_to_pi_star_W} from the universal
case to the general case is encoded as a closed embedding of representable
subfunctors via the kernel-vanishing scheme.

\paragraph{Cohomology and base change.}
\cref{thm:flat_base_change_cohomology} is the Mathlib-equivalent of
Stacks tag 02KH. Mathlib carries a partial form
(\texttt{AlgebraicGeometry.flat\_base\_change}) for affine bases; the
full \(R^i f_*\) statement is a project-side bridge if not present at the
pinned commit.

\paragraph{Dependencies summary.}
This file directly consumes
\cref{chap:Picard_RelativeSpec} (A.1.a, on disk) for relative-spectrum
gluing infrastructure and
\cref{thm:flattening_stratification_exists} from
\cref{chap:Picard_FlatteningStratification} (A.2.a, written this iter)
for the flattening stratum. Downstream, \cref{chap:Picard_FGAPicRepresentability}
(A.2.c) consumes \cref{thm:quot_representable} via the special case
\(E = \mathcal{O}_C\) on a smooth proper curve \(C/k\), giving the Hilbert
scheme \(\Hilb_{C/k} = \Quot^{\Phi,\mathcal{O}_C(1)}_{\mathcal{O}_C/C/k}\)
and its open subscheme \(\mathrm{Div}_{C/k}\) of relative effective divisors.

\section{Out of scope}
\label{sec:quot_out_of_scope}

This chapter packages \emph{only} the Quot-foundations layer: the Hilbert
polynomial, the Quot functor, and the Grassmannian scheme with its
representability. The following are explicitly out of scope here:
\begin{itemize}
  \item \textbf{Representability of the Quot scheme} -- the
    Grothendieck--Altman--Kleiman existence theorem and its proof skeleton
    (boundedness, Grassmannian embedding, flattening stratification,
    valuative criterion) are not constructed in this chapter.
  \item \textbf{Flat base change of cohomology} (Stacks 02KH) and its
    associated prerequisites -- not constructed in this chapter.
  \item \textbf{Hilbert scheme as a separate chapter}. The Hilbert scheme
    \(\Hilb_{X/S}\) is the special case
    \(\Hilb_{X/S} = \Quot_{\mathcal{O}_X/X/S}\) of the Quot functor and is
    introduced inline (\cref{def:quot_functor}); no separate chapter.
  \item \textbf{Snapper's Lemma} -- the polynomial-eventually property of
    Euler characteristics, used in \cref{def:hilbert_polynomial}. Cited
    via Nitsure \S 1 and (textbook reference) Hartshorne III.5.2;
    Lean formalization is a prerequisite sub-task.
  \item \textbf{Identification with \(\mathbb{P}^n\) and relative
    Grassmannian variants} (Nitsure \S 1, Exercises (1)--(4)) --
    optional applications, not consumed by Route A.
  \item \textbf{Group-scheme structure on \(\Pic_{C/k}\)} (A.2.c) -- the
    abelian-group-scheme refinement is downstream. Its blueprint home is the
    FGA Picard-representability assembly (the chapter that bundles the relative
    Picard functor, its representability, and the abelian-group-scheme
    structure on \(\Pic_{C/k}\)); that assembly lives in the parent project
    which consumes this Quot-foundations layer and is not part of the present
    subproject.
  \item \textbf{Identity-component \(\Pic^0\) and abelian-variety
    structure} (A.3) -- a separate Route A sub-row, fed by but not
    constructed in this chapter.
\end{itemize}
